\documentclass[letterpaper,12pt]{article}
% Version 7 (2026-06-15): adds JPE Policy 1.6 data citations -- \nocite of the
% data_* @misc entries in paper/FDAPaperv1.bib (MarketScan, Dimensions, CRSP,
% Compustat, Nexis Uni, FDA MAUDE/510(k)/PMA/classification/recall, GRID), so the
% data sources appear in the paper's reference section. No other change from v6;
% no estimate, table, figure, or prose changes.
%
% Version 6 (2026-06-10): packaging/reproduction pass. Two table values change
% from v5, both data-vintage / display-only: (1) Table D.14 (class-specific
% summary stats) -- "FDA Administrative Data" adverse-event counts and the
% associated means/SDs updated to the freshly regenerated replication-package
% values (per-class N off by 1; e.g. 4,510,435->4,510,436 and 976,693->976,692);
% and (2) the 3_2 market-composition "excluding-zero" robustness table, matched-
% column Sample Size 1408->1188 (minimum-N-across-outcomes convention). Table 1
% is unchanged (the difference rounds away). No estimate, sign, significance, or
% conclusion changes; no prose changes. Done so the deposited tables match a
% fresh from-raw reproduction exactly (JPE Data Editor reproduction).
%
% Version 5 (2026-06-02): second-round R&R edits relative to the last JPE
% submission (v3; v4 was identical to v3 in substance). Changes: (1) Figure 1
% made black-and-white legible; (2) spillover language softened to "do not
% detect"; (3) added a market-age vs. utilization sentence in Sec. 6; (4)
% cost-benefit and other in-text numbers updated to the reproducible
% replication-package values. See my_paper/MANUSCRIPT_CHANGE_LOG.md and the
% one-page change note to the editor.

\usepackage{threeparttable}
\usepackage{geometry}
\geometry{letterpaper,tmargin=1in,bmargin=1in,lmargin=1in,rmargin=1in}
\usepackage[format=hang,font=normalsize,labelfont=bf]{caption}
\usepackage{amsmath}
\usepackage{multirow}
\usepackage{array}
\usepackage{delarray}
\usepackage{dsfont}
\usepackage{mathrsfs}
\usepackage{enumitem}
\usepackage{tikz}\usepackage[percent]{overpic}
\newcommand{\inputtight}[1]{\begingroup\renewcommand{\vspace}[1]{}\input{#1}\endgroup}
\usepackage{relsize}
\usepackage{amssymb}
\usepackage{rotating}
\usepackage{booktabs}
\usepackage{pdflscape}
\usepackage{pdfpages}
\usepackage{amsthm}
\usepackage{tabularx}
\usepackage{lscape}
\usepackage[lofdepth,lotdepth]{subfig}
\usepackage{graphicx}
\usepackage{natbib}
\usepackage{setspace}
\usepackage{float,color}
\usepackage{placeins}
\usepackage{listings}
\usepackage[autostyle, english = american]{csquotes}
\MakeOuterQuote{"}
\doublespacing
\synctex=1
\usepackage{xurl}
\usepackage{hyperref}
\hypersetup{colorlinks,linkcolor=black,urlcolor=black,citecolor=black}
\usepackage{bm}
\usepackage{etoolbox}

% put in preamble (no makeatletter needed)

\usepackage{adjustbox,makecell}
\newcolumntype{L}{>{\raggedright\arraybackslash}X}
\newcolumntype{C}{>{\centering\arraybackslash}X}
% helper for repeated estimator blocks with 4 numeric columns
\newcommand{\robblockIV}[2]{%
  \multicolumn{5}{@{}l@{}}{\emph{Estimator: #1}}\\
  #2
  \addlinespace[.5mm]\cmidrule(lr){2-5}
}
% helper for repeated estimator blocks with 5 numeric columns (Market Structure)
\newcommand{\robblockV}[2]{%
  \multicolumn{6}{@{}l@{}}{\emph{Estimator: #1}}\\
  #2
  \addlinespace[.5mm]\cmidrule(lr){2-6}
}

\AtBeginEnvironment{quote}{\singlespacing\small}
\newcommand{\E}{\mathbb{E}}

\newcounter{daggerfootnote}
\newcommand*{\daggerfootnote}[1]{%
    \setcounter{daggerfootnote}{\value{footnote}}%
    \renewcommand*{\thefootnote}{\fnsymbol{footnote}}%
    \footnote[2]{#1}%
    \setcounter{footnote}{\value{daggerfootnote}}%
    \renewcommand*{\thefootnote}{\arabic{footnote}}%
    }

\newtheorem{assumption}{Assumption}

\theoremstyle{definition}
\newtheorem{theorem}{Theorem}
\newtheorem{acknowledgement}[theorem]{Acknowledgement}
\newtheorem{algorithm}[theorem]{Algorithm}
\newtheorem{axiom}[theorem]{Axiom}
\newtheorem{case}[theorem]{Case}
\newtheorem{claim}[theorem]{Claim}
\newtheorem{conclusion}[theorem]{Conclusion}
\newtheorem{condition}[theorem]{Condition}
\newtheorem{conjecture}[theorem]{Conjecture}
\newtheorem{corollary}[theorem]{Corollary}
\newtheorem{criterion}[theorem]{Criterion}
\newtheorem{definition}{Definition} % Number definitions on their own
\newtheorem{derivation}{Derivation} % Number derivations on their own
\newtheorem{example}[theorem]{Example}
\newtheorem{exercise}[theorem]{Exercise}
\newtheorem{lemma}[theorem]{Lemma}
\newtheorem{notation}[theorem]{Notation}
\newtheorem{problem}[theorem]{Problem}
\newtheorem{proposition}{\normalfont\scshape Proposition} % Number propositions on their own
\newtheorem{remark}[theorem]{Remark}
\newtheorem{solution}[theorem]{Solution}
\newtheorem{summary}[theorem]{Summary}
\bibliographystyle{aer}
\newcommand\ve{\varepsilon}
\renewcommand\theenumi{\roman{enumi}}
\newcommand\norm[1]{\left\lVert#1\right\rVert}
\defcitealias{us2009fda}{GAO, 2009}
\defcitealias{NAP12960}{IOM, 2010}

\defcitealias{NIM:2011}{(IOM, 2011)}
\defcitealias{FR63:3142}{63 FR 3142 (1998)}
\defcitealias{FR60:38903}{60 FR 38903 (1995)}
\defcitealias{MartinVMedtronic:2001}{Martin v. Medtronic Inc. (2001)}
\defcitealias{KelseyVAlcon:2019}{Kelsey v. Alcon Laboratories Inc. (2019)}
\defcitealias{RiegelVMedtronic:2008}{Riegel v. Medtronic Inc. (2008)}
\defcitealias{MedtronicvLohr:1996}{Medtronic Inc. v. Lohr (1996)}





\begin{document}

% This version combines the changes in shorten second and shorten (figure changes). If don't like figure change up from shorten, revert back to shorten second

\begin{titlepage}
\title{Regulating the Innovators: \\ Approval Costs and Innovation in Medical Technologies}
% \title{Regulating Mature Markets: Impacts on Medical Technologies}


\author{Parker Rogers\thanks{Indiana University, Kelley School of Business. I thank the editor, John Asker, and the anonymous referees for thoughtful comments and suggestions that substantially improved the paper. Many thanks to Yue Chu, Yilan Jiang, and Yutong Wu for excellent research assistance. Thanks to Jeffrey Clemens, Josh Graff Zivin, Paul Niehaus, Craig McIntosh, Gordon Dahl, John Van Reenen, Heidi Williams, Matthew Grennan, Ariel Stern, Daniel Carpenter, Judson Boomhower, Itzik Fadlon, Karthik Muralidharan, Matthew Fiedler, Alexander Gelber, Mark Jacobsen, Aaron Kesselheim, Josh Makower, Jeffrey McCullough, Nicholson Price, Maggie Shi, Paul Kindsgrab, Ryan Hill, Andrew Johnston, Alex Everhart, Nathaniel Bechhofer, Yunan Ji, Max Edwards, and Rachel Olvera for excellent feedback. I also thank participants at the NBER Summer Institute Innovation Session 2023, NBER Economics of Health Fall Meeting 2023, NBER Economic Analysis of Regulation Conference 2023, FTC Microeconomics Conference 2023, BFI Health Economics Initiative Conference 2023, WEAI Graduate Student Workshop 2022, ASHEcon 2021, Young Economist Symposium 2021, SOCAE 2021, and UCSD Applied Micro. I gratefully acknowledge support from the National Institute on Aging, grant number T32-AG000186.
}}
\date{\today}





\maketitle

% \vspace{-24pt}

% \begin{center}
%      \href{https://parkerrogers.github.io/Papers/RegulatingtheInnovators_Rogers.pdf}{Click here for the latest version}. \\
%     % \textcolor{red}{\textit{Preliminary and incomplete.}}   
% \end{center}

\begin{singlespace}

\vspace{-24pt}

\begin{center}
\textbf{Abstract}
\end{center}


\noindent 



I examine how FDA regulation shapes innovation, market structure, and product safety by leveraging deregulation events that affected some established medical device types but not others. Linking FDA records, patents, and health insurance claims, I find that deregulation increases the quantity and quality of new technologies---especially among smaller, less-experienced firms---increases entry, and reduces medical procedure prices. Severe adverse events do not increase, and for mature technologies in higher-liability settings, they decline, consistent with postmarket measures (standards and liability) substituting for premarket review. These results clarify when targeted deregulation can expand innovation without sacrificing safety. (JEL I18, L51, O31)\\


% How does FDA regulation shape innovation, market structure, and product safety? I examine this question by exploiting FDA deregulation events that affected certain established medical device types but not others. A comprehensive analysis of data spanning innovation, safety, firm entry, prices, and regulatory changes reveals three key results.  First, deregulation significantly increased the quantity and quality of new technologies, particularly benefiting smaller and less regulation-experienced firms. Second, the events increased firm entry and reduced prices for medical procedures using deregulated devices. Finally, deregulation did not increase severe device-related injuries or deaths. Interestingly, some deregulated device types exhibit decreased adverse events, potentially due to firms increasing their emphasis on product safety in response to heightened litigation risk. (JEL I18, L51, O31)\\
% \noindent \textbf{JEL Codes:} I1, Health; O31, Innovation and Invention; L51, Economics of Regulation 
% K13, Tort Law and Product Liability
\end{singlespace}



\thispagestyle{empty}
\end{titlepage}




New technologies can both improve well-being and increase risk. Regulatory agencies such as the U.S. Food and Drug Administration (FDA) manage risk through premarket review---a compelling approach for high-risk, groundbreaking innovations \citep{Grennan:2020}. For lower-risk or mature technologies, where innovation is often incremental, the case for intensive review is less clear. Critics argue that stringent review raises entry costs and chills innovation \citep{peltzman1973evaluation}; proponents counter that it builds demand and fosters entry \citep{carpenter2010approval}. A key insight is that this trade-off is dynamic: as a technology matures, regulators and firms accumulate clinical evidence and refine safety standards, reducing the informational value of premarket review even as its costs persist \citep{henry2022regulation}. The question becomes not whether to regulate, but when lighter oversight can replace intensive premarket review without sacrificing safety. Given that FDA-regulated products comprise a \$2.8 trillion market \citep{FDAFactsheet2020}, empirical evidence on this question is important.

I advance this debate by measuring the impact of FDA regulation on innovation, market structure, and safety in established medical device types. My approach is twofold. I first consider instances of ``down-regulation,'' where moderate-risk devices, such as soft contact lenses, were moved from stringent Class III to lighter Class II review, replacing intensive premarket approval with standards-based review. Next, I analyze instances of ``deregulation,'' where lower-risk items like ventilator tubing transitioned from Class II to minimal Class I requirements, eliminating premarket review. This dual focus lets me compare the effects of stringent FDA protocols with those of lighter alternatives across a diverse range of products---from generic, low-risk items to differentiated, moderate-risk innovations---and estimate local average treatment effects at two distinct points on the regulatory spectrum.


To estimate the causal effects of reclassification, I employ a stacked difference-in-differences strategy. I compare affected device types to multiple control groups: devices with similar baseline characteristics, devices deregulated at a later time, and devices intuitively resembling the treated group. This approach yields consistent results across controls, with no evidence of divergent pre-existing trends. These results align with industry consensus that reclassifications are ``unpredictable'' \citep{Makower:2010, powell2018changing}, often occurring after long delays and opaque decision-making.


The down-regulation of daily-wear soft contact lenses provides a clear illustration of these dynamics. The FDA internally identified these lenses as candidates for reclassification in the mid-1980s, but key safety evidence was locked in proprietary premarket approval submissions and ``not publicly available for use'' (FDA Panel Minutes, October 1993), delaying reclassification until 1994. After down-regulation in 1994, annual patent filings linked to soft contact lenses rose more than tenfold; the number of unique firms active in this product category grew fivefold; and device submissions roughly quadrupled---with variety expanding from conventional hydrogel materials into silicone hydrogel formulations, advanced polymer coatings, and antimicrobial designs. One post-reclassification patent alone was valued at \$760 million (US-6478423-B1; \citealp{kogan:2017}), underpinning the market-leading ACUVUE brand. Extended-wear lenses, which remained in Class~III, show no comparable surge (Figure~\ref{fig:ContactLensCaseStudy}). This case captures the pattern I document across dozens of reclassifications: for mature technologies with well-understood risk profiles, lighter oversight spurs entry and innovation.

A distinct mechanism also operates through these regulatory transitions: the shift from ex-ante safety regulation to ex-post legal liability. FDA clearance historically provided manufacturers a compliance defense in tort; reclassification removes that safe harbor, raising expected liability exposure, which can reach up to 3.8\% of annual revenues \citep{MckinseyCostSafety:2018} and influence safety investments.\footnotemark\footnotetext{Legal liability has been shown to alter innovation strategies within the medical device sector \citep{galasso2017tort,Galasso:2018}. The financial toll can be heavy: the \$1 billion Stryker hip implant settlement and the \$3.2 billion Dow Corning breast implant settlement (2023 USD) illustrate the scale of exposure.} 


An important contribution of this paper is the assembly of novel, comprehensive data linking regulatory changes to market outcomes. Regulation affects outcomes across multiple domains, but data on these outcomes are siloed, disorganized, and unconnected to medical device types---limitations that have constrained prior research. I overcome these challenges by combining programmatic online text extraction, text analysis algorithms, and manual linkages to compile, merge, and harmonize the required data. This unified dataset spans four decades of regulatory changes affecting thousands of device types, as well as corroborative measures of innovation quantity and quality, firm characteristics, market structure, device pricing, and product safety.

Analysis of these data shows that down-regulation substantially increases the quantity and quality of innovation. Relative to control groups, moving from Class III (high regulation) to Class II (moderate) requirements increases patent filings and FDA submissions by 140--350\% under the preferred specification, and I do not detect spillover effects on similar device types.\footnotemark \footnotetext{These localized effects can be attributed to a high degree of specialization, where a significant number of device inventions stem from practicing physicians who operate within their highly specialized medical fields and actively employ these devices themselves (\citetalias{NAP12960}; \citealp{chatterji2016does}).}  Patents filed after down-regulation also garner more citations and higher market valuations, indicating improved technological quality. By contrast, the shift from Class II to Class I (deregulation) spurs a modest increase in the quantity and quality of patent filings. Across both pathways, the largest increases in innovation come from smaller, less regulatory-experienced firms, which have historically driven radical breakthroughs \citep{wu2019large}.

These regulatory events also reshaped market structure. Down-regulation spurred a 600\% increase in new entry (i.e., firms with no prior device submissions) and a fivefold increase in incumbent entry (i.e., firms expanding into newly down-regulated products). Meanwhile, deregulation significantly increased new entrants by 72\% without influencing incumbent entry. These market structure changes expanded product variety. Deregulation predominantly affected generic, low-risk devices where new entry compressed margins, yielding a 32\% decline in outpatient procedure prices.


The standard case for premarket review is that it prevents unsafe products from reaching patients. The preceding results show what regulation costs; the critical question is whether deregulation compromises safety. For deregulation, it does not---serious adverse events decline significantly and safety-focused patenting rises.\footnotemark \footnotetext{I show, using a subset of devices for which I have claims data, that this finding is robust to normalizing by utilization.} Down-regulation shows a modest, statistically insignificant rise in adverse events. For these mature technologies, accumulated market experience and postmarket enforcement mechanisms appear sufficient to maintain safety without intensive premarket review.



Two mechanisms help explain why safety is maintained---and, for deregulation, improved---after reclassification. First, a firm-level liability channel: when Class II requirements are lifted---viewed by the Institute of Medicine as insufficient to assure safety \citetalias{NIM:2011}---manufacturers lose a regulatory-compliance safe harbor in tort (a legal defense), raising expected liability and strengthening incentives to reduce injury risk. I test this using cross-firm variation in effective liability: smaller firms can more readily cap exposure via bankruptcy and thus face lower effective liability \citep{shavell1986judgment}, whereas larger firms cannot. Consistent with this channel, safety gains concentrate among larger firms. A placebo test reinforces the interpretation: a subsequent deregulation that did not alter firms' liability exposure shows no detectable change in safety. These patterns suggest that tort liability can substitute for formal premarket regulation in promoting safety, especially for generic or low-risk products. 

Second, a device-level lifecycle channel: the impact of reclassification depends on technological maturity. Early reclassification, when failure modes are still being discovered, can spur innovation but raise risk; once a device type is mature, reclassification lowers entry costs without compromising safety, enhancing innovation and improving safety. Many device types currently in Class III or II meet these maturity benchmarks, implying that standards-based oversight may outperform premarket authorization in those settings. The channels are complementary: liability pressure predicts where safety improves (across firms), while maturity predicts when it improves (across device types). Similar compliance-liability and lifecycle trade-offs appear in generic drugs, genetically modified foods, aviation, automobiles, and over 15{,}000 consumer products regulated by the CPSC \citep{schwartz2020s,schauzu2000concept,pisanifederal}.



My findings contribute to two broad literatures. First, I advance the growing literature on the effects of public policy on medical innovation.\footnotemark \footnotetext{See \citet{mulligan2021peltzman,Grennan:2020,clemens2020demand,stern2017innovation,budish2015firms,AcemogluLinn:2004,Finkelstein:2004}.} While the global medical device market is valued at \$500 billion and projected to reach \$1 trillion by 2030 \citep{stewart_2022}, relatively little empirical work examines how regulation shapes innovation in this industry. Unlike studies that focus on high-risk cardiovascular technologies \citep{stern2017innovation,Grennan:2020}, I exploit quasi-exogenous FDA regulatory shocks that span low- to moderate-risk device types and both high- and low-regulation environments. This design recovers not only innovation responses but also the underexplored safety impacts of regulation. 



Second, the paper contributes to the literature on regulatory design under uncertainty, particularly the trade-off between intensive ex-ante oversight and greater reliance on ex-post liability and monitoring \citep{Shavell:1984,shavell1986judgment,Kolstad:1990}.\footnote{See also \citet{ehrlich1974economic, glaeser2001coase, kessler2010regulation, shavell2018liability}.} More recent work shows that, as information about product risks accumulates over time, the optimal mix of ex-ante approval, ex-post withdrawal, and liability may shift toward lighter premarket control and greater reliance on postmarket mechanisms \citep{henry2022regulation,philipson2010effects}.\footnote{\citet{Henry:2019} also analyzes the conditions under which stringent approval processes are socially optimal.} Related work on management-based regulation highlights how compliance systems can support such regimes \citep{coglianese2003management}. I provide causal evidence on this trade-off using FDA reclassifications: deregulation need not reduce safety; its effects hinge on firm characteristics, product maturity, and liability risk. In “safe-candidate” categories---those with substantial clinical evidence---firms increase safety-related innovation post-deregulation. These findings suggest that the effects of deregulation depend on where a technology sits in its lifecycle and on the extent to which ex-post liability can substitute for premarket review \citep{galasso2017tort}. My results also speak to the broader effects of product regulation on market competition.\footnotemark \footnotetext{See \citet{buettner2006entry, aghion2009effects, aghion2005competition, djankov2006regulation, hahn1991costs} on regulation and competition, and \citet{acharya2014wrongful, acharya2013labor, vanreenan19} on labor regulation and innovation.} Deregulation disproportionately benefits firms with less prior regulatory experience, implying that regulatory knowledge remains localized rather than freely diffusing across firms---frictions that raise the effective costs of regulation and advantage experienced multiproduct incumbents.




This paper is organized as follows. Section \ref{sec:Background} provides background on the FDA regulatory process, Section \ref{sec:Concept} develops the conceptual framework, Section \ref{sec:Data} discusses my data, Section \ref{sec:Empirical} describes my empirical strategy, Section \ref{sec:Results} presents my empirical results, Section \ref{sec:RegulatoryDesign} presents implications for policy, and Section \ref{sec:Conclusion} concludes. 




\section{Background}\label{sec:Background}


This section describes the structure and legal consequences of FDA medical device regulations. Medical devices include products like COVID-19 tests, pacemakers, X-ray machines, and spinal implants.

\subsection{Enactment of Medical Device Regulations}



The 1976 Medical Device Amendments (MDA) extended the FDA's oversight to devices and organized products into \emph{generic device types} (e.g., “daily-wear soft contact lenses” vs. “extended-wear soft contact lenses”). The policy variation I study occurs at the level of these generic device types.


Device types fall into a three-tier risk system. Class I (low risk) generally requires facility registration alone. Class II (moderate risk) typically proceeds via the 510(k) ``substantial equivalence'' pathway; this pathway has been criticized as insufficient to assure safety \citetalias{NIM:2011} and, on average, costs firms \$24 million and delays commercialization by ten months \citep{Makower:2010}.\footnotemark \footnotetext{Estimates in \citet{Makower:2010} come from a survey of 204 predominantly early-stage firms reporting their aggregate FDA-related costs (see Appendix \ref{sec:MoreRegDetails}).} Class III (high risk) generally requires time-intensive premarket approval (PMA) supported by clinical evidence, costing, on average, \$75 million \citep{Makower:2010}. Figure \ref{pic:DeviceDeregulation} provides an illustration of these different pathways and Appendix \ref{sec:MoreRegDetails} provides more details.

\subsection{Reclassification of Medical Device Types}





The FDA may reclassify medical device types to a lower class as postmarket evidence clarifies risks.\footnotemark \footnotetext{Additionally, manufacturers can file a petition for reclassification, bringing the FDA's attention to particular device types for further investigation. My analysis, however, focuses on reclassification events explicitly enacted by the FDA's initiative (rather than a petition).} The FDA regulates new, markedly novel devices in Class III to ensure safety under unknown risks.\footnotemark \footnotetext{In 1997, the FDA began allowing manufacturers of markedly novel devices to petition for a direct Class II or I classification under the ``De Novo'' process by showing that best practices assure the safety and efficacy of their device. However, all the device types I consider existed before 1997 and thus were either automatically or intentionally classified into Class III.}  While it is probable that the FDA will deregulate some lower-risk device types within a class, the specific types affected and the timing of such events are difficult to anticipate \citep{Makower:2010,powell2018changing}.\footnotemark \footnotetext{See Appendix \ref{sec:BackgroundIIItoII} for additional details. Figure \ref{fig:event_distribution}, Panel (a), presents the distribution of these events over time.} My empirical analysis supports this assessment, as I do not find evidence of divergent pre-existing trends when comparing reclassified device types to control groups. 


By contrast, Class II to I deregulations mostly resulted from episodic, statute-driven agency initiatives. For example, following the Safe Medical Devices Act, the FDA reclassified a broad set of Class II types using a previously undisclosed adverse-event cutoff computed from pre-period averages \citep{FR60:38903} (Appendix~\ref{sec:FDADecisionRule}; Figure~\ref{fig:event_distribution}). Because the review was triggered by legislative timing and the cutoff referenced only baseline levels---not contemporaneous technology changes or trends---the timing across device types is plausibly ``as-good-as-random.'' A regression-discontinuity design is infeasible (too few types lie near the cutoff), so I compare treated types with unaffected Class I types---both previously deregulated and always-Class I---selected to have similar pre-period scores under the same rule. Event-study estimates show no differential pre-trends, supporting this identification. These reclassifications occur in established device types; accordingly, the estimates speak to improvements in existing technologies rather than the emergence of radically new ones.




\subsection{Regulation versus Litigation: Federal Preemption}

U.S. device manufacturers face substantial product-liability exposure---estimated at up to 3.8\% of revenues \citep{MckinseyCostSafety:2018}---which can chill innovation and threaten smaller firms \citep{Galasso:2018}. The U.S. legal environment (class actions, punitive damages, limited caps) heightens this risk \citep{Guendling2016}, creating a “powerful incentive for improving product safety” (Bravman v.\ Baxter Healthcare Corp., 1994).\footnotemark \footnotetext{Bravman v. Baxter Healthcare Corp., 842 F. Supp. 747, 761 (S.D.N.Y. 1994).}


Regulatory compliance, however, can limit tort exposure. In \citetalias{RiegelVMedtronic:2008}, the Supreme Court held that PMA (Class III) preempts many state-law claims. In the 1990s, many 510(k) (Class II) devices also received substantial protection \citep{flaherty2008defending}, whereas Class I devices---lacking FDA approval---do not. Some Class II devices remain protected today; Appendix~\ref{sec:Placebo} documents the evolving litigation environment.




\section{Conceptual Framework}\label{sec:Concept}

I build a two-stage model (formal details in Appendix~\ref{sec:AppendixConcept}) in which firms first invent and finance development, then commercialize and choose safety effort. Reclassification shifts three primitives that govern these choices: the approval delay, the mandated safety level, and effective exposure to tort. Under Class~III/II (regulated), firms face positive approval delays, must meet a mandated level of safety effort, and enjoy a regulatory-compliance safe harbor in tort that lowers expected liability. Under Class~I (deregulated), firms avoid approval delays, choose safety effort to maximize profits, and lose the compliance defense, raising expected liability.


Delays shorten the effective life of innovations and raise external financing needs, especially for inexperienced and capital-constrained firms. Easing review through deregulation, therefore, increases R\&D and entry, with disproportionately strong responses among smaller or less regulatory-experienced firms.


Product safety responds differently across regimes because, historically, effective liability changes only after deregulation. When Class~II requirements are lifted, firms lose the compliance defense, and expected liability increases, potentially enhancing safety effort. If deregulation increases safety effort, the magnitude of this response is heterogeneous: when a firm’s seizable assets are low relative to worst-case damages (partly “judgment-proof”), the marginal return to safety is attenuated and safety effort increases less, while larger, better-capitalized firms face higher effective exposure and increase safety effort more. By contrast, down-regulation (III to II) leaves the compliance defense in place and, thus, produces little systematic change in safety.



Although not formalized in the model, device-level technological maturity complements this firm-level mechanism. Early in the lifecycle, unknown failure modes make safety effort hard to target, lowering its marginal return and increasing the informational value of premarket testing. Once failure modes are understood and codified in standards, easing review can reduce delays and financing costs without compromising safety.



This framework delivers three predictions borne out in the data: (i) both forms of reclassification increase innovation and entry via the delay/finance channel, especially among smaller and inexperienced firms; (ii) safety improves after deregulation (II to I) because expected liability rises when the compliance defense disappears, with larger effects at better-capitalized firms; and (iii) these safety gains are most pronounced for mature device types, where the informational value of premarket review is lower.

\section{Data}\label{sec:Data}

To conduct my empirical analysis, I compile data from eight sources to provide an expansive view of the costs and benefits of medical device regulations. Summary statistics for these data are provided in Table \ref{tab:SummaryStatistics}, and a data catalog is presented in Figure \ref{fig:DataCatalog}.

\textit{FDA Reclassification Events.} To conduct a comprehensive analysis of FDA reclassification events, I compile a complete record of all reclassifications from 1980 to 2015.\footnotemark \footnotetext{Table \ref{tab:ClassIIIRegnums_combined}, Panel A enumerates the device types in my analysis that were down-regulated, including Federal Register citations. Appendices \ref{sec:BackgroundIIItoII} and \ref{sec:FDADecisionRule} elaborate on the deregulation events and cite sources listing the affected devices.} For Class III to II events, I also indicate whether the event was motivated by the FDA's ``own initiative'' or in response to an industry petition, as this distinction is empirically important.\footnotemark \footnotetext{Of the 29 Class~III to~II reclassifications, 12 were FDA-initiated and 17 were petition-driven. Firm lobbying in the reclassification process operates primarily through this petition mechanism, whereby manufacturers formally request the FDA to reclassify a device type. Excluding petition-driven events from the main Class~III to~II sample addresses the primary channel through which firm-initiated advocacy could introduce endogeneity.} Figure \ref{fig:Petitioned_IIItoII} shows that device types subject to a petitioned reclassification demonstrate a notable jump in patent filings in the year of the announcement, with long-run impacts more than double those observed in non-petitioned events.\footnotemark \footnotetext{Table \ref{tab:ClassIIIRegnums_combined}, Panel B enumerates the petitioned device types that were down-regulated, while Appendix \ref{sec:IndustryPetition} provides more information on the petitioned events.} The Class III to II events I consider are those enacted by the FDA's own initiative where the affected device types had at least one PMA submission or patent filing prior to reclassification.\footnotemark \footnotetext{Many Class III ``preamendment'' devices were never officially required to submit PMA documentation and had no pre-reclassification patenting activity; excluding them ensures a measurable pre-treatment baseline.} For Class II to I events, I consider affected device types with at least one pre-reclassification patent filing (Table~\ref{tab:SampleConstruction}).

\textit{FDA Device Submissions (PMA and 510(k) Databases).} I derive a primary dataset from FDA administrative data on the universe of medical devices submitted for approval/clearance. These data are drawn from the FDA's PMA and 510(k) databases, covering both Class III and II devices. Submissions include the submitting company name, device brand name, medical device type, and submission and approval dates. I use fuzzy matching to form three measures of market dynamics and innovation. First, I measure ``new entry'' by identifying firms submitting approval documents for the first time. Second, I form a measure of ``incumbent entry'' by locating firms that have filed prior approval documents but are starting to submit for approval in a given device type. Third, I isolate the first occurrence of unique device brand names within a device type to form the ``unique device submissions'' measure. These variables are aggregated to the device-type-year level. To measure each firm's regulatory proficiency, I calculate the total approval delays (in days) the submitting firm has experienced up to the date of each submission. 



\textit{FDA Adverse Event Reports (MAUDE).} The FDA’s Manufacturer and User Facility Device Experience (MAUDE) database contains roughly 9.2 million device-related reports. I construct a device-type-by-year series of adverse event reports from 1992 to 2019. I focus on “serious” outcomes---deaths, hospitalizations, and life-threatening events---for which manufacturers face mandatory reporting obligations enforced through criminal liability \citep{bragg_florence_chan_2018}. To study heterogeneity by firm size, I link MAUDE-reporting manufacturers to CRSP/Compustat to obtain asset totals for public firms, working down the report-count ranking: the 300 largest reporters account for 84\% of all reports, of which 88 are private firms with no public-firm match and the remainder are linked by automated name clustering with hand-coded corrections. I construct three complementary safety measures, each suited to a different purpose in the analysis. First, simple annual counts of serious events retain the full treated sample and avoid denominator noise; I use these for baseline event-study figures as a benchmark, while acknowledging that levels are sensitive to utilization and reporting intensity (see Appendix \ref{appendix:safety} for more details and robustness checks). Second, where reliable procedure mappings exist, I normalize by exposure---serious events per linked procedure---so that estimates reflect risk conditional on use rather than shifts in volume or case mix; these exposure-adjusted rates serve as a robustness check for treated groups that have linked procedures. Third, I report a “severity share”---serious events divided by total MAUDE reports---which absorbs broad shifts in reporting intensity and utilization to the extent they affect both numerator and denominator. Because reductions in serious events must outpace reductions in overall reports to manifest as product safety improvements, this is a conservative metric that I deploy for cross-device-type comparisons (e.g., by device age) to maintain the full treated sample while aiming to account for utilization. Across these measures, the conclusions are consistent.

\textit{Dimensions Patent Grants Extract.} 
Patents complement the ``unique device submissions'' metric as a measure of innovation. Unlike pharmaceuticals, which have the Orange Book to link drugs to patents, no equivalent database exists for medical devices. To construct this linkage, I use \textit{Dimensions}---a tool for programmatic patent text searches \citep{Dimensions}---to identify relevant USPTO-granted patents. The process involves extracting keywords from FDA device type descriptions and using them to retrieve patents containing these terms. The resulting patent dataset contains over 1.2 million medical device patents from 1976 to 2019. I then compute annual patent counts per device type (i.e., ``patenting rates'') based on initial filing dates. This approach yields a panel of yearly patent counts for over 5,000 FDA-defined device types from 1976 to 2019. Patents serve as a valuable innovation measure, enabling analysis of how Class II to I deregulation affects innovation---a margin not captured by the ``unique device submissions'' metric. They also allow for direct comparisons between down-regulation and deregulation. Moreover, the consistency of results across patent filings and device submissions, as shown in Section~\ref{sec:Results}, reinforces the robustness of the findings, demonstrating that the observed increase in innovation is not dependent on a specific metric. Appendix~\ref{sec:appendixPatentCollection} provides further details on the patent collection process and robustness checks.




\textit{Patent and Patent Applicant Characteristics.} I enrich the patent data with measures of innovation quality and applicant characteristics. Patent quality is assessed using forward citations from other patents and patent market value.\footnotemark \footnotetext{I exclude examiner citations and assign zero values to patent citations and market values when no patents were filed in a given device-type-year.} Market values, derived from \citet{kogan:2017}, are expressed in millions of 2019 dollars (CPI-deflated) and are based on the increase in a patent assignee’s stock price following a USPTO issuance announcement, making them available only for publicly traded firms. I also construct \textit{emphasis on safety}---a device-type-year indicator equal to one if at least one patent filed that year explicitly references safety features---following \citet{clemens2020demand}.\footnotemark \footnotetext{Explicit safety mentions are rare, so I code the measure as binary. See Appendix~\ref{appendix:safety} for details and robustness checks.} This measure allows me to assess how reclassification influences inventors' focus on improving device safety, complementing the adverse event analyses. Lastly, to examine how reclassification affects innovation across firms of different sizes, I link patent applicants to total firm asset holdings from the CRSP/Compustat database.


\textit{Truven Health MarketScan Database (``MarketScan'').} MarketScan is a monthly panel of employer-sponsored commercial claims from 1996 to 2013. Because no nationally representative source reports hospital device purchase prices before 2011,\footnote{ECRI's \textit{PriceGuide} offers a rolling 10-year window of hospital purchase prices, limiting historical coverage.} I use MarketScan outpatient paid amounts---adjudicated under the Healthcare Common Procedure Coding System (HCPCS)---as a proxy for device prices, since device costs are typically bundled into the procedure payment.\footnote{NIH, \emph{Reimbursement Knowledge Guide for Medical Devices}, pp.~10, 14, \url{https://seed.nih.gov/sites/default/files/2024-01/Reimbursement-Knowledge-Guide-for-Medical-Devices.pdf}.} To estimate the price effects of FDA reclassification, I link FDA product codes to outpatient procedures with a large language model (GPT-5-mini). I instructed the model with strict guardrails, including linking only when the device defines the procedure or materially drives payment.\footnotemark
\footnotetext{Additional guardrails included: (i) outpatient-only scope; (ii) analyte-specific mappings for in vitro diagnostics; and (iii) hard exclusions for E/M, visit/collection/handling, and unlisted/miscellaneous HCPCS families (e.g., 992\textit{xx}). Each proposed linkage required a brief rationale and supporting source(s).}
I then benchmark the LLM-based device-procedure concordance against an independently produced medical-coder concordance (two AAPC-certified coders with third-coder reconciliation) and find it more accurate.\footnotemark
\footnotetext{Accuracy was adjudicated by a three-model panel (GPT-5 Thinking, Gemini~2.5~Pro, Claude~Sonnet~4.5) applying the same rules and selecting the better mapping in blinded A/B tests. Across the compared device types, the panel scored GPT-5-mini higher. Estimated accuracy (correct links / total links provided) from Claude~Sonnet~4.5: GPT-5-mini = 80\%, expert coders = 75\%. GPT-5 Thinking judged GPT-5-mini correct on 78/100 items vs.\ expert coders 43/100; Gemini~2.5~Pro judged 85/100 vs.\ 48/100. I therefore adopt the LLM concordance at baseline and report robustness to human encodings.} However, the results are robust when using the human encodings. The final concordance (Table~\ref{tab:DeviceToProcedureLinkages}) covers 15 deregulated device types (Class II to I)---a subset of the 56 device types reclassified in the 1998 event, the only deregulation cohort with pre-period claims data (Table~\ref{tab:SampleConstruction})---and 3 of the 12 down-regulated device types (Class III to II); the remaining device types either lack a clear outpatient-procedure mapping or are linked to procedures not observed in the MarketScan claims. For controls, I use (i) CPT procedures matched to treated procedures on pre-event mean prices from the full CPT universe and (ii) a second matched set restricted to procedures tied to device-type controls used in the innovation and market-structure analysis. I then aggregate to an annual panel of mean paid amounts by linked device type to measure price dynamics around reclassification. The final analytic sample comprises approximately 17,400 procedure codes and 182,200 procedure-code-by-year observations.






\section{Empirical Strategy}\label{sec:Empirical}
I estimate the effects of reclassification using stacked difference-in-differences and event-study designs. After describing each design, I underscore how I address potential issues when generating causal estimates in this context.

The first regression specification uses a staggered difference-in-differences design. I estimate a ``stacked'' regression, following \citet{cengiz2019effect}, which avoids potential biases from staggered treatment timing when treatment effects are heterogeneous \citep{goodman2018difference,de2019two}. The results are robust to alternative estimators and outcome transformations.\footnotemark \footnotetext{I re-estimate the main specifications using the imputation estimator of \citet{borusyak2021revisiting} (in levels), a log-transformed specification for count outcomes ($\log(y+1)$), and a quasi-Poisson model. Robustness results for the three broad outcome categories---Innovation (Table~\ref{tab:Robustness_Innovation}), Market Structure (Table~\ref{tab:Robustness_MarketStructure}), and Adverse Events (Table~\ref{tab:Robustness_Adverse})---are reported in Appendix Tables.} This approach assembles event-specific panel data, which includes a given treated unit $r \in \{1,\dots, N^1\}$ and all admissible controls for that unit. Then, all event-specific panels are stacked while allowing unique time and group fixed effects for each panel. Thus, the estimating equation is given by



\begin{equation}
Y_{t,c,r}=\gamma_{c,r} + \gamma_{t,r} + \beta_1 1\text{\{reclass\}}_{t,c,r} + \varepsilon_{t, c,r} \label{eq:OLS}.
\end{equation}


In Equation \ref{eq:OLS}, $c$ denotes the medical device type, $t$ denotes time, $r$ denotes the event, and $1\text{\{reclass\}}_{t,c,r}$ is an indicator equal to one when reclassification has occurred in device type $c$. The outcomes of interest are denoted by $Y_{t, c,r}$. Event-by-time fixed effects ($\gamma_{t,r}$) and event-by-device-type fixed effects ($\gamma_{c,r}$) are included. The coefficient of interest, $\beta_1$, estimates the differential change in the outcome variable for treated device types relative to control device types after reclassification. I estimate Equation \ref{eq:OLS} separately for Class III to II events and Class II to I events. 


The number of FDA-initiated Class III to II events is relatively low ($N^1=12$). Thus, I follow \citet{conley2011inference}, who provide a method of constructing reliable confidence intervals for difference-in-differences estimates in the presence of a small number of policy changes. This approach uses information from control group residuals to form confidence intervals.\footnotemark \footnotetext{Conclusions do not hinge on Conley--Taber standard errors; estimates and significance patterns are similar under the \citet{borusyak2021revisiting} imputation and quasi-Poisson specifications (Tables~\ref{tab:Robustness_Innovation}, \ref{tab:Robustness_MarketStructure}, \ref{tab:Robustness_Adverse}).}




Like all difference-in-differences designs, my specification relies on the assumption that differential trends in the outcomes of interest do not pre-date the reclassification events. To test this assumption, I estimate a stacked event-study design using OLS, given by

\begin{equation}
Y_{t,c,r}=\gamma_{c,r} + \gamma_{t,r} + \sum_{k(t,c) \not = -1} \beta_{k} 1\text{\{Treated\}}_{c,r} \times I_{k(t,c)} + \varepsilon_{t,c,r} \label{eq:EventStudy}.
\end{equation}

In Equation \ref{eq:EventStudy}, $1\text{\{Treated\}}_{c,r}$ is an indicator for whether device type $c$ in event panel $r$ has been subject to reclassification.\footnotemark \footnotetext{Note that for treated units, there is only one device-type fixed effect; however, control device types may appear in multiple event panels $r$, resulting in multiple device-type-event fixed effects. For example, the outcome for the number of patents filed for daily-wear soft contact lenses in 1994 (a treated group) appears once and is associated with one device-type fixed effect, whereas the outcome for extended-wear soft contact lenses in 1994, used as an intuitive control, appears in multiple panels corresponding to different treated events, each with its own device-type-event fixed effect.} $I_{k(t,c)}$ are indicators for years relative to the reclassification. I define $k(t,c) = 0$ as the year of reclassification. Thus, coefficients $\beta_{k}$ represent the difference-in-differences estimate of the change in the outcome in a given period relative to the year before reclassification (i.e., the omitted year $k(t,c) = -1$). Standard errors for each $\beta_{k}$ are calculated using \citet{conley2011inference}. 

Reclassification rulings for Class III to II events are typically announced two years before enactment, while the Class II to I events analyzed here were announced one year before enactment. Because innovators could respond to a reclassification announcement, $1\text{\{reclass\}}_{t,c,r}$ equals one for forward-looking outcomes (e.g., patenting) in all device-type-years after an announcement occurs in device type $c$. However, FDA administrative data will not reflect changes until enactment, because firms cannot introduce devices under new regulations beforehand. Thus, for FDA-derived outcome data, the indicator $1\text{\{reclass\}}_{t,c,r}$ is equal to one for all device-type-years after a reclassification is enacted in device type $c$. For the event study, the event-time $k(t,c) = 0$ follows accordingly.


Identifying control device types that track the counterfactual development of the outcome variables is a central challenge in my empirical context. Controls could be unsuitable for several reasons: they could be affected by unique scientific developments, have lower scientific potential, or face different market forces. Additionally, some device types could be affected by spillovers from treated device types. Lastly, the FDA selects device types for reclassification based on inherent risk. Thus, reclassified devices may be less dangerous than those not chosen.

I construct three distinct control groups to capture different counterfactual scenarios and address various sources of potential bias. The first control group consists of “later-treated” device types that were reclassified only after the treated device types and after the latest sample year.\footnotemark \footnotetext{Specifically, for Class III to II events, I gather controls from all Class III to II events that occurred after 2015, censoring the outcome data after 2015. For Class II to I events, I use device types moved from Class II to I in late 2019. The 21st Century Cures Act drove this Class II to I event and was the first time FDA-initiated reclassifications of Class II devices occurred since 1998 (the year of the event I analyze). Importantly, the FDA used the same explicit reclassification criteria in both events.}
Because the FDA ultimately deemed these later-treated devices suitable for the same kind of reclassification, they can be viewed as “on the margin.” By comparing earlier-treated devices to these borderline reclassifications, I approximate a scenario in which the earlier-treated devices remained at their higher classification somewhat longer, assuming their outcomes would have evolved similarly.

Second, I include an “intuitive” control group, composed of device types that treat similar ailments and exhibit comparable risk profiles to those in the treatment group. By focusing on devices within the same medical domain, I control for market-level changes (e.g., shifts in demand or clinical practice) that might otherwise confound the effect of reclassification. Although these intuitive controls could plausibly be exposed to reclassification spillovers, I find no evidence of such effects.\footnotemark \footnotetext{See Table \ref{tab:Innovation_Spillovers} for spillover estimates.} Tables \ref{tab:ClassIIIRegnumswControls} and \ref{tab:ClassIIRegnumsTreated} provide detailed profiles of the intuitive control groups and the treated devices. 


Finally, I construct ``matched'' control groups. For each reclassification pathway, I match each treated product code to its two nearest never-reclassified codes using 5-year pre-event means.\footnotemark  
\footnotetext{Implementation details: For each treated product code, I compute 5-year pre-event means (relative to its reclassification year) and select the two nearest-neighbor controls from the treated group's pre/post regimes (III to II: 510(k)/PMA; II to I: 510(k)/510(k)-exempt). All matched units have at least one pre-period patent. Distance metrics and covariate sets are: (i) \emph{Adverse-event analysis}---Mahalanobis distance on standardized pre-event means of all considered outcomes; using all outcomes guards against noise in pre-event adverse-event rates. (ii) \emph{Innovation/market-structure analyses}---squared Euclidean distance on raw pre-event means of patent counts and new firms patenting, augmented with device-submission counts for Class III to II. (iii) \emph{Patent-quality analysis}---squared Euclidean distance on raw pre-event means of forward-citation rates and real patent-value measures.} Although I do not find evidence of spillovers, I exclude device types that treat the same indications as the treated code. Because many outcomes (e.g., patenting) naturally trend upward over time, aligning pre-event levels via matching improves comparability; accordingly, this matched design is my preferred specification, though results are similar under alternative control groups.\footnotemark \footnotetext{Each event panel contains one treated group and its own matched controls; controls can reappear across panels. In inference, I account for such reuse when clustering standard errors.}







My findings are robust to excluding device types that never exhibit adverse events or innovative activity across the entire sample period.\footnotemark \footnotetext{Table~\ref{tab:Combined_DropNoCounts} restricts the sample to device types with at least one positive count of a given outcome. Because this restriction conditions on the full panel---including post-treatment years---it could, in principle, select on treatment-affected outcomes; these estimates therefore complement the full-sample baseline rather than replace it.}

As with any non-experimental research design, selection into treatment is a primary concern. Because the FDA's reclassification of medical device types is based on baseline adverse event rates, I cannot ascertain how reclassification would impact adverse event rates for a randomly chosen device type.\footnotemark \footnotetext{Importantly, the FDA does not choose device types based on innovative potential or the potential for market structure effects. See Appendix \ref{sec:FDADecisionRule} for more details.} Consequently, my estimates represent local average treatment effects among lower-risk device types within a given class.

\section{Results}\label{sec:Results}
 
 This section follows the paper's central trade-off. I first show that lighter oversight increases innovation (Section~\ref{sec:innovationchange}). I then show that it reshapes market structure and prices (Section~\ref{sec:marketcomposition}), with heterogeneous effects across firms (Section~\ref{sec:heterogeneityresults}). I finally ask whether these gains come at a safety cost (Section~\ref{sec:adverseeventresults}). Throughout, I report the preferred matched-control estimates in the text and treat alternative specifications as supporting evidence in footnotes and the appendix.


 \subsection{Changes in Innovation}\label{sec:innovationchange}
 Table \ref{tab:Innovation} reports estimates of Equation \ref{eq:OLS} for the innovation outcomes.\footnotemark \footnotetext{Table~\ref{tab:Combined_DropNoCounts} restricts to device types with at least one positive count of the relevant outcome across the full panel; see Section~\ref{sec:Empirical} for discussion. Note that sample sizes differ across Tables \ref{tab:Innovation}, \ref{tab:MarketComposition}, and \ref{tab:AdverseEvent} due to variations in the availability of associated patents, medical procedures, and adverse events for each device type.} Panel A presents estimates for Class III to II events, and Panel B provides estimates for Class II to I. I focus on the preferred matched-control specification (Col.~2); results are consistent across the intuitive and later-treated controls (Cols.~3--4).\footnotemark \footnotetext{Column~(1) reports the treated group's 5-year pre-event mean. Conley-Taber standard errors are reported below the coefficients throughout.}


Table \ref{tab:Innovation}, Panel A indicates that down-regulation led to statistically significant increases in patenting rates, unique device submissions, forward citations per patent, and patent market values. The matched-control estimates imply approximately 140--350\% more patents and device submissions per device-type-year (baseline means: 8 patents/yr; 1 submission/yr). Post-event patents receive nearly three times as many forward citations and exhibit double the market value. Panel B reports smaller but economically meaningful responses for deregulation: patenting increased by 62\%, forward citations by 87\%, and market value increased modestly, though these effects are not statistically significant under the matched-control specification.


I study dynamics by estimating the event-study specification in Equation \ref{eq:EventStudy}. Figure \ref{fig:InnovationES} plots the event-time coefficients $\beta_{k}$ for down-regulation and deregulation events using the ``matched'' control groups. The results of this analysis provide several insights for interpreting my findings. First, the estimated leads ($k<0$) are small and statistically indistinguishable from zero over the decade prior to reclassification. This insight strengthens the identifying assumptions that (i) treatment and control groups would have exhibited similar trends in outcomes absent the policy change,\footnotemark \footnotetext{Event-study coefficients are similar to those estimated using the intuitive and later-treated controls (not shown).} (ii) policies were not anticipated, and (iii) policies were not endogenous to increases in innovative activity. Second, post-event, $\beta_k$ rises gradually and remains elevated, consistent with reclassification inducing new R\&D investment (with gestation lags) rather than merely pulling filings forward in time.

  

Lastly, Figure \ref{fig:InnovationES} reveals different dynamics across the two innovation margins. Device submissions spike immediately---consistent with firms commercializing existing technologies once regulatory barriers fall---while patenting rises gradually, reflecting R\&D gestation lags.\footnotemark \footnotetext{Existing technologies marketed anywhere in the world are generally already patented: international filing conventions require firms to file within 12 months of an initial application \citep{popp2005using}, so reclassification cannot unlock a backlog of unpatented but commercialized devices. Several mechanisms can generate the early spike in device submissions: (i) firms commercialize “on-the-shelf” products previously stalled by costly FDA approval; (ii) existing technologies are repurposed to new indications; (iii) streamlined review lowers frictions for products at various stages; and (iv) more lenient E.U. rules previously enabled firms to bring E.U.-approved devices to the U.S. after reclassification \citep{Grennan:2020}.}




\subsection{Changes in Market Structure}\label{sec:marketcomposition}

This section examines how reclassification reshapes market structure through four channels: (i) overall firm entry, (ii) incumbent expansion, (iii) procedure prices, and (iv) product variety.


I begin by estimating Equation~\ref{eq:OLS} on two key outcomes: the number of new firms entering each device category (i.e., ``new entry'') and the number of incumbent firms expanding into reclassified device types (i.e., ``incumbent entry''). Table~\ref{tab:MarketComposition} presents the difference-in-differences estimates.\footnotemark \footnotetext{Table~\ref{tab:Combined_DropNoCounts} presents a robustness check that restricts to device types with at least one nonzero observation of a given outcome across the entire panel. Many device types are structurally zero on certain outcomes---for example, types that never generate a patent filing or adverse event report in any year---and including them may attenuate the estimated effects. Because this restriction is defined over the full panel, including post-treatment years, it could in principle select on treatment-affected outcomes. I therefore present these estimates as a complement to the full-sample baseline, not a substitute.}

Panel~A shows that down-regulation (Class III to II) leads to a sharp rise in both new and incumbent entry. New-firm entry significantly increases by 600\% in the FDA submissions data (baseline: 0.1 firms/yr) and by 164\% in the patent data. Incumbents also expand into newly down-regulated device types, though less than the inflow of new entrants.

Panel~B indicates that deregulation (Class II to I) generates a more modest but statistically significant 72\% increase in new-firm entry in the patent data. Increases in incumbent entry are not statistically significant.

Event-study estimates in Figure~\ref{fig:MarketStructure} reinforce these patterns. Effects, when present, are persistent, and pre-treatment coefficients are flat, supporting the identifying assumptions. Panels~(a)--(b) highlight a gradual increase in new-firm entry, consistent with R\&D and commercialization lags; additional event-study plots for incumbent and new entry appear in Figure~\ref{fig:MktCompostition_IIItoII_Patents}.




% A natural question is whether the observed entry patterns dilute the market presence of previously dominant incumbents. To assess this, I examine changes in patenting rates of incumbents (those previously patenting in an affected category) and newcomers before and after reclassification. Figure~\ref{fig:MktCompostition_IIItoII_Patents_Incumbent} shows that, following down-regulation and deregulation, incumbents experience a relatively quick uptick in patent filings suggesting they promptly capitalize on lower hurdles to bring existing R\&D projects to market. Over the medium term, however, increases in new entrant patenting approaches and eventually outpaces incumbent patenting rate increases, implying that incumbents’ share of patenting declines as more firms enter. This result is consistent with the high new-firm entry effects in Panel~A of Table~\ref{tab:MarketComposition}. These results suggest that despite intensified competition, reducing regulatory barriers boosts innovation incentives across the board, with newcomers particularly advantaged. While incumbents lose some patenting share relative to entrants, they remain active innovators.

Having established effects on entry, I turn to medical procedure prices linked to reclassified devices. The top row of each panel in Table~\ref{tab:MarketComposition} reports difference-in-differences estimates for log procedure prices.\footnotemark \footnotetext{Results for deregulation are robust to using a human-expert device-type-procedure concordance that allows less conservative linkages and identifies $52$ treated device types. Holding the control construction fixed at the baseline, the deregulation estimate is $-0.18$ (SE $0.07$), slightly attenuated but statistically significant. For down-regulation, the expert concordance yields an estimate of $-0.09$ (SE $0.03$), compared with an insignificant baseline estimate. Unlike deregulation, however, the estimated price effect of down-regulation is not consistently signed across concordances.}

Panel~A shows no economically meaningful change following down-regulation: the preferred specification yields a 0.06 log-point increase in prices, which is statistically indistinguishable from zero. The event study in Figure~\ref{fig:MarketStructure}, Panel~(c), plots a possible post-event temporary increase that subsequently fades, consistent with short-run price adjustments rather than a persistent level shift.


By contrast, Panel~B indicates a large and statistically significant 32\% decline in procedure prices following deregulation. The preferred specification matches treated procedures to control procedures drawn from the full pool of CPT codes (``All CPTs'').\footnotemark \footnotetext{Results are similar when restricting to the matched device-type controls used in the innovation and market structure analyses (``Matched'' column in Table~\ref{tab:MarketComposition}).} Event-study estimates (Figure~\ref{fig:MarketStructure}, Panel~(d)) show an immediate drop that deepens over time, consistent with price renegotiation dynamics in medical technology markets \citep{reinhardt2006pricing,grennan2020transparency}. Overall, the evidence suggests that competition induced by deregulation lowers associated procedure prices, whereas down-regulation leaves prices relatively unchanged.\footnote{The claims data do not permit a full decomposition of the price decline into supply-side competition versus shifts in product mix toward lower-cost devices; disentangling these channels remains an important direction for future work.}

I also examine whether reclassification affects utilization of linked medical procedures. Using the same MarketScan claims data, I estimate Equation~\ref{eq:OLS} with log annual procedure counts as the outcome. Utilization estimates are sensitive to the choice of control group and not stable across specifications (Appendix Table~\ref{tab:Utilization}): the preferred specification yields a positive and significant estimate for deregulation, but the matched-control specification does not. Down-regulation shows no significant effect under either approach. This instability is consistent with utilization being largely determined by clinical indication; because insurance substantially insulates patients from procedure prices, the price declines documented above have limited scope to shift demand.

Did reclassification simply produce more patents in the same technology space, or did it expand the scope of inventive activity? To assess this, I use LDA topic modeling on patent texts to identify dominant technology dimensions within each device type and estimate treatment effects separately by topic (Figure~\ref{fig:ProductVariety}).\footnotemark \footnotetext{Specifically, I train an LDA model on patent titles and abstracts linked to treated device types for each pathway, choosing the number of topics by maximizing a coherence criterion. Each patent is assigned a dominant topic to construct a topic-device-year panel. For each treated device type and topic, I match to the never-reclassified device type whose pre-event count in that topic is closest and estimate via Equation~\ref{eq:OLS}. Preprocessing includes lower-casing, tokenization, stop-word filtering, lemmatization, bigram/trigram detection, and dictionary trimming.} Panel~(a)---Class~III to II---shows that post-event patenting growth appears across most estimated topic clusters rather than concentrating in one. Panel~(b)---Class~II to I---exhibits a similar pattern. In both pathways, reclassification expands the scope of patenting activity rather than simply intensifying it within a single technology area.




In sum, reclassifying devices from Class~III to II induces significant entry by both incumbents and new firms. Procedure prices remain largely unchanged in this pathway. Moving from Class~II to I also stimulates new-firm entry and is associated with substantially lower procedure prices. In both settings, product variety widens. These patterns illustrate the regulatory trade-off: premarket review generates safety information, but at the cost of deterring entry, limiting variety, and---for generic devices---sustaining higher prices.




\subsection{Heterogeneity by Firm Size and Regulatory Proficiency}\label{sec:heterogeneityresults}

Average treatment effects can mask important cross-firm differences. If the costs of premarket review---regulatory expertise, financing, time---fall unevenly across firms, the gains from lighter oversight should as well. I therefore estimate Equation~\ref{eq:OLS} separately by (i) firm size and (ii) regulatory proficiency, linking results to the conceptual framework in Section~\ref{sec:Concept}.

 
\textit{Regulatory Proficiency.} I proxy proficiency using firms’ cumulative exposure to FDA approval delays up to each event date and form terciles of this measure. I then re‐estimate device submission responses within tercile bins and report percent changes relative to pre‐event means (Figure~\ref{fig:FirmExperienceHeteroSimulation}, Panel~(a)).\footnotemark \footnotetext{Baseline submission rates are similar across terciles; percent changes normalize each tercile's estimate by its pre-event baseline, facilitating comparison across bins.} Down‐regulation raises new device submissions across all terciles, but the response is steeply decreasing in proficiency: least‐experienced firms exhibit roughly a $280\%$ increase, compared with about $65\%$ for the most experienced.\footnotemark \footnotetext{To assess whether large incumbents simply created subsidiaries to evade liability, I sampled 20 new spinal‐implant manufacturers (a high‐innovation segment) submitting 510(k)s after the 2008 down‐regulation. Only one was a subsidiary; the remainder were new firms (founded in 2000 or later; median founding year = 2011, mean = 2011). Over the subsequent 13 years, only two of the 20 were acquired; 80\% were U.S.-based.} This pattern aligns with Proposition~\ref{prop:2} (Appendix~\ref{sec:AppendixConcept}), described in Section \ref{sec:Concept}, which predicts larger gains for firms facing higher ex-ante regulatory delays.\footnotemark \footnotetext{This pattern is robust to measuring proficiency as cumulative product count rather than cumulative approval time (Figure~\ref{fig:FirmExperienceHeteroSimulation}, Panel~(b)).}


\textit{Firm Size.} I next examine differential innovation responses by asset terciles, using patenting rates as the outcome. Patents are already linked to firm‐level assets via \cite{kogan:2017} and provide a common metric across both reclassification pathways. Figure~\ref{fig:FirmSizeHetero}, Panels (a)--(b) show that both transitions (Class~III to II and II to I) result in the smallest firms (lowest asset tercile) displaying the largest post‐event increases in patenting, consistent with Proposition~\ref{prop:1} in Appendix~\ref{sec:AppendixConcept}: when regulatory and financing frictions fall and expected profits rise, marginal returns to R\&D are highest for capital‐constrained firms.\footnotemark \footnotetext{While unobserved factors correlated with size or experience could, in principle, confound these patterns, two diagnostics mitigate concern. First, in my data, firm size is empirically uncorrelated with FDA experience, allowing the two channels to be distinguished. Second, industry documentation and manufacturer statements repeatedly cite regulatory know‐how and financing constraints as pivotal inputs into R\&D decisions \citep{Makower:2010,emergo_2019}, reinforcing the interpretation that relaxed requirements disproportionately benefit smaller and less regulation‐proficient firms.}
 


\subsection{Changes in Device Safety}\label{sec:adverseeventresults}

The preceding sections show that lighter oversight spurs innovation, entry, and lower prices---but premarket review exists to generate information about device safety and performance before products reach patients. Whether postmarket mechanisms can substitute for premarket review is the central empirical question. I examine this by re-estimating Equation \ref{eq:OLS} for two outcomes: (i) serious adverse‐event rates (events per year) and (ii) the probability that any patent in a device-year explicitly emphasizes safety. Table \ref{tab:AdverseEvent}, Panel A indicates that down-regulation (Class~III to II) is not associated with consistent, statistically significant changes in adverse-event rates across control groups.\footnotemark \footnotetext{Table~\ref{tab:Combined_DropNoCounts} restricts to device types with at least one positive count of the relevant outcome across the full panel.} The preferred specification does, however, suggest an economically meaningful rise in hospitalizations. The likelihood of any safety-emphasis patent increases under the baseline identification but is not robust to alternative safety-innovation measures.\footnotemark \footnotetext{See Appendix~\ref{appendix:safety} and Table~\ref{tab:ProbabilitySafetyMentionEffects_Robustness} for details on alternative methods for identifying safety innovations.}


By contrast, Table \ref{tab:AdverseEvent}, Panel B shows that deregulation (Class~II to I) is associated with marked safety improvements: hospitalization rates, life-threatening incidents, and fatalities fall significantly under most control comparisons, and safety-focused innovation rises. Magnitudes imply an approximately 77\% reduction in annual hospitalizations and a suggestive 40--70\% decline in deaths (pre-mean: 0.2 deaths/yr), relative to counterfactual levels. The probability that any patent emphasizes safety increases by about 80\%, and this effect is robust to alternative measurement approaches (Appendix~\ref{appendix:safety}; Table~\ref{tab:ProbabilitySafetyMentionEffects_Robustness}), mirroring the pattern in FDA adverse‐event data.\footnotemark \footnotetext{Inventors also show a heightened focus on quality in patent texts post-deregulation, implying that removing regulatory controls on efficacy and quality does not diminish quality.}


Figure \ref{fig:ESProductSafety} plots the event-time coefficients $\beta_k$ from Equation \ref{eq:EventStudy}.\footnotemark \footnotetext{Event studies disaggregated by serious-event subtypes show no divergent pre-trends between treated and control device types (not shown).} Panel (a) suggests a possible increase in serious events after down-regulation, consistent with the static point estimate on hospitalizations (though the pooled DiD effect is statistically insignificant). Panel (b) shows a clear post-deregulation decline across serious events. Consistent with increased manufacturer attention to safety, the probability of safety‐emphasis in patents rises in tandem (Figure~\ref{fig:Safety_IItoI}).

The adverse-event rates in Figure~\ref{fig:ESProductSafety} are raw annual counts, unadjusted for utilization. As discussed in Section~\ref{sec:marketcomposition}, utilization estimates are not stable across specifications; if utilization did increase, these raw counts are a conservative measure of the per-procedure safety improvement. To assess this directly, I normalize annual adverse-event counts by the volume of medical procedures linked to each device type, using MarketScan claims available for 1996--2013. This normalization is feasible for 15 deregulated device types that can be linked to outpatient procedure codes; too few Class~III devices are linked to allow a comparable analysis. Figure~\ref{fig:ESProductSafety}, Panels~(c)--(d), compare the normalizable subset before and after normalization: both series exhibit similar dynamics, and the normalized estimates are somewhat stronger and more statistically precise, confirming that the safety improvements reflect genuine reductions in per-use risk rather than compositional shifts in utilization.



What explains safety gains after deregulation? A leading mechanism is legal liability. Removing ex-ante regulation shifts the locus of enforcement to ex-post litigation, potentially improving safety; larger, asset‐rich firms face higher collectible damages and thus stronger incentives to invest in safety (Propositions~\ref{prop:3}--\ref{prop:4}, Appendix~\ref{sec:AppendixConcept}). Two tests support this mechanism. First, an analysis of heterogeneity by firm assets shows that larger firms exhibit a $>$150\% increase in the probability of at least one safety innovation (Figure~\ref{fig:SafetybyAsset}, Panel~(a)), whereas smaller firms respond less; correspondingly, reductions in serious adverse events are attenuated among small firms (Panel~(b)).\footnotemark \footnotetext{Figure \ref{fig:DeviceTortCasesOverTime}, Panel (b) illustrates that the share of medical device tort claims culminating in bankruptcy among small firms sharply increases after deregulation.} Second, a 2015 deregulation in which devices were \emph{not} liability‐exempt prior to the policy change yields no post‐deregulation safety improvement (Appendix~\ref{sec:AppenixSafetyValid}), pointing to liability exposure---rather than competition or innovation alone---as the key driver of safety gains.






Several caveats are in order. First, Class III device types are reclassified when prospective regulation is deemed sufficient to mitigate harm; therefore, the average safety effect for down‐regulation should be interpreted as a local average treatment effect for device types meeting that criterion. Second, while regulatory changes could, in principle, influence reporting behavior, stringent enforcement of serious-event reporting and the null findings after the 2015 policy change argue against reporting bias as the main explanation. Additional discussion and robustness appear in Appendix Section \ref{sec:Adverse_Event_Concerns}.

% For Class II to I events, however, I assess the FDA's decisions using the described decision rule in Appendix \ref{sec:FDADecisionRule} and estimate equation \ref{eq:OLS} separately for each treated device type matched with a control based on the rule's assigned score. Figure \ref{fig:IItoIAdverseTopDecile} demonstrates more pronounced safety improvements among marginal device types relative to less dangerous treated device types. Thus, my results could extend to most current Class II device types, with 95\% exhibiting fewer adverse events than the most marginal deregulated type. 


% Second, regulatory decisions could potentially influence the reporting of adverse events. However, several factors mitigate this concern. The reporting of serious adverse events is stringently enforced through criminal and financial penalties \citep{bragg_florence_chan_2018,emergo_2022}, ensuring data reliability. Moreover, the lack of impact observed following the 2015 deregulation event suggests that deregulation itself does not skew reporting. Additionally, Appendix \ref{sec:AppenixSafetyValid} includes an analysis of device recalls, which are difficult to conceal, where a post-deregulation decrease in recalls supports the adverse event and inventor emphasis findings. For a comprehensive discussion on these issues and additional details that support my conclusions, see Appendix \ref{sec:Adverse_Event_Concerns}.

% Second, the FDA does not normalize adverse event rates by device utilization due to data limitations. Growth in utilization would increase the likelihood of adverse events. Thus, fluctuations in adverse event rates reflect changes in product safety \textit{and} utilization. Hence, using adverse event rates as a signal of product safety provides a conservative estimate of the net benefit of deregulation as deregulation increases utilization. Figure \ref{fig:Utilization_IIItoII} shows that, although no pre-trends are present, utilization rates of treated medical device types significantly increase three years after Class III to II deregulations, plausibly due to increased supply. Although I do not have similar utilization data for Class II to I events, treated device types also exhibit increased supply after deregulation. All else equal, if the demand curve is not perfectly inelastic, an outward shift in the supply curve would increase utilization. 

% Second, the FDA explitly cautions against using adverse event reports as a measure of safety due to potential under-reporting, inaccuracies, lack of verification, and lack of utilization data. For example, the number of adverse event reports could increase simply because more patients are using a device, reflecting changes in utilization rather than product safety. Hence, using adverse event rates as a signal of product safety provides a conservative estimate of the net benefit of deregulation as deregulation increases utilization. Figure \ref{fig:Utilization_IIItoII} shows that, although no pre-trends are present, utilization rates of treated medical device types significantly increase three years after Class III to II deregulations, plausibly due to increased supply. Although I do not have similar utilization data for Class II to I events, treated device types also exhibit increased supply after deregulation. All else equal, if the demand curve is not perfectly inelastic, an outward shift in the supply curve would increase utilization. 


% Lastly, media and regulatory decisions may influence adverse event reports. Manufacturers, for example, could be less likely to report adverse events if they are subject to less regulatory scrutiny or if reports are more likely to make news after deregulation. However, I focus on mandatory reports of deaths or severe injuries from hospitals and device manufacturers, which are less sensitive to these factors than voluntary reports of less severe injuries \citep{FDAMDR2020}. The FDA enforces the reporting of serious events using financial penalties and criminal resolution \citep{bragg_florence_chan_2018,emergo_2022}.\footnotemark \footnotetext{Both user facilities (i.e., hospitals) and manufacturers are required to report serious adverse events to the FDA. Thus, if either entity fails to report an event, but the FDA is notified by the other (or other sources like end users), then it is implicated in incompliance. Additionally, the FDA \textit{increased} its monitoring of deregulated device types to take appropriate remedial action if products had become less safe, which would make it it more difficult for firms marketing affected devices to hide adverse events relative to those marketing unaffected devices \citep{FR60:38903}.} Additionally, after a placebo Class II to I event in which the litigation environment did not change, I did not find a significant changes in adverse event reporting, suggesting that deregulation does not mechanically lead to lower reporting. Lastly, when the FDA announced the largest Class II to I reclassification event in 1995, it created new authorities that enabled closer monitoring of the affected devices to ``take appropriate remedial action, if necessary'' \citep{FR60:38903}, suggesting that adverse event reports would be more challenging to conceal.



% References (for below)

% Baldwin, R., Cave, M., & Lodge, M. (2012). Understanding Regulation: Theory, Strategy, and Practice. Oxford University Press.
% Carpenter, D., & Moss, D. (Eds.). (2014). Preventing Regulatory Capture: Special Interest Influence and How to Limit It. Cambridge University Press.
% Coglianese, C., & Lazer, D. (2003). Management-Based Regulation: Using Private Management to Achieve Public Goals. Law & Society Review, 37(4), 691–730.
% Kessler, D. P. (2004). Regulation of the Medical Device Industry. In Frontiers in Health Policy Research, Vol. 7. NBER.
% Ogus, A. (1994). Regulation: Legal Form and Economic Theory. Clarendon Press Oxford.
% Viscusi, W. K., Harrington, J. E., & Vernon, J. M. (2005). Economics of Regulation and Antitrust. MIT Press.



\section{Implications for Regulatory Design}\label{sec:RegulatoryDesign}

A central policy question is under which conditions lighter-touch oversight can safely replace stringent premarket review. As product classes mature, regulators and firms accumulate information about failure modes, refine guidance, and improve designs. I formalize this idea by linking post-reclassification outcomes to a device’s \emph{market age} and translating the estimates into practical thresholds to guide reclassification decisions. Once a device type’s risk profile is well understood, targeted special or general controls can replace more resource-intensive premarket requirements without compromising safety \citep{baldwin2011understanding,coglianese2003management,viscusi2018economics}.

\textit{Measuring market age.} For each device type $i$, \emph{market age} at reclassification is defined as the reclassification year minus the earliest verifiable introduction. I take the minimum of (i) the first peer-reviewed clinical study, (ii) the earliest patent filing linked to the device type, and (iii) the first FDA submission.\footnote{Introduction years are assembled via semi-automated extraction with human verification; in a hold-out validation sample, consisting of all down-regulated devices, automated and hand-coded dates align closely ($r=0.76$).}



\textit{Device-level effects.} I study heterogeneity in two outcomes: (i) the share of serious events among all reported adverse events (a severity-composition measure that accounts for utilization) and (ii) annual patent counts. For each treated device type, I estimate a device-level difference-in-differences effect relative to matched controls.\footnotemark \footnotetext{I focus on the severity share rather than serious-event levels to reduce mechanical differences from heterogeneous post-event utilization paths---newer device types typically ramp up adoption while older types plateau. A similar maturity gradient obtains when estimating using $\log(\text{serious events}+1)$ counts.}


\textit{Linking effects to market age.} I relate each device-level DiD estimate to market age using OLS, controlling for baseline severity, baseline patenting, device-category fixed effects, and pathway-specific covariates.\footnotemark \footnotetext{For Class~III to II transitions, I include an FDA-initiated indicator; for Class~II to I, I include the Device Priority Model score (Appendix~\ref{sec:FDADecisionRule}).}




\textit{Findings.} More mature device types become safer post-reclassification. Regressing device-level DiD effects on market age yields a negative age slope for both transitions (Appendix Tables~\ref{tab:impacts_vs_exp_3_2} and \ref{tab:impacts_vs_exp_2_1}). The linear fits cross zero at roughly 30 years (III to II) and roughly 60 years (II to I), beyond which the fitted post-reclassification change in the serious-event measure is negative conditional on covariates (Appendix Figure~\ref{fig:Impact_by_experience}). These crossings should be read as empirical benchmarks rather than rigid decision rules, especially because the II-to-I age slope is estimated imprecisely.\footnote{Using $\log(\text{serious events}+1)$ yields a similar maturity gradient. For Class~III to II, the age slope is negative and marginally significant ($\hat\beta_{\text{age}}=-0.0170$, s.e.\ $0.008$, $p=0.060$); for Class~II to I it is negative but imprecise ($\hat\beta_{\text{age}}=-0.0015$, s.e.\ $0.001$, $p=0.218$). Innovation shows the following pattern: the age slope is $\hat\beta_{\text{age}}=-0.0396$ (s.e.\ $0.013$, $p=0.008$) for Class~III to II and is near zero for Class~II to I ($\hat\beta_{\text{age}}=0.0002$, s.e.\ $0.001$, $p=0.876$).} Innovation effects are largest at lower ages, with steeper negative slopes in the patent-impact regressions for Class~III to II, but remain economically meaningful well beyond the safety thresholds. I measure market age in calendar time. Because safety information likely accumulates with use, cumulative market volume (utilization) is a natural alternative metric; the calendar-time thresholds are therefore best read as benchmarks for device types of roughly average utilization, with higher-volume (lower-volume) types plausibly reaching the safety threshold sooner (later).




\textit{Policy implications.} Many current Class~III and II device types already exceed these market-age thresholds (Figure~\ref{fig:DeviceExperienceCurrent}), with over 80\% of Class III devices exceeding 30 years of age and about 29\% of Class II devices exceeding 60 years, among device types with available invention-year data. These patterns indicate a sizable pool of candidates for reclassification that, according to the estimates, could reduce entry costs and spur innovation without increasing serious adverse events. A back-of-the-envelope cost-benefit calculation for deregulation (Class~II to~I), where the evidence on both innovation and safety effects is strongest, implies net benefits of roughly \$28 million per year per reclassified device type (Table~\ref{tab:CostBenefits}).\footnotemark \footnotetext{On the benefit side, I estimate the effect of deregulation directly on the total market value of public-firm patents per device type per year, using the same difference-in-differences design and matched sample as the rest of the paper (\citealp{kogan:2017} patent values, deflated to 2019 dollars). The point estimate is roughly \$118 million; retaining 20\% after netting out business stealing and creative destruction yields roughly \$23.5 million in annual innovation surplus. The 80\% discount is conservative: \citet{bloom2013identifying} estimate that business-stealing effects roughly offset positive technology spillovers, suggesting a smaller discount may be appropriate. On the cost side, deregulation reduces serious adverse events---hospitalizations fall by 1.8 per year and deaths by 0.4---so there are no safety costs to net against the innovation gains; valued at \$22{,}000 per hospitalization \citep{CostsofHospitalStay:2019} and \$10 million per statistical life (EPA), the safety improvements add approximately \$4.3 million. The calculation omits private-firm patent values, the scientific value of innovation, and gains from firm entry, making it conservative. About 29\% of Class~II device types with available invention-year data exceed the 60-year maturity threshold---approximately 650 device types to which this per-unit estimate could apply.}


\textit{Complementary mechanism (liability).} Section~\ref{sec:adverseeventresults} shows that larger, well-capitalized firms improve safety more after reclassification. A natural interpretation, consistent with standard liability models, is that when damages are collectible, tort liability provides additional discipline on safety; for smaller, thinly capitalized firms, this channel is weaker because they are partly judgment-proof. Policymakers could strengthen liability-based enforcement for these smaller firms by requiring liability coverage for firms below an asset threshold so underwriters, who bear loss risk, can price safety and impose conditions (e.g., maintenance and operator-protocol requirements, incident reporting, recall participation). Where such coverage is credible, earlier transitions than the market-age thresholds may be justified; where it is not, higher cutoffs may be appropriate.






% \section{Back-of-the-Envelope Calculation: Costs \& Benefits}\label{sec:CostsandBenefits}
% This section assesses the costs and benefits associated with reclassification, building on the key findings presented in Section \ref{sec:Results}. The first key finding is that reclassification increases patenting rates.  Estimating the social value of this innovation is inherently challenging, yet for this analysis, the value is approximated based on the combined market value of each additional patent, taking into account the effects of creative destruction and business stealing by conservatively assuming these represent 80\% of the patent's value. Additionally, I found that deregulation reduced adverse event rates. The valuation of this decrease incorporates the statistical value of lives saved and a monetary assessment of prevented hospitalizations, as per the estimations in \citet{CostsofHospitalStay:2019}. The assumptions and detailed calculations underpinning these valuations are thoroughly outlined in Table \ref{tab:CostBenefits}.

% Table \ref{tab:CostBenefits} summarizes the costs and benefits of reclassification events. For Class III to II down-regulation, the unmeasured costs would need to be twice as large as the measured costs to outweigh the benefits. In contrast, Class II to I deregulations exhibit no measurable costs, demonstrating instead a combination of decreased adverse events and increased innovation. The benefits of deregulation equate to approximately \$53 million annually per treated device type. Extrapolating these findings, if all 2,500 current Class II devices were deregulated, the potential annual net benefits could reach \$132 billion, equivalent to 77\% of the value of the US medical device market.

% I do not include all the costs and benefits of deregulation in these calculations. For costs, I do not account for the potential loss or attenuation of the informational value of FDA approvals resulting from reclassification (see \citet{Grennan:2020}). However, given that one criterion for reclassification is the ease of verifying and maintaining device efficacy post-deregulation, these costs are likely minimal. Additionally, the value of waiting to down-classify to learn more about a device type's inherent risk is not considered (i.e., the option value of waiting). Nonetheless, given that deregulation leads to lower adverse event rates, delaying deregulation may not offer significant value. Lastly, the potential political costs of misguided deregulation are not quantified.

% There are several unmeasured potential benefits of deregulation. These include reductions in FDA administrative costs, health care cost reductions from deregulation, the economic value of new jobs created following firm entry, the benefits of innovation from private firms, and the broader scientific value of such innovation. 


  
\section{Discussion and Conclusion}\label{sec:Conclusion}


This paper quantifies the effect of FDA medical device regulation on innovation, market structure, and safety. With FDA-regulated products comprising a \$2.8 trillion annual market \citep{FDAFactsheet2020}, understanding these impacts is crucial. By assembling a comprehensive dataset on innovation, market dynamics, firm profiles, and safety records, I provide causal evidence on the consequences of ``down-regulation'' (Class III to II) and ``deregulation'' (Class II to I).

Lighter oversight substantially boosts both the quantity and quality of innovation. Patent filings and unique device submissions increase by approximately 140--350\% under the preferred specification, with pronounced benefits accruing to less regulation-proficient and smaller firms. These policy transitions increase market entry, widen product variety, and---for generic devices---lower procedure prices. Product safety does not systematically change after down-regulation, while deregulation yields sizable declines in serious adverse events alongside a robust rise in safety emphasis in patent texts.

This paper also demonstrates that two forces---legal liability and device market age---systematically influence how reclassification affects safety. First, I exploit heterogeneity by firm assets and find that safety improves more for well-capitalized firms---where damages are collectible---than for smaller, judgment-proof producers, indicating that courts discipline safety when exposure is credible. For small firms, mandating liability coverage can delegate monitoring to insurers. Second, I quantify maturity benchmarks at which lighter-touch oversight is predicted not to worsen safety: about 30 years for Class III to II and roughly 60 years for Class II to I. These are informative benchmarks, not literal policy cutoffs; their main value is to show that the safety consequences of reclassification vary systematically with technological maturity. The fitted relationships suggest that when reclassification occurs beyond these ages, serious adverse events tend to decline and innovation rises, implying economically meaningful net benefits from targeted, maturity-contingent deregulation.


These insights extend beyond medical devices, but the boundaries of extrapolation matter. The estimates here apply to low- and moderate-risk technologies with established safety records. For high-risk or dramatically novel devices---AI-enabled diagnostics, brain-computer interfaces---premarket evidence requirements serve a different function: they generate safety information that does not yet exist and that tort liability alone cannot efficiently produce. The maturity thresholds I estimate (about 30 years for Class~III to~II, roughly 60 years for Class~II to~I) offer one empirical benchmark for when accumulated market experience can substitute for premarket review.

The medical device sector has features that make it especially responsive to deregulation: incremental innovation, low fixed development costs relative to pharmaceuticals, and short product cycles that allow rapid competitive entry. In pharmaceuticals or biologics---where development costs are orders of magnitude larger and clinical trials span years---removing regulatory barriers would likely yield smaller supply responses.\footnotemark \footnotetext{\citet{tabarrok2000assessing} offers evidence regarding the optimality of FDA pharmaceutical regulations.} The liability channel also differs: pharmaceutical manufacturers already face well-developed mass tort litigation that disciplines safety, whereas device manufacturers have historically enjoyed broader preemption protections, making the regulatory-to-liability substitution studied here particularly salient in the device context. Parallels with E.U.\ device oversight, generic drug regulation, and GM food approvals suggest broader value in regulatory designs that (i) tailor requirements to market maturity and (ii) rely on liability-based enforcement where damages are collectible. Future work could test these mechanisms in settings with different cost structures, liability regimes, and innovation dynamics.


% This paper analyzes the effect of FDA regulation on innovation, market structure, and safety. Quantifying these effects is important given that FDA-regulated products represent a \$2.8 trillion market annually \citep{FDAFactsheet2020}. I focus my analysis on the medical device industry. I first present a theoretical framework highlighting the disproportionate burden of regulation on small and inexperienced firms, primarily due to "learning by doing" and high financing costs. The framework suggests that deregulation could enhance the profitability of innovation for these firms and potentially promote product safety through increased liability exposure. Empirically, I explore these propositions---and my broader questions---within the medical device sector, characterized by stringent regulations that preclude litigation. By assembling a comprehensive dataset on innovation, market dynamics, firm profiles, and safety records, the analysis reveals that down-regulation favors smaller and less-experienced firms and broadly accelerates technological progress and firm entry, which, in turn, may lower health care prices. Moreover, deregulation is associated with a significant reduction in adverse events, underscoring the potential of legal liability to motivate safety improvements, especially in larger firms with significant assets vulnerable to liability cases.


% A back-of-the-envelope calculation suggests that the net benefits of deregulation, from Class II to I, approximate \$53 million annually per affected device type. These benefits are higher for marginal, higher-risk device types, suggesting my results may generalize to other Class II devices. These results align with critiques from the National Institute of Medicine and physicians regarding the efficacy of Class II regulations and present deregulation as a viable alternative to simulate quality and innovation. 

% Evaluating the shift from Class III to II presents challenges. While down-regulation significantly boosts new technology---a 470\% increase---this benefit could be tempered by potential risks to the FDA's reputation from increased device-related incidents. For instance, the FDA faces asymmetric costs and benefits; an increase in salient device-related deaths could degrade the regulator's reputation and undermine its more cost-effective efforts elsewhere \citep{carpenter2004political,carpenter2004protection}. By contrast, the technological benefits that come from deregulation are more abstract, leading to a strategy that is ``too conservative'' \citep{isakov2019fda} relative to the social optimum.\footnotemark \footnotetext{This asymmetry is evident in FDA regulatory proceedings outlining the criteria for reclassification as the value of forgone innovation is not considered.} This study aims to clarify these forgone benefits, though further empirical work is needed to fully understand the costs of regulatory mistakes and the value of regulator reputation.

% Although centered on the large and growing medical device market, the implications of this study may extend to other regulated sectors. For instance, FDA regulations for Class III devices are similar to those in the EU, and requirements for these devices resemble those for brand-name drugs in the US and other countries \citep{VanNorman:2016}.\footnotemark \footnotetext{\citet{tabarrok2000assessing} offers evidence regarding the optimality of FDA pharmaceutical regulations.} Additionally, Class II device regulations are similar to those used abroad and resemble those for generic drugs---which are also protected from product design tort claims after FDA approval---and genetically modified (GM) foods \citep{schwartz2020s,schauzu2000concept}. These similarities suggest that medical technology and food regulations may slow innovation and increase market concentration worldwide. Lastly, my analysis highlights the potential issues that arise when regulators use imperfect proxies or heuristics to evaluate product quality, such as the ``substantial equivalence'' heuristic used for Class II devices, generic drugs, tobacco products, and GM foods. These heuristics may be particularly pervasive when product quality is hard to verify or when regulators are under-resourced. In such situations, a robust legal system with impartial judges and high damage caps may better incentivize product safety through litigation. 



 % My results provide several specific, albeit suggestive, insights. First, simulation results suggest that simplifying and standardizing Class III regulations could increase the flow of innovation by 63\%.\footnotemark \footnotetext{If learning by doing is primarily on the part of the FDA, as it seeks to learn which firms conduct reliable device testing, then more uniform standards would provide additional assurance that the evidence provided by new firms is reliable.} Second, resourcing new MedTech firms during the approval process, helping them overcome upfront capital costs, may prove socially beneficial. Some European governments, for example, provide financial assistance during the early phases of medical device approvals \citep{propel2017}. Lastly, deregulation could be a meaningful tool to combating rising health care costs \citep{cutler2005your,smith2009income}. 


% Lastly, deregulating when the benefits of such actions outweigh the costs could improve the dynamism of medical device markets. Firm creation brings job creation. Descriptive evidence suggests that new high-paying MedTech jobs create ripple effects: For every 100 MedTech jobs created, 447 indirect jobs are also created \citep{Makower:2010}. 


% Resourcing the FDA to dedicate staff directly to balancing this decision process could lead to more efficient outcomes (e.g., in contrast to the FDA's Center for Drug Evaluation and Research, the FDA's Center for Devices and Radiological Health does not currently have an economics group).

% Different things I could get across:
% -Need to actually provide a quick summary of the different aspects of the paper (very important)
% -My results are generalizable to different areas: EU and other medical device regulations and pharam. But also any regulation that uses proxies (I think this is important: good to make connections of these findings to other things besides just FDA MDR.
% One way to do this is to also talk about how small and inexperienced firms face difficulties with regulation in any sector (hospitals, producers (environ), etc.))
% -European fixes to learning by doing and capital costs (moderately important)
% -caveats (very important)
% -deregulation improves market dynamism (least important--what is different is speculative, so just emphasize this in the results)
% -assymetries (slightly important); but emphasis on FDA's lack of emphasis on innovation is important because that speaks to the importance of the results


\newpage
\typeout{}
\begin{singlespace}

% Data citations (JPE Policy 1.6): ensure data sources appear in the reference list
\nocite{data_marketscan,data_dimensions,data_crsp,data_compustat,data_nexisuni,data_maude,data_fda510k,data_fdapma,data_fdaclass,data_fdarecall,data_grid,data_kpst,data_kpss,data_fr84_71794,data_fr82_31976}

\bibliography{FDAPaperv1.bib}
\end{singlespace}

\newpage
\begin{center}
\section*{Figures and Tables}
\end{center}

\clearpage
\newpage

\begin{figure}[!ht]
\centering
\caption{\small{\textbf{Contact Lens Case Study}}}

    \subfloat[Patent Filings]{%
        \includegraphics[trim=0cm 0cm 0cm 0.5cm, clip, width=0.49\textwidth]{../replication_package/output/figures/ContactLens_patents.png}
    }\hfill
    \subfloat[Device Submissions]{%
        \includegraphics[trim=0cm 0cm 0cm 0.5cm, clip, width=0.49\textwidth]{../replication_package/output/figures/ContactLens_submissions.png}
    }

    \vspace{1em}

    \subfloat[New-Firm Entry]{%
        \includegraphics[trim=0cm 0cm 0cm 0.5cm, clip, width=0.49\textwidth]{../replication_package/output/figures/ContactLens_firms.png}
    }

\caption*{\footnotesize{Note: This figure plots annual patent filings, device submissions, and new-firm entry for daily-wear soft contact lenses (product code LPL, reclassified from Class~III to~II in 1994) and extended-wear soft contact lenses (product code LPM, which remained in Class~III). The dashed vertical line marks the year of reclassification.}}
\label{fig:ContactLensCaseStudy}
\end{figure}

\clearpage
\newpage

    \begin{figure}
       \centering \caption{\textbf{Background on Medical Device Regulations}}
\vspace{4mm}

   \includegraphics[trim=0cm 0cm 0cm 0cm,scale = .49]{../replication_package/output/figures/Deregulation_May31_23_Paper.png}
   \vspace{4mm}

\caption*{\footnotesize{Note: This figure presents background on FDA medical device regulations and the reclassification policy changes I leverage in my analysis. The FDA assigned device types to one of three classes, each corresponding to a level of perceived risk, which requires a certain amount of time spent approving the device and associated costs. The time and cost values are averages within the given class derived from \citet{Makower:2010}. The first column includes examples of Class III to II and Class II to I reclassifications. For Class III to II, examples include spinal fusion devices and full-field digital mammography machines. For Class II to I, examples include hospital beds and ventilator tubing. *While Class II devices were protected from legal liability during the historical era analyzed, only contemporary Class II medical devices with ``special controls'' requirements are typically protected from product liability \citep{Costello:2016}.}}
        \label{pic:DeviceDeregulation}
\end{figure}  


    
\clearpage
\newpage




% \begin{figure}
%      \centering   \caption{\textbf{Spinal Implant Use Case---III to II reclassification}}
%      \subfloat[Panel A: Number of New Devices Submitted for Approval]{
%   \centering \begin{overpic}[trim=0cm 0cm 0cm 0cm,scale=0.4]{../Figures/Spinal_Firms.png}
%      \put(13,22){\includegraphics[trim=0cm 0cm 0cm 0cm,scale=0.35]{../Figures/SpinalFusion.png}} 
%      \end{overpic}
%  }
%  \\
%  \subfloat[Panel B: Annual Count of Serious Adverse Events]{
%   \includegraphics[trim=0cm 0cm 0cm 0cm,scale=0.4]{../Figures/Spinal_Firms_Adv.png}
% }
% \caption*{\footnotesize{Note: This figure presents an example of a Class III to II reclassification event. In 2007, the FDA down-classified ``intervertebral body fusion devices'' to Class II. As a control case, ``thoracolumbosacral pedicle screw systems'' remained in Class II. Panel A measures the number of unique contact lens devices submitted to the FDA for approval in a given year, and panel B measures the annual count of serious adverse events. The teal line represents the outcome for spinal fusion devices, and the dark blue line represents the same for spinal screw systems. The vertical black line represents the year of reclassification. The imposed picture on panel A shows both device types implanted into the spine in the same spinal fusion procedure. The text on panel B gives an example of a serious adverse event report that resulted in an injury.}}
% \label{fig:ContactLensUseCase}

% \end{figure}  




% \begin{sidewaysfigure}[!ht]
% \centering \caption{\small{\textbf{Innovation Event-Study Class III to II}}}

% \centering \includegraphics[trim=8cm 0cm 0cm 6cm,scale=0.25]{../Figures/3_2_logged_Main_Panel.png}
% \caption*{\footnotesize{Note: This figure presents the estimates of the $\beta_k$ coefficients from equation \ref{eq:EventStudy} for my innovation measures. Only Class III to II reclassification events are considered. Data are analyzed at an annual frequency. Following \citet{de2019two}, $\beta_k$ is the average of all differences of the outcome evolution of medical device types that have been down-classified from $t-1$ to $t$ to groups for which treatment does not change. The top-left subfigure illustrates the evolution of the patenting rate of treated device types relative to matched control groups. The top-right subfigure describes the evolution of unique device submissions by the FDA for treated device types relative to control groups. The bottom-left subfigure describes the evolution of the average citations-per-patent rate. When no patents are filed in a given year, the citations-per-patent rate is set to zero. The bottom-right subfigure presents the evolution of the average patent value in treated device types relative to controls. Patent values are derived from \citet{kogan:2017}, who calculate the change in a firm's stock market valuation upon patent grant announcements to measure patent value.}}
% \label{fig:InnovationES_IIItoII}
% \end{sidewaysfigure}

\global\pdfpageattr\expandafter{/Rotate 90}

\begin{sidewaysfigure}
\centering
\centering \caption{\small{\textbf{Effect of Reclassification on Innovation}}}


 \subfloat[Patenting Rate, Down-Regulation (Class III to II)]{\begin{overpic}[trim=0cm 0cm 0cm 1.5cm,clip,scale=0.35]{../replication_package/output/figures/3_2_nolog_real_pats_oct22_24.png}
     \put(13,60){\tiny Base rate: 6.3 patents/yr}
  \end{overpic}
 }\hspace{6mm}
 \subfloat[Device Submission Rate, Down-Regulation]{\begin{overpic}[trim=0cm 0cm 0cm 1.5cm,clip,scale=0.35]{../replication_package/output/figures/3_2_nolog_dev_count_oct22_24.png}
     \put(13,60){\tiny Base rate: 0.9 dev/yr}
  \end{overpic}
}
 \\
  \noindent
  \subfloat[Patenting Rate, Deregulation (Class II to I)]{\begin{overpic}[trim=0cm 0cm 0cm 1.5cm,clip,scale=0.35]{../replication_package/output/figures/2_1_nolog_real_pats_oct22_24.png}
     \put(13,60){\tiny Base rate: 14.8 patents/yr}
  \end{overpic}
 }\hspace{6mm}
 \subfloat[Device Submission Rate, Deregulation]{
  \includegraphics[trim=0cm 0cm 0cm 1.5cm,clip,scale=0.35]{../replication_package/output/figures/2_1_nolog_dev_count_solo.png}
}

\caption*{\footnotesize{Note: This figure presents the $\beta_k$ coefficient estimates from the event-study equation (refer to Equation \ref{eq:EventStudy}) for innovation outcomes. Device types in control groups are matched based on pre-event outcome averages (levels). The coefficient $\beta_{-1}$ serves as the reference and is therefore omitted. Data are analyzed annually. Every panel traces the trajectory of outcomes in treated groups compared to controls surrounding the reclassification event. Panel (a) depicts the annual patent filing trends before and after down-regulation events (Class III to II). Panel (b) shows the annual evolution of device submissions to the FDA. Panel (c) portrays patenting trends for deregulated device types (Class II to I). For clarity, Panel (d) lacks data on Class I device submissions due to their non-mandatory nature for market entry. Standard errors for down-regulation events follow the methodology of \citet{conley2011inference}. 95\% confidence intervals are provided.}}
\label{fig:InnovationES}
\end{sidewaysfigure}


% took out: 
%  \noindent
%  \subfloat{
%   \includegraphics[scale = .35]{../Figures/3_2_nolog_absnew_firm_count_solo_CTs.png}
%  }
%   \subfloat{
%   \includegraphics[scale = .35]{../Figures/3_2_nolog_firm_expansion_solo_CTs.png}
%  }

% \newpage

% \global\pdfpageattr\expandafter{/Rotate 90}


% \begin{sidewaysfigure}[!ht]
% \centering \caption{\small{\textbf{Innovation Event-Study Class II to I}}}

%  \subfloat{
%   \includegraphics[scale = .35]{../Figures/2_1_nolog_real_pats_solo_CTs.png}
%  }\hspace{6mm}
%  \subfloat{
%   \includegraphics[scale = .35]{../replication_package/output/figures/2_1_nolog_dev_count_solo.png}
% }
%  \\
%  \noindent
%  \subfloat{
%   \includegraphics[scale = .35]{../Figures/2_1_nolog_times_cited_solo_CTs.png}
%  }
%   \subfloat{
%   \includegraphics[scale = .35]{../Figures/2_1_nolog_xi_real_solo_CTs.png}

%  }\caption*{\footnotesize{Note: This figure presents the estimates of the $\beta_k$ coefficients from event-study equation \ref{eq:EventStudy} for my innovation measures. Only Class II to I reclassification events are considered. Controls are device types matched on baseline averages of the outcome. The coefficient $\beta_{-1}$ is omitted and serves as the reference period. Data are analyzed at an annual frequency. The top-left subfigure illustrates the evolution of the patenting rate of treated device types relative to matched control groups. The top-right subfigure is blank as there is no comparable data for Class I approved devices. The bottom-left subfigure describes the evolution of the average citations-per-patent rate. When no patents are filed in a given year, the citations-per-patent rate is set to zero. The bottom-right subfigure presents the evolution of the average patent value in treated device types relative to controls. Patent values are derived from \citet{kogan:2017}, who calculate the change in a firm's stock market valuation upon patent grant announcements to measure patent value.}}
% \label{fig:Innovation_IItoI}
% \end{sidewaysfigure}


\clearpage
\newpage

\begin{sidewaysfigure}[!ht]
\centering \caption{\small{\textbf{Effect of Reclassification on Market Structure
}}}

 \subfloat[New Firms Patenting, Down-Regulation (Class III to II)]{
 \begin{overpic}[trim=0cm 0cm 0cm 1.5cm,clip,scale=0.35]{../replication_package/output/figures/3_2_nolog_unique_firms_solo_CTs.png}
     \put(13,60){\tiny Base rate: 3.0 firms/yr}
  \end{overpic}
 }\hspace{6mm}
 \subfloat[New Firms Patenting, Deregulation (Class II to I)]{ \begin{overpic}[trim=0cm 0cm 0cm 1.5cm,clip,scale=0.35]{../replication_package/output/figures/2_1_nolog_unique_firms_oct22_24.png}
     \put(13,60){\tiny Base rate: 6.8 firms/yr}
  \end{overpic}
}
 \\
 \noindent
 \subfloat[Log Price, Down-Regulation]{
  \includegraphics[trim=0cm 0cm 0cm 1.5cm,clip,scale=0.365]{../replication_package/output/figures/3_2_nolog_Price_solo_CTs_sept20_25.png}
 }
  \subfloat[Log Price, Deregulation]{
  \includegraphics[trim=0cm 0cm 0cm 1.5cm,clip,scale=0.365]{../replication_package/output/figures/2_1_nolog_Price_solo_CTs_sept20_25.png}

 }\caption*{\footnotesize{Note: This figure presents the $\beta_k$ coefficient estimates from the event-study equation (refer to Equation \ref{eq:EventStudy}) for outcomes related to market structure. Device types in control groups are matched based on pre-event outcome averages. The coefficient $\beta_{-1}$ serves as the reference and is therefore omitted. Data are analyzed annually. Each panel contrasts the treated groups' outcomes with control groups around the time of the reclassification event. Specifically, Panel (a) depicts the annual count of new firms patenting around the time of down-regulation events (Class III to II). Panel (b) shows the count of new firms patenting around deregulation events (Class II to I). Panel (c) plots estimates of log outpatient payments for device-linked medical procedures associated with down-regulated device types, using MarketScan claims. Panel (d) provides a similar analysis for procedures associated with deregulated device types for which I only have two years of pre-event claims data. Standard errors for down-regulation events follow the methodology of \citet{conley2011inference}. 95\% confidence intervals are provided.}}
\label{fig:MarketStructure}
\end{sidewaysfigure}


% \begin{sidewaysfigure}[!ht]
% \centering 
% \caption{\small{\textbf{Effects of reclassification on Market Structure}}}

% % First row
% \subfloat[New Firms Patenting, Down-Regulation]{
%  \begin{overpic}[trim=0cm 0cm 0cm 1.5cm,clip,scale=0.28]{../replication_package/output/figures/3_2_nolog_unique_firms_solo_CTs.png}
%      \put(13,60){\tiny Base rate: 3.0 firms/yr}
%   \end{overpic}
% }
% \subfloat[New Firms Submitting, Down-Regulation]{
%  \begin{overpic}[trim=0cm 0cm 0cm 1.5cm,clip,scale=0.28]{../Figures/3_2_nolog_absnew_firm_count_solo_CTs.png}
%      \put(13,60){\tiny Base rate: 1 firm/10 yrs}
%   \end{overpic}
% }
% \subfloat[New Firms Patenting, Deregulation]{ 
%  \begin{overpic}[trim=0cm 0cm 0cm 1.5cm,clip,scale=0.28]{../Figures/2_1_nolog_unique_firms_solo_CTs.png}
%      \put(13,60){\tiny Base rate: 7 firms/yr}
%   \end{overpic}
% }


% % Second row
% \noindent
% \subfloat[Logged Price, Down-Regulation]{
%   \includegraphics[trim=0cm 0cm 0cm 1.5cm,clip,scale=0.35]{../Figures/IIItoIIPrices_latest.png}
% }\hfill
% \subfloat[Logged Price, Deregulation]{
%   \includegraphics[trim=0cm 0cm 0cm 1.5cm,clip,scale=0.35]{../Figures/IItoIPrices_latest.png}
% }

% \caption*{\footnotesize{Note: This figure presents the estimates of the $\beta_k$ coefficients from the event-study equation \ref{eq:EventStudy} for some innovation and market structure outcomes. Only Class II to I reclassification events are considered. Controls are device types matched on baseline averages of the outcome. The coefficient $\beta_{-1}$ is omitted and serves as the reference period. Data are analyzed at an annual frequency. The top-left subfigure illustrates the evolution of the patenting rate of treated device types relative to matched control groups. The top-middle subfigure is blank, as there is no comparable data for Class I approved devices. The top-right subfigure illustrates the evolution of the rate of new firm entry (measured by new firms patenting) relative to matched control groups. The bottom-left subfigure illustrates the evolution of the rate of incumbent firm entry (firms that have received a granted patent), entering treated device types relative to matched controls. I do not include FDA-approved device measures of new and incumbent entry as I do not have reliable data on new Class I devices from FDA sources. 95\% confidence intervals are provided.}}
% \label{fig:Innovation_IItoI}
% \end{sidewaysfigure}




\clearpage
\newpage

\global\pdfpageattr\expandafter{/Rotate 90}


\begin{sidewaysfigure}[!ht]
\centering \caption{\small{\textbf{Effect of Reclassification on Product Safety}}}

 \subfloat[Serious Adverse Events, Down-Regulation (Class III to II)]{
  \begin{overpic}[trim=0cm 0cm 0cm 1.5cm,clip,scale=0.35]{../replication_package/output/figures/3_2_nolog_v_serious_oct22_24.png}
     \put(55,60){\tiny Base rate: 0.6 events/yr}
  \end{overpic}
 }\hspace{6mm}
 \subfloat[Serious Adverse Events, Deregulation (Class II to I)]{
  \begin{overpic}[trim=0cm 0cm 0cm 1.5cm,clip,scale=0.35]{../replication_package/output/figures/2_1_nolog_v_serious_oct22_24.png}
     \put(40,60){\tiny Base rate: 0.5 events/yr}
  \end{overpic}
 }
 \\
 \noindent
 \subfloat[Normalizable Subset, Raw Counts (Class II to I)]{
  \includegraphics[trim=0cm 0cm 0cm 1.5cm,clip,scale=0.35]{../replication_package/output/figures/2_1_nolog_v_serious_normalizable_oct22_24.png}
 }\hspace{6mm}
 \subfloat[Normalizable Subset, Normalized by Procedure Volume (Class II to I)]{
  \includegraphics[trim=0cm 0cm 0cm 1.5cm,clip,scale=0.35]{../replication_package/output/figures/2_1_nolog_v_serious_restrict_oct22_24.png}
 }
\caption*{\footnotesize{Note: This figure plots event-study estimates $\beta_k$ from Equation~\ref{eq:EventStudy} for serious adverse events (life-threatening incidents, deaths, hospitalizations, or disabilities) using annual data. Controls are device types matched on pre-event outcome averages; $\beta_{-1}$ is the omitted reference period. Panel~(a): down-regulation (Class~III to II) vs.\ matched controls. Panel~(b): deregulation (Class~II to I) vs.\ matched controls. Panels~(c)--(d) restrict the sample to the 15 deregulated device types that can be linked to outpatient procedure codes in MarketScan claims (1996--2013). Panel~(c) reports raw event counts for this subset; Panel~(d) normalizes by annual procedure volume. The pre-deregulation mean for the normalizable subset is 0.002 serious events per procedure. Too few Class~III devices are linked to procedure codes to allow a comparable normalization. Conley--Taber 95\% confidence intervals are reported.}}
\label{fig:ESProductSafety}
\end{sidewaysfigure}





\clearpage
\newpage

\pdfpageattr\expandafter{/Rotate 0}


\begin{figure}[!ht]
\centering \caption{\small{\textbf{Change in Product Safety by Firm Asset Terciles (Class II to I)}}}
\subfloat[Change in Safety Effort by Firm Assets]{
  \includegraphics[trim=0cm 0cm 0cm 1.7cm,clip,scale=0.42]{../replication_package/output/figures/2_1_Safety_Top_Decile_Present_oct22_24.png}
 }

 \subfloat[Change in Serious Events by Firm Assets]{
  \includegraphics[trim=0cm 0cm 0cm 1.7cm,clip,scale=0.42]{../replication_package/output/figures/2_1_Deaths_Top_Decile_oct22_24.png}
}
    \caption*{\footnotesize{Note: This figure presents separate difference-in-differences (DiD) estimates from Equation \ref{eq:OLS}, detailing the change in the probability of firms having at least one instance of specified outcomes within treated device types, categorized by firm asset terciles. Given the low frequency of safety mentions and serious events at the firm level, all non-zero instances are coded as one, thus transforming the analysis into a Linear Probability Model (LPM). Baseline outcome frequencies are similar across asset terciles. Panel (a) examines changes in the probability of safety-related innovations, and Panel (b) addresses changes in the probability of serious adverse events (such as death, hospitalization, or life-threatening incidents). The terciles represent asset sizes of publicly traded firms. The x-axis shows the asset terciles (first, second, or third), while the y-axis reflects the percentage change in likelihood. Bars signifying 95\% confidence intervals are also included.}}
    \label{fig:SafetybyAsset}%


\end{figure} 


    
% \begin{sidewaysfigure}
%       \centering
% \centering
% \centering \caption{\small{\textbf{Innovation Event-Study Class II to I}}}


%  \subfloat{
%   \includegraphics[width=0.5\linewidth]{../Figures/2_1_logged_Main_Panel_real_pats.png}
%   \label{subfig:fig1}
%  }
%  \subfloat{
%   \includegraphics[width=0.5\linewidth]{../Figures/2_1_logged_Main_Panel_times_cited.png}
%   \label{subfig:fig2}
% }
% \vspace{-6mm}
 
%  \subfloat{
%   \includegraphics[width=0.5\linewidth]{../Figures/2_1_logged_Main_Panel_xi_real.png}
%   \label{subfig:fig3}
%  }

% \caption*{\footnotesize{Note: This figure presents the estimates of the $\beta_k$ coefficients from equation \ref{eq:EventStudy} for my innovation measures. Only Class II to I reclassification events are considered. Data are analyzed at an annual frequency. Following \citet{de2019two}, $\beta_k$ is the average of all differences of the outcome evolution of medical device types that have been down-classified from $t-1$ to $t$ to groups for which treatment does not change. The top subfigure illustrates the evolution of the patenting rate of treated groups relative to matched control groups. The middle subfigure illustrates the evolution of the citations-per-patent rate of treated groups relative to controls groups. The bottom subfigure describes the evolution of the average patent value of treated groups relative to control groups. Average patent values are derived from \citet{kogan:2017} and describe the change in a firm's stock price upon patent grant announcements. I do not include a new device measure as no reliable FDA data exists for new Class I devices.}}
% \label{fig:Innovation_IItoI}
%     \end{sidewaysfigure}   



% \clearpage
% \newpage

% \pdfpageattr\expandafter{/Rotate 0}

% \begin{figure}[!ht]
% \centering \caption{\small{\textbf{Annual Costs and Benefits of reclassifications (Dollars in Millions)}}}
%     \centering
%             \caption*{Panel A: Class III to II reclassifications
% }
% \vspace{-12pt}
%     \subfloat{{\includegraphics[trim=30cm 35cm 30cm 35cm,scale=0.3]{../Figures/IIItoIICostsBenefits.png} }}%
%     \qquad
%       \caption*{Panel B: Class II to I reclassifications
% }
% \vspace{-12pt}

%     \subfloat{{\includegraphics[trim=30cm 35cm 30cm 35cm,scale=0.3]{../Figures/IItoICostsBenefits.png} }}%
%     \caption*{\footnotesize{Note: This figure presents the back-of-the-envelope calculations for the average yearly costs and benefits of reclassifications, in millions of dollars, across down-classified device types. Table \ref{tab:CostBenefits} details these calculations and the underlying assumptions. Dollar amounts are provided in millions. Panel A presents the average yearly costs and benefits from down-classifying Class III devices. Panel B presents the average yearly costs and benefits from down-classifying Class II devices. Since I did not identify any costs associated with Class II reclassifications (as treated medical device types exhibited fewer adverse events), I treat the decline in adverse event rates as a benefit. Both panels A and B do not consider the informational value of FDA assurances regarding device quality and efficacy---although these benefits are likely low as the FDA chooses device types for which efficacy is easily verified. The cost of mortality is EPA's VSL of \$10 million. I assume an average inpatient hospital stay costs \$22,000 \citep{CostsofHospitalStay:2019} while abstracting away from other costs. I assume creative destruction of 4/5 from the value of patents \citep{kogan:2017}. I do not consider private firm patent values. I also do not consider the scientific value of innovation. I assume no value from firm entry (e.g., not considering the value of new jobs). Using calculations from Medicare and UCSDH data, I assume UCSDH performs .08\% of total U.S. procedures to identify the change in national procedure costs from deregulation.}}
%     \label{fig:CostsandBenefits}%
% \end{figure}




\clearpage
\newpage

\input{../replication_package/output/tables/Table_1_summary_stats.tex}






% \clearpage
% \newpage

% \begin{table}[h!]
% \begin{center}
% \caption{\textbf{Summary Statistics---Class III Devices}}
% \begin{tabular}{lcccc}
% \hline \hline 
% \addlinespace
%  & N & Mean  & SD & Range \\
%  \hline

%  \addlinespace

% % &&&&\\


% \multicolumn{4}{l}{\textit{FDA Admin. Data---Device Submissions (PMA and 510(k) Databases)}} \\
%  \hspace{6mm} Total & 168,880 & - & - & - \\
%   \hspace{6mm} per Device Type & 4,710 (Types) & 35.5	&  110.8 & [1, 2,457]\\
%   \hspace{6mm} Total Submitting Firms& 20,343 & 
% - & -& -	\\
% \hspace{6mm}  Firms per Device Type & 4,710 (Types) & 15.7	&  39.5& [1, 1,048]\\
% \hspace{6mm} Firm Regulatory Proficiency &  4,660 (Types) & 19.5yrs	&  65.4yrs & [0, 686.2yrs]*\\
  
%    \addlinespace
% \addlinespace

% \multicolumn{4}{l}{\textit{FDA Admin. Data---Adverse Event Reports (MAUDE)}}  \\
%  \hspace{6mm} Total & 9,238,733 & - & -  & -\\
 
%   \hspace{6mm} per Device Type & 4,111 (Types) &  2,353.3	&  18,939.9 & [1, 0.6M]	\\
%   \hspace{6mm} Serious Events per Dev. type & 2,400 (Types) & 571.7	&  5186.8	& [1, 0.15M]\\

%         \hspace{6mm} Assets of Offending Firm & 7,139,727 &  \$3.76B &       \$5.77B& [\$0, \$0.79T]	\\
% \addlinespace
%  \addlinespace

% \multicolumn{4}{l}{\textit{USPTO Device Patents}} \\
%  \hspace{6mm} Total & 1,248,292 & - & -  & -\\
%   \hspace{6mm} per Device Type & 2,113 (Types) & 590.8	&  2077.4	& [1, 	23,056]\\
%   \hspace{6mm} Citations & 1,248,292 & 
% 14.6 & 88.8	& [1, 5,817]\\
%     \hspace{6mm} Market Valuation  & 377,465 & \$13.1M & \$30.7M & [\$45, \$1.9B]	\\
%         \hspace{6mm} Applicant Assets  & 377,465 &  \$26.7B &       \$54.8B & [\$0.07M, \$1.1T]	\\
%  \addlinespace
%  \addlinespace

% \multicolumn{4}{l}{\textit{UCSD Healthcare Claims Extract}} \\
%  \hspace{6mm} Total & 495,519  & - & - & - \\
%      \hspace{6mm} per Procedure Code & 528 (Codes) & 880.4		& 2397.5& [1, 18,915]	\\
%   \hspace{6mm} Unique Patients & 55,621 & -	&  -& -	\\

  
  
%   % \hspace{6mm} Amount Billed & 197,407 & \$710.4 & \$2455.2& [\$0, \$0.62M]	\\
%  \hspace{6mm} Price & 453,079 &  \$135.7& \$389.0& [\$0, \$0.01M]	\\
%      % \hspace{6mm} Amt. Billed per Proc. Code & 462 (Codes) & \$1,650.4	 & \$1,779.8 & [\$8.06, \$0.02M]\\
%     \hspace{6mm} Price per Proc. Code & 528 (Codes) & \$354.8	 & \$576.1 & [\$0, \$5,401]\\




% % &&&&\\

% % &&&&\\




% % &&&&\\
% \addlinespace

% \hline 
% \end{tabular}
% \caption*{ \begin{footnotesize}Note: See \citet{kogan:2017} for more information on the patent market valuation data, which was merged into my patent dataset. The CRSP/Compustat database was used to derive the total assets of the firms applying for patent protection and is a proxy for firm size. Market values and applicant assets are only available for patents filed by publicly traded firms, representing roughly 25\% of the total sample of patents. Missing observations account for the discrepancies between (i) the number of total FDA device types (5,542) and the number of device types represented in device submissions, adverse event reports, and patents (many device types have no associated patents), (ii) the total number of patents and the number of patents with market valuations and applicant assets, and (iii) the total number of claims and claims containing amounts paid. *``Regulatory proficiency'' indicates the total number of days a firm has experienced approval delays across all its submitted devices.
% \end{footnotesize} }
% \label{tab:SummaryStatistics}
% \end{center}
% \end{table}



% \begin{table}[h!]
% \begin{center}
% \caption{\textbf{Summary Statistics}}
% \begin{tabular}{lccc}
% \hline \hline 
% \addlinespace
%  & N & Mean & SD  \\
%  \hline

%  \addlinespace

% % &&&&\\


% \textit{FDA Admin. Data---Device Submissions} & & &  \\
%  \hspace{6mm} Total & 168,880 & - & -  \\
%   \hspace{6mm} per Device Type & 4,710 (Types) & 32.62	&  106.72\\
%    \addlinespace

%   \addlinespace

% \textit{FDA Admin. Data---Submitting Firms} & & &  \\
%   \hspace{6mm} Total & 20,343 & 
% - & -	\\
% \hspace{6mm}  per Device Type & 4,710 (Types) & 15.66	&  39.48\\
%    \addlinespace
% \addlinespace

% \textit{FDA Admin. Data---Adverse Event Reports} & & &  \\
%  \hspace{6mm} Total & 9,238,733 & - & -  \\
 
%   \hspace{6mm} per Device Type & 4,111 (Types) &  2,353.27	&  18,939.92	\\
%   \hspace{6mm} Serious Events per Device type & 4,111 (Types) & 333.79	&  3,972.71	\\

%         \hspace{6mm} Assets of Offending Firm & 7,139,727 &  \$3.76B &       \$5.77B	\\
% \addlinespace
%  \addlinespace

% \textit{USPTO Device Patents} & & &  \\
%  \hspace{6mm} Total & 1,248,292 & - & -  \\
%   \hspace{6mm} per Device Type & 2,113 (Types) & 378.17	&  1377.45	\\
%   \hspace{6mm} Citations & 1,248,292 & 
% 14.59 & 88.84	\\
%     \hspace{6mm} Market Valuation (Kogen et al. 2017) & 377,465 & \$13.09M & \$30.72M	\\
%         \hspace{6mm} Applicant Assets (CRSP/Compustat) & 377,465 &  \$26.68B &       \$54.82B	\\
%  \addlinespace
%  \addlinespace

% \textit{UCSD Healthcare Claims Extract} & & &  \\
%  \hspace{6mm} Total & 495,519  & - & -  \\
%   \hspace{6mm} Unique Patients & 55,621 & -	&  -	\\
%     \hspace{6mm} Claims per Procedure Code & 528 (Codes) & 880.42		& 2397.53	\\
  
  
%   \hspace{6mm} Amount Billed & 197,407 & \$710.38 & \$2455.19	\\
%  \hspace{6mm} Amount Paid & 453,079 &  \$135.68& \$388.99	\\
%     \hspace{6mm} Amount Paid per Proc. Code & 528 (Codes) & \$354.83	 & \$576.09 \\




% % &&&&\\

% % &&&&\\




% % &&&&\\
% \addlinespace

% \hline 
% \end{tabular}
% \caption*{ \begin{footnotesize}Note: For information on the collection of all FDA policymaking across the last 40 years, see table \ref{table:ReclassificationEvents}. See \citet{kogan:2017} for more information on the patent market valuation data, which was merged into my patent dataset. The CRSP/Compustat database was used to derive the total assets of the firms applying for patent protection and is a proxy for firm size. Market values and applicant assets are only available for patents filed by publicly traded firms, representing roughly 25\% of the total sample of patents. Missing values account for the discrepancies between (i) the number of total FDA device types (5,542) and the number of device types represented in device submissions, adverse event reports, and patents (many device types have no associated patents), (ii) the total number of patents and the number of patents with market valuations and applicant assets, and (iii) the total number of claims and claims containing amounts billed and paid. 
% \end{footnotesize} }
% \label{tab:SummaryStatistics}
% \end{center}
% \end{table}



\clearpage
\newpage

\pdfpageattr\expandafter{/Rotate 0}


\begin{table}[h!]
\begin{center}
\caption{\textbf{Effect of Reclassification on Innovation}}
\resizebox{.9\textwidth}{!}{\begin{tabular}{lcccc}
\hline \hline &   & \multicolumn{3}{c}{DiD Estimates}\\
\cmidrule{3-5}
 & Pre-mean & Matched & Intuitive& Later\\
 Reclassification& \multicolumn{1}{c}{(1)}  & \multicolumn{1}{c}{(2)}&\multicolumn{1}{c}{(3)} &\multicolumn{1}{c}{(4)}\\
\hline 
&&&&\\
\vspace{5mm}
\textbf{A. Class III to II} &&&&\\

\input{../replication_package/output/tables/Formatted/PaperRegTable3_2_Main_Panel_NoLog_oct22_24.tex}
&&&&\\
\hline

&&&&\\
\textbf{B. Class II to I} &&&&\\
&&&&\\
\input{../replication_package/output/tables/Formatted/PaperRegTable2_1_Main_Panel_NoLog_oct22_24.tex}

&&&&\\

\bottomrule 
\end{tabular}}
\caption*{ \begin{footnotesize}Note: The table presents estimates of Equation \ref{eq:OLS}, a DiD OLS model. Outcomes are derived from \textit{Dimensions} patent data, FDA administrative data, and \citet{kogan:2017}. Column (1) reports 5-year pre-treatment averages for treated device types. Columns (2)--(4) report DiD estimates using (i) matched controls, (ii) intuitively similar devices (treating similar diseases), and (iii) ``later-treated'' devices (reclassified after the sample window), respectively. Device submissions are unavailable for Class I devices. For Column (4), Class III to II control device types are treated after 2015; thus, all post-2015 observations are dropped. Confidence intervals use Conley--Taber inference. Symbols +, *, **, and *** denote statistical significance at 0.10, 0.05, 0.01, and 0.001 levels, respectively.
\end{footnotesize} }
\label{tab:Innovation}
\end{center}
\end{table}




\clearpage
\newpage



\begin{table}[h!]
\begin{center}
\caption{\textbf{Effect of Reclassification on Market Structure}}
\resizebox{.9\textwidth}{!}{
   \small % Reducing font size
\begin{tabular}{lccccc}
\hline \hline 
&   & \multicolumn{4}{c}{DiD Estimates}\\
\cmidrule{3-6}
 & Pre-mean & All CPTs & Matched & Intuitive & Later \\
 Reclassification& \multicolumn{1}{c}{(1)}  & \multicolumn{1}{c}{(2)}&\multicolumn{1}{c}{(3)} & \multicolumn{1}{c}{(4)} & \multicolumn{1}{c}{(5)}\\
 \addlinespace

\midrule 
% &&&&\\
\addlinespace
\textbf{A. Class III to II} & & & & & \\
\addlinespace

% &&&&\\
\input{../replication_package/output/tables/Formatted/RegTable_Prices_3_2_sep25.tex}
% &&&&\\
\addlinespace

\cmidrule{2-6}\\

\input{../replication_package/output/tables/Formatted/PaperRegTable3_2_MarketComposition_NoLog_oct22_24.tex}

% &&&&\\
\addlinespace

\hline

\addlinespace


\textbf{B. Class II to I} & & & & & \\
\addlinespace

% &&&&\\
\input{../replication_package/output/tables/Formatted/RegTable_Prices_2_1_sep25.tex}
% &&&&\\
\addlinespace

\cmidrule{2-6}\\

% &&&&\\
\input{../replication_package/output/tables/Formatted/PaperRegTable2_1_MarketComposition_NoLog_oct22_24.tex}

% &&&\\
\addlinespace

\bottomrule 
\end{tabular}}

\caption*{ \begin{footnotesize}Note: The table reports estimates from Equation \ref{eq:OLS}, a DiD OLS model. Column (1) reports 5-year pre-event means for treated device types. Columns (2)--(5) show DiD estimates using controls that match on (2) baseline prices across all procedure codes, (3) device outcomes, (4) intuitively similar diseases, and (5) later-treated devices. Price data are available only for 1996--2013, limiting sample size. Due to data constraints, Columns (4)--(5) omit price results. † Pre-mean prices are unlogged, while estimates are logged. Confidence intervals use Conley--Taber inference. +, *, **, and *** denote statistical significance at 0.10, 0.05, 0.01, and 0.001 levels, respectively.
\end{footnotesize} }
\label{tab:MarketComposition}
\end{center}
\end{table}



\clearpage
\newpage


\begin{table}[h!]
\begin{center}
\caption{\textbf{Effect of Reclassification on Adverse Events}}
\resizebox{.9\textwidth}{!}{\begin{tabular}{lcccc}
\hline \hline &   & \multicolumn{3}{c}{DiD Estimates}\\
\cmidrule{3-5}
 & Pre-mean & Matched & Intuitive & Later\\
 Reclassification& \multicolumn{1}{c}{(1)}  & \multicolumn{1}{c}{(2)}&\multicolumn{1}{c}{(3)} &\multicolumn{1}{c}{(4)}\\
\hline 
&&&&\\
\vspace{5mm}
\textbf{A. Class III to II} &&&&\\
\hspace{-1mm}\input{../replication_package/output/tables/Formatted/PaperRegTableSafety_IIItoII_NoLog.tex}
\input{../replication_package/output/tables/Formatted/PaperRegTable3_2_AdverseEvent_NoLog_oct22_24.tex}
&&&&\\
\hline

&&&&\\
\textbf{B. Class II to I} &&&&\\
&&&&\\
\hspace{-1mm}\input{../replication_package/output/tables/Formatted/PaperRegTableSafety_IItoI_NoLog.tex}
\input{../replication_package/output/tables/Formatted/PaperRegTable2_1_AdverseEvent_NoLog_oct22_24.tex}

&&&&\\

\bottomrule 
\end{tabular}}
\caption*{ \begin{footnotesize}Note: The table presents estimates of Equation \ref{eq:OLS}, a DiD OLS model. Column (1) reports 5-year pre-event means for treated device types. Adverse event outcomes are derived from the FDA MAUDE database. Columns (2)--(4) show DiD estimates using matched, intuitively similar, and ``later-treated'' control groups, respectively. For Column (4), Class III to II controls are treated after 2015; thus, post-2015 observations are dropped. Confidence intervals use Conley--Taber inference. +, *, **, and *** correspond to statistical significance at 0.10, 0.05, 0.01, and 0.001 levels, respectively.
\end{footnotesize} }
\label{tab:AdverseEvent}
\end{center}
\end{table}









\clearpage
\newpage




\begin{center}
\section*{Appendix Material: For Online Publication}
\end{center}

\appendix

%\setcounter{section}{0}
%\setcounter{table}{0}
%\setcounter{figure}{0}
\renewcommand{\figurename}{\bf Appendix Figure}
\renewcommand{\tablename}{\bf Appendix Table}
\numberwithin{table}{section}
\numberwithin{figure}{section}
\numberwithin{equation}{section}

\vspace{1em}
\noindent\textbf{Contents}
\begin{itemize}[leftmargin=*,itemsep=2pt]
  \item[{\ref*{sec:AppendixRegulations}.}] \hyperref[sec:AppendixRegulations]{Additional Details on FDA Regulations}
  \item[{\ref*{sec:AppendixConcept}.}] \hyperref[sec:AppendixConcept]{Detailed Conceptual Framework}
  \item[{\ref*{app:data}.}] \hyperref[app:data]{Data and Variable Construction and Validation}
  \item[{\ref*{sec:SuppFigTab}.}] \hyperref[sec:SuppFigTab]{Supplemental Figures and Tables}
\end{itemize}

\newpage

\section{Additional Details on FDA Regulations}\label{sec:AppendixRegulations}


This appendix provides context on the timing and rationale of FDA reclassification decisions and a concise reference on device classes. First, I summarize agency motivations for selected Class~III to II down-regulations, emphasizing that timing typically reflected process-driven steps (statutory directives, evidence accumulation, or piggybacking on related actions) rather than contemporaneous technological developments. Second, I contextualize industry-petitioned down-regulations. Third, I describe how the FDA selected Class~II to I deregulations. Finally, I give a brief reference on Class~I, II, and III requirements.


\subsection{Background on Class III to II Events}\label{sec:BackgroundIIItoII}

The cases below illustrate that Class~III to II reclassifications were predominantly process-driven and difficult to anticipate.



\subsubsection*{Daily-Wear Soft Contact Lenses}
\textbf{FR citation:} 59~FR~10284 \quad \textbf{Broad area:} Ophthalmic \quad \textbf{Year Announced:} 1994


\textit{Rationale.} The down-regulation was primarily based on accumulated PMA evidence over the prior decade, much of it not publicly available at the time of deliberation. As noted in the October~1993 panel minutes, ``Key data that were needed to support reclassification were contained in PMAs and were not publicly available for use.'' This procedural hurdle led to the FDA ``waiting ten years to make this [reclassification] announcement.'' Notably, only \emph{daily-wear} lenses were reclassified; extended-wear lenses remained in Class~III (which serves as an ``intuitive'' control group). The timing therefore reflected evidentiary maturity, not a discrete technology shock.




\subsubsection*{Vascular Graft Prosthesis}
\textbf{FR citation:} 64~FR~12774 \quad \textbf{Broad area:} Cardiovascular \quad \textbf{Year Announced:} 1999 


\textit{Rationale.} This action formed part of the Food and Drug Administration Modernization Act (FDAMA, 1997) follow-through on legacy (pre-1976) Class~III types with substantial clinical history. FDAMA~\S515(i) directed the agency to resolve the status of pre-amendment Class~III devices lacking final regulations or PMA requirements. The reclassification was therefore deadline- and process-driven, leveraging long-run safety/effectiveness experience rather than recent technological advances.





\subsubsection*{Tacrolimus Test System}
\textbf{FR citation:} 67~FR~58329 \quad \textbf{Broad area:} Clinical chemistry \quad \textbf{Year Announced:} 2002


\textit{Rationale.} The tacrolimus assay reclassification followed, and was harmonized with, a petition-initiated action for the analogous cyclosporine assay. The FDA identified a common special-controls guidance (“Cyclosporine and Tacrolimus Assays”) as sufficient to assure safety and effectiveness for both, citing similarities in mechanism and clinical management. The timing of this reclassification thus reflected alignment with a related assay rather than a tacrolimus-specific breakthrough.




\subsubsection*{Vascular Embolization Device}
\textbf{FR citation:} 69~FR~77899 \quad \textbf{Broad area:} Cardiovascular \quad \textbf{Year Announced:} 2004

\textit{Rationale.} Panel records emphasize that, decades earlier, clinical information was insufficient to set performance standards; only after extended experience did evidence support robust standards and special controls. The reclassification culminated a long evidence-accumulation process rather than a sudden technological shock.





\subsection{Background on Industry-Petitioned Class III to II Events}\label{sec:IndustryPetition}


The main analysis focuses on FDA-initiated actions; petitioned cases are informative about strategic incentives. Sponsors sometimes petition to (i) reduce first-entry burden, (ii) facilitate follow-on equivalence-based entry, or (iii) broaden access once evidence is mature. Examples include Alto Development’s request to reclassify stainless-steel sutures on the basis of existing materials standards \citep{KahanJonathan:1987}; Centocor’s efforts to reclassify tumor-associated antigen tests (e.g., CA125, PSA) following an initial PMA, likely to enable subsequent 510(k) pathways; and a petition by the Orthopaedic Surgical Manufacturers Association (with AAOS support) to reclassify pedicle screws for pediatric deformities (Docket FDA-2009-M-0101), citing accumulated clinical evidence. These illustrate the mix of motives: accelerating initial entry, smoothing iterative innovation, and expanding access when special controls suffice.




\subsection{Background on Class II to I Events}\label{sec:FDADecisionRule}

Two agency actions---in 1996 and 1998---removed 510(k) requirements for selected Class~II devices. The 1996 action was part of a resource-reallocation strategy to focus review on higher-risk submissions \citep{FR58:37378,FR60:38903}. To prioritize candidates, the FDA used an internal \emph{Device Priority Model} (DPM) that combined injury frequencies (deaths, serious, less-serious) with expected use, health benefit, and effectiveness \citep{FR60:38903}. Injury components increased the score; higher benefit/effectiveness or lower use reduced it. Because exact cutoffs and some inputs were not public, I proxy the DPM with a weighted index of average annual adverse-event counts,

\[
\text{Proxy DPM} \;=\; 0.38 \times \text{Deaths} \;+\; 0.30 \times \text{Serious Injuries} \;+\; 0.12 \times \text{Less-Serious Injuries},
\]

which preserves the risk ordering most relevant to the 1996 action. In the data, top-decile device types have roughly four times the proxy score of the 89th-percentile type, and high-death categories also rank high on serious and less-serious injuries.



The 1998 deregulation implemented FDAMA~\S510(m)(1) (1997), which directed the FDA to list Class~II devices no longer needing 510(k) \citep{FR63:3142}. Internal scoring was paired with four criteria: (i) no reported deaths in prior years; (ii) well-understood characteristics for safe/effective performance; (iii) device changes are either readily detectable by users or unlikely to raise risk; and (iv) modifications are unlikely to create a new device type. These criteria did not hinge on short-run technological improvements or outcome trends, and the agency did not publish exact scores.



\subsection{Class I, II, and III Device Regulation Reference}\label{sec:MoreRegDetails}

\textbf{Class~I (low risk).} Subject to general controls.\footnotemark \footnotetext{``Low risk'' refers to devices that do not support/sustain life and do not present unreasonable risk (e.g., tongue depressors).}~About 41\% of product codes are Class~I, and roughly 90\% of those are exempt from premarket submissions beyond facility registration. New devices may be Class~I if intended use and key characteristics match an existing Class~I type; novel intended use or technology can trigger a new device type requiring review.\footnotemark \footnotetext{Manufacturers can also seek a De Novo classification when appropriate; see \citet{fda:CFR}.}


\textbf{Class~II (moderate risk).} Subject to special controls and typically the 510(k) pathway, which requires substantial equivalence to a predicate. About 56\% of product codes are Class~II. While shorter and less intensive than PMA, 510(k) review is not costless in practice; survey evidence puts FDA-related and adjacent costs around \$24 million \citep{Makower:2010}. If substantial equivalence fails, the sponsor must pursue PMA.



\textbf{Class~III (high risk or insufficient information for special controls).} New high-risk types generally require PMA with clinical evidence; some pre-amendment Class~III types have historically proceeded via 510(k).\footnotemark 
\footnotetext{A small share of 510(k)s includes limited clinical data; requirements vary by context \citep{FDAMD}.}~Class~III codes comprise roughly 2\% of product codes but a disproportionate share of spending \citep{Meier09}. End-to-end premarket testing costs from concept to PMA approval have been estimated at about \$94~million, of which roughly \$75~million reflects FDA-related activities, with pivotal studies the largest component \citep{Makower:2010}. PMA supplements cover many post-approval changes with timelines calibrated to the scope of modification \citep{johnson2012fda}.




 \section{Detailed Conceptual Framework}\label{sec:AppendixConcept}

 In this section, I detail the theoretical model introduced in Section \ref{sec:Concept} for R\&D decisions under regulation and litigation.

I model R\&D as a two-stage process: development and commercialization. First, firms invent and patent a new product, improve its safety profile, and raise capital to cover commercialization. Second, firms bring their products to market by attaining regulatory approval/clearance, forming distribution networks, etc. The model builds on that of \citet{budish2015firms}, who formalize the impacts of commercialization lags on innovation. For comparability, I follow their notation closely wherever possible. I introduce into their framework two alternative policy regimes (i.e., regulation and litigation) that differ in commercialization lags, liability risk, and financing costs.

The model's purpose is to illustrate the range of possible consequences of deregulation, to connect these to underlying fundamentals, and, in particular, to relate these effects to firm traits. In turn, the insights from this model will be helpful for interpreting my empirical results. The model is set in the medical device industry, though the implications may apply to other regulated products.


\subsection{Model Preliminaries}\label{sec:Concept_prelim}

Undirected R\&D yields stochastic inventions to a representative, profit-maximizing firm. Upon realizing the new technology, the firm decides if it will allocate capital for directed R\&D to (i) improve the product's safety profile during the development phase and (ii) commercialize the invention. The firm makes this decision in one of two environments: regulation ``R'' or litigation ``L.'' 
The model is characterized by the following parameters:


\vspace{10pt}

\textit{Timing Parameters.}---The year a firm realizes and develops an invention is given by $t_{invent}$, which I normalize to zero. The years it takes to commercialize the product are $t_{comm,f}$. In the medical device industry, FDA approval plays a key role in delaying commercialization \citep{Makower:2010,pietzsch2012medical}.\footnotemark \footnotetext{Approval delays in other areas of health care, like delays in securing medical procedure reimbursement codes, have also been shown to play a key role in innovation  \citep{dranove2022economics}.} Thus, for concreteness, think of $t_{comm,f}$ as the approval delay. Under litigation $L$, there are no approval delays (i.e., $t_{comm,f,L} = 0$). In the regulated environment, approval delays are positive but decrease with regulatory experience \citep{olson1997firm,Makower:2010,chatterji2009spawned}.\footnotemark \footnotetext{Two factors may explain this pattern, both of which are driven by the complexity of the regulatory process. First, inexperienced firms report difficulty benefiting from hired regulatory experts and must instead learn the process independently \citep{ycombinator_2016}.
% Grabowski and Vernon (1976 Structural Effects of Regulation on Innovation in the Ethical Drug Industry GOOD)
% Thomas 1990 (Regulation and Firm Size: FDA Impacts on Innovation), Olson (1997; Firm Characteristics and the Speed of FDA Approval),. 
From the regulator's perspective, having prior experience with a firm reduces the uncertainty about the quality of its products, which may merit shorter review times \citep{olson1997firm,carpenter2004protection}.} Following \citet{arrow1971economic}, I model this relationship by equating the present delay $t_{comm,f}$ to the learning curve $\beta \, T_f^{-\gamma}$, where $T_f$ is prior experience, $\beta$ is the delay with no prior experience (i.e., $T=1$) and $\gamma > 0$. Delay costs are given by $\chi t_{comm,f}$, where $\chi$ is the yearly cost of approval delays.\footnotemark \footnotetext{\citet{Makower:2010} find an average monthly cost of \$1.3 million for Class III approval delays (e.g., clinical trial costs, etc.). I assume $t_{comm,f}$ and several other parameters below are deterministic for simplicity.}



\vspace{10pt}

\textit{Financing Costs.}---Smaller firms must raise external capital to cover the costs of development and commercialization at time $t_{invent}$.\footnotemark \footnotetext{For simplicity, I assume firms finance their project instantaneously. Although fundraising could prolong commercialization delays, removing this assumption does not change my theoretical results.} Fundraising can be difficult: 56\% of small medical device firms claim funding as a central challenge \citep{emergo_2019}.
% \footnotemark \footnotetext{These firms traverse a ``funding chasm,''  a period of capital scarcity resulting from long approval delays and uncertainty \citep{gaglani_2014,propel2017}.} 
% Following \citet{stein2003agency}, I capture these financing frictions by assuming deadweight costs given by $\theta_e C(e)$, where $C(~)$ is an increasing convex function of external funds $e$, and $\theta_e \in [0, 1]$ is the degree of financing friction in the market.\footnotemark \footnotetext{A firm does not incur financing frictions if external capital is nonpositive (i.e., $C(e \leq 0) = 0$).} 
Following \citet{stein2003agency}, I capture these financing frictions by assuming deadweight costs given by $C(e_f)$, where $C(\cdot)$ is an increasing convex function of external funds $e_f$ (similar to the R\&D model of \citealp{stern2017innovation}).
% \footnotemark \footnotetext{A firm does not incur financing frictions if external capital is nonpositive (i.e., $C(e \leq 0) = 0$).} 
External funds $e_f$ are equal to the difference between the non-financing costs and internal capital $K_f$. I omit other costs of commercialization for simplicity.



% \textit{Probability of Commercialization.}---
% \vspace{10pt}




\vspace{10pt}

\textit{Regulated and Deregulated Effective Lives.}---A successfully commercialized product becomes less relevant over time. For expositional ease, I describe the neoclassical risk-adjusted discount factor of the R\&D project as $\delta$, which includes obsolescence and commercialization risk.\footnotemark \footnotetext{A product may also face a probability of successful commercialization $p$, which may be appropriately modeled as a function of safety effort; however, the FDA approves 80\%--90\% of all medical device submissions \citepalias{us2009fda}. Thus, for simplicity, I assume that approval is certain given a firm achieves the mandated safety effort, and I abstract away from other non-approval-related commercialization uncertainty. Including product denial and commercialization risks does not meaningfully change my theoretical insights.}\textsuperscript{,}\footnotemark \footnotetext{Although obsolescence risk is more appropriately modeled as endogenous to R\&D investments, I follow the patent literature and take it as exogenous \citep{budish2015firms}.} Firms enjoy longer or shorter effective product lives depending on the regulatory environment. Under regulation, I define an invention's \textit{Regulated Effective Life (REL)} as the expected years it will be commercialized and non-obsolete in present value terms as discounted by the regulated firm. The effective life of the regulated product begins at time $t_{comm,f}$, yielding an effective life of $REL_f = \sum_{t=t_{comm,f}}^\infty \delta^t = \delta^{t_{comm,f}}/(1 - \delta)$.
% \begin{equation}\label{eq:regprofit}
% REL~ =  ~ \sum_{t_{comm}}^\infty \delta^t ~ =  ~ \frac{\delta^{t_{comm}}}{1 - \delta}.
% \end{equation}
By contrast, in a deregulated environment $N$, I define an invention's \textit{Effective Life (EL)} similar to \textit{REL}, except the lifespan of the product starts at $t_{invent}$, given by $EL = \sum_{t=0}^\infty \delta^t = 1/(1-\delta)$. Notice that $REL_f < EL$ by definition, as regulated profit flows are delayed.


\vspace{10pt}


\textit{Expected Damages and Safety Effort Costs.}---Borrowing from \citet{shavell1986judgment} and \citet{boomhower2019drilling}, if a firm chooses to commercialize its product, it exerts $x_f$ effort to improve product safety, costing $\psi$ per unit, at $t_{invent}$.\footnotemark \footnotetext{For simplicity, I assume firms exert safety effort instantaneously. Alternatively, safety efforts could prolong commercialization delays. Modeling such delays, however, would not change the model implications.} 
% Firms exert effort to prevent expected legal damages incurred after $t_{comm}$ or to meet the mandated minimum $\underline{x}$. 
Under litigation $L$, a commercialized product generates stochastic adverse events that yield $\phi(x_f;\vec{Z})$ legal damages per year, a random variable with expected value $D(x_f;\vec{Z})$ and vector $\vec{Z}$ containing other factors that influence damages in expectation (e.g., firm seizable assets $K_f$, the litigation environment, damage caps). The expected damages function $D(\cdot)$ is a positive decreasing convex function of safety effort $x_f$. The firm exerts effort to maximize the returns to commercialization by equating the marginal cost of effort $\psi + C_x(\psi x_f^* - K_f)$ to the present value of its marginal benefits $-EL  \cdotp  D'(x_f^*;\vec{Z})$ (i.e., marginal abatement of expected damages).
By contrast, under regulation $R$, the firm is exposed to no legal damages, due to federal preemption. Thus, firms exert the mandated level of safety effort $\underline{x}$, as any further effort yields no return.


\vspace{10pt}



\textit{Profits.}---If the product is successfully commercialized and non-obsolete, it generates profits $\pi$ per year for the innovating firm. Although regulation can affect profits by altering market structure, I do not model this relationship, focusing instead on motivating my firm composition and product safety results. Thus, for simplicity, I assume that deregulation increases the aggregate level of R\&D, consistent with my empirical findings, which implies that deregulation does not cut profits enough to outweigh declines in commercialization costs.\footnotemark \footnotetext{Note that this assumption also places an upper bound on the value of legal damages and safety effort costs after deregulation.} 
% \footnotemark \footnotetext{This approach allows me to focus on modeling costs---the central aim of this study.} 
I assume only expert regulators can perceive safety effort (i.e., asymmetric information); hence, safety effort does not affect profits once a product is approved.   




\subsection{Characterization of the Investment Decision}
In the regulated environment $R$, firm $f$ expects to receive profits from commercializing a device for $REL_f$ years. The firm will develop and commercialize its invention if and only if these expected profits exceed the combined delay, safety effort, and financing costs:\footnotemark \footnotetext{Notice the implicit assumption that firms do not consider the future benefits of regulatory experience (i.e., learning by doing) in their investment decisions. This assumption is consistent with a large literature documenting that managers maximize short-term rather than long-term firm value \citep{budish2015firms}.}
% \protect\daggerfootnote{In BRW, the right-hand sand is a constant cost of $c$.}
\begin{equation}\label{eq:firmprofit}
\text{Regulated Firm Invests} \iff  \underbrace{REL_f}_{\substack{\text{Regulated} \\ \text{effective life}}}~  \cdot ~\underbrace{\pi_R}_{\text{Profits}}  ~ \geq  ~  \underbrace{\chi t_{comm,f}}_{\text{Delay costs}}  +  \underbrace{\psi \underline{x}}_{\substack{\text{Mandated safety} \\ \text{effort costs}}}  + \underbrace{C(e_{f,R})}_{\substack{\text{Financing}\\ \text{costs}}}.
\end{equation}
 The amount of external capital $e_{f,R}$ needed to finance the project is given by the difference between the non-financing commercialization costs and the firm's internal capital $K_f$ (i.e., $e_{f,R} = \chi t_{comm,f}    +  \psi \underline{x} - K_f$ if $e_{f,R} \geq 0$, and $0$ otherwise). 


In the litigation environment $L$, firm $f$ will choose to commercialize if and only if the net expected profits (less expected damages) are greater than the combined safety effort and financing costs:\footnotemark \footnotetext{Note that financing frictions do not affect the payment of damages since they can be financed with profits (i.e., in expectation, damages will always be less than profits if a firm chooses to commercialize).}
\begin{equation}\label{eq:deregfirmprofit}
\text{Deregulated Firm Invests} \iff  \underbrace{EL}_{\substack{\text{Effective} \\ \text{life}}}~ \cdot ~[\underbrace{\pi_L}_{\text{Profits}} - \underbrace{D(x_f^*;\vec{Z})}_{\substack{\text{Expected} \\ \text{damages}}}] ~ \geq  ~   \underbrace{\psi x_f^*}_{\substack{\text{Optimal safety} \\ \text{effort costs}}}  + \underbrace{C(e_{f,L})}_{\substack{\text{Financing}\\ \text{costs}}}.
\end{equation}
The amount of external capital $e_{f,L}$ needed to finance the project is given by the difference between safety effort costs $\psi x_f^*$ and the firm's internal capital $K_f$.

Notice the key differences between the investment incentives in environments $R$ and $L$: firms that commercialize in $L$ (i) expect legal damages, (ii) choose and pay for an optimal level of safety effort, (iii) enjoy a longer effective life of their products, and (iv) do not incur delay costs.\footnotemark \footnotetext{Profits and financing costs also differ across these environments; however, the direction of the difference is ambiguous (e.g., if expected damages are large, safety effort costs could increase financing costs).} 

\subsection{Distortions from Regulation}

I focus on model implications related to distortions in firm participation and safety efforts resulting from regulation. 
Throughout, I assume that deregulation increases the level of R\&D activity.
% \footnotemark \footnotetext{Note, though, that deregulation could decrease R\&D through decreased profits (e.g., increased competition) and increased exposure to legal damages.} 
This assumption is supported by my empirical results and allows me to more clearly motivate the less intuitive results I find in my analysis. 

First, I explore how deregulation can improve product safety. If mandated levels of safety effort are low enough, deregulation can improve safety by increasing the net incentives for safety improvements. I state this formally as follows:
\begin{proposition}\label{prop:3} \textit{(Deregulation can increase firm safety efforts) If the marginal cost of regulated effort is less than the ex-post marginal benefit of that effort (i.e., $\psi + C_x(\psi\underline{x} - K_f) < -EL \cdot D'(\underline{x})$), then deregulation will increase firm safety effort.}
\end{proposition}

% Part 1 of proposition \ref{prop:3} describes the tradeoff between innovation and safety effort introduced by preemption. Absent preemption, if the marginal benefit of mandated safety effort is greater than the marginal costs, then preemption leads to less safety effort. Preemption, however, protects firms from damages, increasing the returns to commercialization. If we assume the firm draws a distribution of inventions, it would commercialize a larger share of them. 

% \footnotemark \footnotetext{Note that (ii) implies (i), but (i) does not imply (ii). NOT SURE THIS IS TRUE}  
This proposition implies that ineffective regulations could make products less safe. I show in Section \ref{sec:Results} that Class II regulations may lead to such an outcome. These insights, however, may be specific to the litigious U.S. environment. For example, if a country aggressively caps damages (represented in $\vec{Z}$), firms would face lower expected damages, and safety efforts could drop relative to regulated levels.

Another factor influencing a firm's expected damages is the value of its seizable assets. Following insights on the ``judgment proof problem'' \citep{shavell1986judgment}, when damages exceed the value of a firm's seizable assets, the difference can be discharged through bankruptcy. This option protects small firms from worst-case damages, lowering expected damages and the marginal benefit of exerting safety effort. Thus, if deregulation increases safety efforts, it will do so most for large firms. I state this as follows (and more formally in Appendix \ref{sec:bankruptcy}):
\begin{proposition}\label{prop:4} 
\textit{(Deregulation introduces bankruptcy distortion)}  
\textit{Assume firm A has fewer assets than firm B (i.e., \(K_A < K_B\)) and that firm A’s assets are insufficient to cover its worst-case damages. Firms A and B are otherwise identical. If deregulation increases both firms' safety efforts above the regulated baseline \(\underline{x}\) (i.e., \(x_A^*, x_B^* > \underline{x}\)), then firm B will increase its safety effort more than firm A, that is, $x_B^* - \underline{x} > x_A^* - \underline{x}$.}
\end{proposition}



The next distortion I detail arises from regulatory complexity (i.e., the delays from complex regulatory requirements). Complexity distorts the composition of firms that commercialize as inexperienced firms reap lower returns from commercialization. Deregulation removes these distortions and disproportionately increases the returns to commercialization for inexperienced firms. To formalize this claim, I present the following proposition:\footnotemark \footnotetext{Proofs are presented in Appendix \ref{sec:proofs}.}
\begin{proposition}\label{prop:2} \textit{(Deregulation disproportionately benefits inexperienced firms) If firm A has less regulatory experience than firm B (i.e., $T_A < T_B$; all else equal), then deregulation increases the returns to commercialization most for firm A.}
\end{proposition}

An example helps illustrate the potentially dramatic implications of Proposition \ref{prop:2}. Consider that firm A has no prior experience, and firm B has one previously commercialized project that was delayed for two years. Firm A must wait out a two-year delay. By contrast, firm B waits out a one-year delay, incurring 50\% lower delay costs than firm A and enjoying a longer effective life of its product. Although deregulation removes delay-related costs for both firms, the increase in returns to commercialization is at least twice as large for firm A. 
% Assuming an identical distribution of invention draws for both firms, deregulation leads firm A to commercialize a wider margin of previously abandoned projects. 



Lastly, I discuss distortions that arise from financing frictions and regulation. Small firms incur deadweight costs when raising capital to commercialize their products \citep{gaglani_2014,emergo_2019}. Deregulation can decrease commercialization costs and financing costs, especially for small firms. I state this claim formally as follows:
\begin{proposition}\label{prop:1} \textit{(Deregulation can disproportionately benefit smaller firms) Assume firm A is smaller than firm B and has non-zero financing costs when regulated (i.e., $K_A < K_B$ and $K_A < \chi t_{comm,A} +  \psi \underline{x}$). Firms A and B are otherwise identical. If deregulation does not increase financing costs for firm A (i.e., $\psi x^*_A < \chi t_{comm,A} +  \psi \underline{x}$), then deregulation increases commercialization returns most for firm A.
}
\end{proposition}

However, deregulation could lead to lower returns to commercialization for small firms if financing costs increase after deregulation. For example, if deregulation induces enough additional safety effort costs to outweigh the decrease in approval delay costs, financing costs could increase for smaller firms. By contrast, if the assumptions hold, Proposition \ref{prop:4} will amplify Proposition \ref{prop:1} as small firms face lower expected damages and lower safety effort costs after deregulation and, thus, even lower financing costs.

\subsection{Proofs}\label{sec:proofs}

\subsubsection*{Proof of Proposition \ref{prop:3}}

\begin{proof}
\textbf{Claim.} If \(\psi + C_x(\psi \underline{x} - K) < -EL \cdot D'(\underline{x})\), then the optimal safety effort under deregulation satisfies \(x_f^* > \underline{x}\).

\noindent\textbf{Proof.}
Suppose, for contradiction, that \(x_f^* < \underline{x}\). Since \(x_f^*\) is optimal,
\begin{equation}
    \psi + C_x(\psi x_f^* - K) \;=\; -\, EL \;\cdot\; D'(x_f^*).
\end{equation}
Because \(x_f^* < \underline{x}\) and \(C_x(\cdot)\) is strictly increasing, 
\[
    C_x(\psi x_f^* - K) \;\le\; C_x(\psi \underline{x} - K).
\]
Also, \(D'(\cdot)\) is strictly increasing, so
\[
    D'(x_f^*) \;<\; D'(\underline{x})
    \quad\Longrightarrow\quad
    -\,EL \cdot D'(x_f^*) \;>\; -\,EL \cdot D'(\underline{x}).
\]
Combining these inequalities yields
\begin{equation}
    \psi + C_x(\psi \underline{x} - K)
    \;>\;
    -\,EL \cdot D'(\underline{x}),
\end{equation}
contradicting the assumption that \(\psi + C_x(\psi \underline{x} - K) < -EL \cdot D'(\underline{x})\). Hence, \(x_f^* > \underline{x}\).
\end{proof}

\subsubsection*{Proof of Proposition \ref{prop:4}}

\begin{proof}
Under deregulation, firm B chooses \(x_B^*\) satisfying
\[
\psi + C_x(\psi x_B^* - K_B) = -\,EL \cdot D'(x_B^*),
\]
and firm A chooses \(x_A^*\) satisfying
\[
\psi + C_x(\psi x_A^* - K_A) = -\,EL \cdot D_T'(x_A^*),
\]
where \(D_T'(x) > D'(x)\) for all \(x\). (Note that both \(D_T'\) and \(D'\) are negative, so this inequality implies that, in absolute terms, the marginal benefit of safety effort is lower for firm A.)

Suppose, for the sake of contradiction, that \(x_A^* \ge x_B^*\). Since \(C_x(\cdot)\) is strictly increasing in its argument and because \(K_A < K_B\), it follows that
\[
\psi + C_x(\psi x_A^* - K_A) > \psi + C_x(\psi x_B^* - K_B).
\]
By the first-order conditions, we have
\[
\psi + C_x(\psi x_A^* - K_A) = -\,EL \cdot D_T'(x_A^*) \quad \text{and} \quad \psi + C_x(\psi x_B^* - K_B) = -\,EL \cdot D'(x_B^*).
\]
Since \(D_T'(x_A^*) > D'(x_B^*)\) and \(EL > 0\), we obtain
\[
-\,EL \cdot D_T'(x_A^*) < -\,EL \cdot D'(x_B^*).
\]
Thus, we deduce that
\[
\psi + C_x(\psi x_A^* - K_A) < \psi + C_x(\psi x_B^* - K_B),
\]
which is a contradiction. Therefore, the assumption \(x_A^* \ge x_B^*\) must be false, implying that
\[
x_A^* < x_B^*,
\]
and consequently,
\[
x_B^* - \underline{x} > x_A^* - \underline{x}.
\]
\end{proof}

\subsubsection*{Proof of Proposition \ref{prop:2}}

\begin{proof}
\textbf{Claim.} If firm A has less regulatory experience than firm B (\(T_A < T_B\)), then deregulation increases A’s returns to commercialization more than B’s.

\noindent\textbf{Sketch of Proof.}
Under regulation \(R\), the more stringent \emph{effective} delays imply
\[
    \text{Returns}_{A,R} < \text{Returns}_{B,R}.
\]
Under a purely litigious environment \(L\), we assume no complexity delays, so
\[
    \text{Returns}_{A,L} = \text{Returns}_{B,L}.
\]
The difference-in-differences is
\[
    (\text{Returns}_{A,L} - \text{Returns}_{A,R})
    \;-\;
    (\text{Returns}_{B,L} - \text{Returns}_{B,R})
    \;=\;
    \]
   \[
    \bigl(\text{Returns}_{A,L} - \text{Returns}_{B,L}\bigr)
    \;-\;
    \bigl(\text{Returns}_{A,R} - \text{Returns}_{B,R}\bigr).
\]
Since \(\text{Returns}_{A,L} = \text{Returns}_{B,L}\), this becomes
\[
    -\,\bigl(\text{Returns}_{A,R} - \text{Returns}_{B,R}\bigr)
    \;>\; 0.
\]
Thus, deregulation increases returns for A more than for B.
\end{proof}

\subsubsection*{Proof of Proposition \ref{prop:1}}

\begin{proof}
\textbf{Claim.} Suppose firm A is smaller than firm B (\(K_A < K_B\)) and has non-zero financing costs only under regulation (i.e., \(K_A < \chi t_{comm,A} + \psi \underline{x}\), but \(K_A \ge \psi x^*_A\) under deregulation). Then deregulation increases A’s commercialization returns more than B’s.

\noindent\textbf{Sketch of Proof.}
Under regulation, A must raise external capital (and thus pay financing costs) because
\[
    K_A < \chi t_{comm,A} + \psi \underline{x}.
\]
By contrast, B may not need to raise external capital if \(K_B \ge \chi t_{comm,B} + \psi \underline{x}\). Hence,
\[
    \text{Returns}_{A,R} < \text{Returns}_{B,R}.
\]
After deregulation, assume \(K_A \ge \psi x^*_A\). Now A does not incur financing costs for its higher (optimal) safety effort. Its returns thus rise more than B’s from deregulation:
\[
    (\text{Returns}_{A,L} - \text{Returns}_{A,R}) 
    \;>\;
    (\text{Returns}_{B,L} - \text{Returns}_{B,R}).
\]
Therefore, A’s deregulated returns increase the most. If smaller firms also have lower expected damages due to bankruptcy, that effect further amplifies A’s gains.
\end{proof}


\subsection{Bankruptcy Protection Model Extension}\label{sec:bankruptcy}

% \subsection {Bankruptcy Protection}\label{sec:bankruptcy}

Following insights from the literature on the ``judgment proof problem'' (see \citealp{shavell1986judgment,boomhower2019drilling}), when damages exceed the value of a firm's seizable assets, the difference can be discharged through bankruptcy. To reflect the bankruptcy option, I augment the model above to allow expected damages to differ by firm assets $K_f$. I define the unspent capital available to cover damages as $u_f$. Unspent capital includes the capital not spent on commercialization costs ($K_f - c_f$) and profits from the current period, given by $u_f = \pi + K_f - c_f$. This term incorporates the simplifying assumption that net profits from the last period are distributed as dividends.\footnotemark \footnotetext{I could relax this assumption by letting $u$ be equal to the unspent capital and the sum of all prior net profits up to a given point in time. This would mean that firms would tend to grow larger and eventually be unable to file for bankruptcy. However, the theoretical insights remain the same as initially smaller firms will face fewer expected damages for some time.} The upper bound of legal damages is given by $\bar{\phi}$.

Let $\nu$ represent the total realized damages from product defects, with probability distribution function $f(\nu; x^*_f,\vec{Z})$. In the presence of bankruptcy, the expected damages are given by 
\begin{equation}\label{eq:firmprofitwbankruptcy}
\text{Expected Damages} ~=~ 
\begin{cases}
D(x_f^*;\vec{Z}) ~~~~~~~~~~~~~~~~~~~~~~~~~~~~~~~~~~~~~~~~~~\text{if}~ u_f\geq \bar{\phi},\\
\underbrace{\left[\int_0^{u_f}\nu f(\nu; x_f^*,\vec{Z})\,d\nu + \int_{u_f}^{\bar{\phi}}u_f\, f(\nu; x_f^*,\vec{Z})\,d\nu \right]}_{D_T(x_f^*;\vec{Z})}  ~~ \text{else}.
\end{cases}
\end{equation}
In words, if the firm's capital stock is at least as high as worst-case damages, the expected damages are the same as above, and the investment decision is unchanged. Smaller firms, however, confront a truncated damages distribution, where all possible damage outcomes exceeding the firm's unspent capital $u_f$ are capped at $u_f$. Thus, instead of paying these outsized damages, the firm declares bankruptcy and contributes the value of its total assets to partially cover its damages. Hence, expected damages $D_T(x_f^*;\vec{Z})$ are determined by the probability-weighted sum of damages from 0 to $u_f$, plus the probability-weighted sum of $u_f$ for all damages higher than $u_f$. Assume that the marginal benefit of safety effort for small firms is less than that of large firms at the same levels of safety effort, as there are fewer damages to abate (i.e., $-D_T'(x_f;\vec{Z}) < -D'(x_f;\vec{Z})$ for all $x_f$).

Bankruptcy protection changes the incentives to improve product safety for small firms. Deregulation introduces firms to legal damages; however, bankruptcy protects small firms from worst-case damages, lowering the marginal benefit of exerting safety efforts. Thus, small firms exert less safety effort than large firms. I state this formally as follows:
\begin{proposition}\label{prop:bankruptcydistortion}
\textit{(Deregulation introduces bankruptcy distortion)} 
Assume firm A has fewer assets than: 
\begin{enumerate}
    \item[(i)] firm B (i.e., \(K_A < K_B\)), and 
    \item[(ii)] its worst-case damages \(\bar{\phi}\) (i.e., \(K_A < \bar{\phi}\)). 
\end{enumerate}
Firms A and B are otherwise identical. Suppose deregulation leads to an increase in both firms' safety efforts relative to a regulated baseline \(\underline{x}\), that is, \(x_A^*, x_B^* > \underline{x}\) (as in Proposition~\ref{prop:3}). Then
\[
    x_B^* \;>\; x_A^*,
    \quad\text{and equivalently}\quad
    x_B^* - \underline{x} \;>\; x_A^* - \underline{x}.
\]
In other words, the larger-capital firm B increases its safety effort \emph{more} than the smaller-capital firm A under deregulation.
\end{proposition}






\section{Data and Variable Construction and Validation}\label{app:data}

\subsection{Patent Data Collection}\label{sec:appendixPatentCollection}


This section describes how I link patents to device types, validate the linkage, and assess robustness to alternative patent samples.


\paragraph{Construction.}


For each of over 5{,}000 FDA device types, I construct a keyword string from (i) the regulation-number description and (ii) the device-name description, removing stopwords, punctuation, and duplicates. For example, “Implantable pacemaker pulse generator” and “Leadless Pacemaker” become “implantable pacemaker pulse generator leadless.” I then search the full text of all U.S. patent documents in \emph{Dimensions} and keep patents that contain \emph{all} keywords. The process is repeated for every device type. When a patent is returned for multiple types, I keep a single, randomly chosen assignment to avoid double-counting.





\paragraph{Validation and error modes.}

Keyword linkages can admit complementary (non‐device) technologies and, less often, spurious linkages. Drug‐related patents are the most frequent complements under Cooperative Patent Classification (CPC) labels. Some are economically relevant complements rather than errors: e.g., MRI platform growth can raise demand for imaging agents detailed in the linked patent US-10428030-B2; searching for “cyclosporine test system” returns US-10011612-B2, CPC-labeled as a drug but describing administration that is routinely monitored by cyclosporine tests. In other cases, links are more tenuous (e.g., a contact‐lens query capturing a drug patent because of the drug's contact-lens delivery mechanism). Below, I show that results are robust to excluding such patents.


\paragraph{Restricted samples and robustness.}\label{sec:appendixPatentCollectionRobustness}

According to the USPTO,\footnotemark
\footnotetext{See \href{https://www.uspto.gov/web/offices/ac/ido/oeip/taf/meddev.htm}{https://www.uspto.gov/web/offices/ac/ido/oeip/taf/meddev.htm}.} which tags medical-device patents by technology classes, 210{,}992 patents were designated medical-device-related during 1995--2015. My keyword method yields 875{,}324 patents over a comparable period, likely reflecting inclusion of complementary technologies. I validate the main findings using two restricted samples.


First, I limit to CPC “A61\*” (Medical or Veterinary Science; Hygiene) and exclude “A61P” (drugs), yielding 171{,}682 patents in 1996–2015. I report the top three CPC labels by device type for Class III and Class II devices at \href{https://parkerrogers.github.io/Data/3_to_2_PCS_w_CPCdesc.csv}{this link} and \href{https://parkerrogers.github.io/Data/2_to_1_PCS_w_CPCdesc.csv}{this link}; the primary CPC codes align well with device descriptions.


Second, I screen all USPTO patents with the large language model GPT-4o. For each patent, I supply title, abstract, and primary CPC and ask whether it describes an FDA-regulated device and, if so, to suggest a product code. The prompt, with two calibrating examples, is as follows:

\begin{quote}
As an expert in patents and FDA medical devices, assess the following patent---provided with its title, abstract, and CPC codes---to determine whether it describes a medical device that is regulated by the FDA Center for Devices and Radiological Health. If it does, identify the corresponding FDA product code.

Return your response in the following list format:

[Is Device (0 or 1); FDA Product Code (if applicable); Short Description (if applicable)]

Examples:

- Patent 1:

  Title: Compositions and methods for inhibiting protein on surfaces; 
  Abstract: The use of NIPAM polymers to prevent or reduce the formation of protein deposits on the surfaces of medical devices is described. The invention is particularly directed to the reduction of the adsorption of proteins on surfaces of contact lenses and other medical prosthetics.; 
  Primary CPC: A61L12/08

  Response: [1; LPN; Accessories, soft lens products]

- Patent 2:

  Title: Gene therapy vector system and prodrug genes; 
  Abstract: The present invention relates to a replication retrovirus vector system comprising thymidine kinase (HSV-TK) gene and cytosine deaminase (CD) gene which make gene transfer to cancer cell tissue for the efficient treatment of cancer. This system of the invention can be advantageously used as a pharmaceutical composition for the treatment of cancer.; 
  Primary CPC: A61K48/0058

  Response: [0; ; ]
\end{quote}



Although suggested product codes are often inaccurate, the device/not-device flag is informative. Restricting to patents flagged as devices narrows the sample to 302{,}140 patents during 1996–2015. I also use this screened set to calibrate the patent-topic model so that non-device patents do not drive topics.


Table~\ref{tab:Innovation_nodrugs_A61only} re-estimates Equation~\ref{eq:OLS} using these samples. Panel (a) shows that, for Class III to II events, both GPT-4o and A61 restrictions yield effects that are larger relative to pre-means and more precisely estimated than the unrestricted Column (2) of Table~\ref{tab:Innovation}. For the GPT-4o restriction, the absolute effect is about two-thirds of baseline, suggesting roughly one-third of the baseline reflects spillovers into complementary technologies. Figure~\ref{fig:Innovation_3_2_Nodrugs} shows the corresponding event study under the A61 restriction; the main pattern remains. For Class II to I events (Panel (b)), the GPT-4o restriction implies a large increase in patenting, whereas the A61 restriction yields an economically and statistically insignificant effect. These results indicate that innovation gains for II to I are smaller and more sensitive than for III to II.

Patents are an imperfect innovation measure. Many patents are not commercialized, and large firms may file strategically. As \citet{argente2020patents} show, a substantial share of patents never produces marketable products. My analysis of device-submission rates in the main text, however, complements these patent measures of innovation and yields consistent results. 

\paragraph{Additional patent quality measure.}\label{appendix:breakthrough}

I use the novelty/radicalness index of \citet{KellySeru:2021} to proxy for “breakthrough” patents. The indicator identifies patents that are textually novel relative to prior art and similar to subsequent patents. The distribution is highly skewed in the device-type-year panel (1{,}584 observations): $\approx$81\% zeros; right tail up to 51; mean 0.66; s.d. 3.20. DiD estimates in Table~\ref{tab:BreakthroughPatentsEffects} broadly mirror the main patent‐quality results: breakthrough counts rise post‐reclassification, especially for Class~III to II. Quasi‐Poisson estimates are similar but imprecise because Conley–Taber standard errors are unavailable for non‐linear models and treated units are few; linear specifications with Conley–Taber errors yield tighter intervals.



\paragraph{Identifying patents emphasizing safety advancements.}\label{appendix:safety}

To capture safety-directed innovation, I analyze patent text for safety content. I train a Word2Vec model on all titles and abstracts and select terms closest to “safety” in the embedding space; the resulting set includes “safety,” “safe,” “safer,” “endangering,” “precautions,” “unsafe,” “hazardous,” “failsafe,” “safely,” “dangerous,” “hazard,” “danger,” “risking,” “harming,” “injuring,” “injury,” “jeopardizing,” “risk,” “complication,” and “life-threatening.” A patent is flagged if its title or abstract contains at least one of these terms. Because only about two percent of patents include such terms, I aggregate to the device-type-year level and create an indicator for whether any patent in that cell mentions safety. The measure rises after Class II to I reclassifications. Typical examples include US-10045712-B2 (FDA product code FRZ), which describes an MRI-compatible neonatal table with a safety barrier and automated checks to avoid ferromagnetic objects, and US-10247362-B2 (BXH), which adds over-pressure alarms to medical gas gauges to reduce explosion risk.


\paragraph{Robustness of the safety-mention analysis.}

To confirm that the safety-mention measure is not overly dependent on a single detection method, I compare two alternative definitions. First, a literal-keyword approach flags any occurrence of the string “safety” in the title or abstract. Second, an embedding-distance approach computes the cosine similarity between each patent’s embedding and the centroid of five short sentences describing device safety in different contexts. These five sentences are:


\begin{enumerate}
\item ``The device incorporates multiple safety mechanisms to prevent accidents.''
\item ``Safety protocols must be strictly followed in all operational procedures.''
\item ``Ensuring user safety is paramount in the design of this system.''
\item ``The material is certified for safety under international standards.''
\item ``Safety features include automatic shut-off and emergency alerts.''
\end{enumerate}


Using the 99th similarity percentile as the threshold (results are similar at the 95th), Table~\ref{tab:ProbabilitySafetyMentionEffects_Robustness} shows post-Class~II to I increases under both methods---statistically significant for the keyword measure and marginally significant for the embedding-distance measure---consistent with Table~\ref{tab:AdverseEvent}; no robust increase appears for Class~III to II.

\subsection{Validation of Class II to I Product Safety Results}\label{sec:AppenixSafetyValid}


I corroborate the product safety findings and the role of liability using (i) device recalls and (ii) a 2015 Class~II to I event that did not materially change liability exposure.


\paragraph{Device recalls.}\label{sec:Recalls}

Recalls directly address serious defects, are supervised/mandated by the FDA, are less sensitive to utilization, and are difficult to conceal. Recall data begin in 2000, whereas the primary II to I events occur in 1996 and 1998, so I focus on long-run dynamics. Administrative recall data by device type from 2000–2022 show patterns consistent with the adverse-event and safety‐text analyses. To contextualize the timing, Figure~\ref{fig:InnovationES}, panel (b), shows an immediate rise in available technologies after reclassification; Figure~\ref{fig:ESProductSafety} shows a small, immediate decline in serious events (with hospitalizations significantly lower), followed by a larger decline five to seven years later as patents in Figure~\ref{fig:InnovationES}, panel (c), diffuse.


Figure~\ref{fig:Recall_Analysis} compares average recalls per year for affected device types with the same matched controls used in the adverse-event analysis. Trends (and levels) are similar through about seven years post-event and subsequently diverge, with a persistent decline in recalls for reclassified devices over the subsequent two decades. An event-study specification, with zeros imputed for the years immediately around the events when recall data are unavailable, shows no pre-trend or early effects by construction, no divergence at event times 4–6 when data exist, and a significant decline beginning around year 7 that grows thereafter. The decline of about 0.45 recalls per device-type-year represents roughly a 50\% reduction relative to the control mean.

\paragraph{2015 Class II to I event with no change in liability.}\label{sec:Placebo}

More than 200 Class~II device types were reclassified in 2015 without materially altering tort exposure, while removing 510(k) requirements.\footnotemark \footnotetext{The 2015 event was the first major Class II reclassification since the 1996/8 events analyzed in the main specification.} Before \citetalias{MedtronicvLohr:1996}, 510(k) often preempted state tort claims \citep{flaherty2008defending}; \citetalias{MedtronicvLohr:1996} weakened preemption for Class~II; later, \citetalias{RiegelVMedtronic:2008} restored some protection, but primarily where the FDA imposed device-specific special controls \citep{Costello:2016}.\footnotemark \footnotetext{See, for example, \citetalias{KelseyVAlcon:2019}, involving a contact-lens disinfectant subject to special controls; the court found preemption barred design-defect claims. Similar outcomes have been found for latex gloves, contact lenses, tampons, condoms, angioplasty catheters, wound dressings, tissue adhesive with wound-closure devices, a hemorrhoid-prevention wedge, and electrical-stimulation devices \citep{Munford_2018}.} Most 2015 Class II devices lacked such controls and thus remained exposed to tort liability.\footnotemark \footnotetext{Special controls are typically issued for Class III to II reclassifications.} Using the main matching/estimation strategy, I find no significant changes in hospitalizations, life‐threatening events, deaths, or the composite serious‐events measure (Figure~\ref{fig:SafetyRecalls_IItoI_2015}, Panel (a)); recalls likewise do not change meaningfully for up to eight years (Figure~\ref{fig:SafetyRecalls_IItoI_2015}, Panels (b)--(c)). These null results support liability exposure---not market structure or innovation alone---as the driver of the earlier safety gains. Notably, entry and patenting responses appear larger in 2015 (Figure~\ref{fig:Innovation_2_1_2017}): new‐entrant patenting rises by roughly 39\% within two years (base 11.9 entrants/yr), versus a statistically insignificant 13.0\% rise within two years after the 1996/8 events.\footnotemark \footnotetext{Data limit the horizon to two years post-policy. The 2015 event adds about five new firms per year relative to pre-event means; the 1996/8 events add about one.} Because the 2015 event reduced regulatory costs without raising litigation exposure, innovation benefits may be larger than they were historically.



\subsection{Adverse Event Data: Examples, Coverage, and Limitations}\label{sec:Adverse_Event_Concerns}

MAUDE reports include severity labels such as life-threatening (L), hospitalization (H), disability (S), and death (D). An illustrative spinal-fusion report (H) reads: “It was reported that a [spinal fusion] device has broken whilst implanted... The representative has confirmed the revision has now taken place and that he has retrieved the device.”



The FDA cautions against naïvely using adverse-event counts as a safety measure because utilization is unobserved, reporting can be incomplete, and reports are unverified \citep{CDRHMAUDE}. I address each concern:




First, unobserved utilization can bias adverse-event levels, as increases in adverse events may reflect increases in utilization rather than changes in product safety. The DiD design partly addresses this concern by netting out any time-invariant level differences in reporting across treated and control groups. To further mitigate this potential bias, for a subset of device types that map to procedures, I normalize by procedure volumes and re-run the product safety analysis, finding that the results are robust, if not stronger, to normalization. 

Second, only serious events and manufacturer reports are mandatory. Under-reporting of non-serious events could increase after deregulation, but serious manufacturer reports, the focal product safety outcome, are legally required and subject to enforcement with financial and criminal penalties \citep{FDAMDR2020,bragg_florence_chan_2018,emergo_2022}. Moreover, the 2015 II to I event (no liability change) shows no change in product safety, suggesting reporting rates do not mechanically change after deregulation. 

% Additionally, the FDA's announcement of the Class II to I events mentioned new authorities that enabled closer monitoring of the affected devices to ``take appropriate remedial action if necessary'' \citep{FR60:38903}, indicating that adverse event reports would be harder to conceal.

 Third, reports are not adjudicated. The absence of safety impacts following the 2015 event, the focus on closely monitored serious outcomes, and concordant evidence from safety-focused patent text and recalls mitigate this concern.



Lastly, adverse-event data coverage expanded over time. Manufacturer reports entered MAUDE in 1996, coinciding with early II to I events \citep{MAUDEprimer:2017}. This expansion likely contributed to large relative declines immediately after deregulation in Table~\ref{tab:AdverseEvent}, panel (b), as reporting rose in both treated and control groups. However, splitting adverse events by reporter type shows non-manufacturer serious reports---without the same data coverage shock---also fall after deregulation (Figure~\ref{fig:AdvEvent_by_Reporter}, Panel (a)). Analyzing manufacturer reports alone reproduces the main pattern with wider later confidence intervals, consistent with greater variance as reporting increased (Panel (b)).



\subsection{Device Market Age Construction}\label{sec:AppendixExperience}

This subsection details the construction of the \emph{market age} variable used in Section~\ref{sec:RegulatoryDesign}. I identify the introduction year using hand coding supplemented by GPT‐4o adjudication for Class~II devices. GPT‐4o was provided with a prompt instructing it to report the earliest introduction date and a plausible source (e.g., a patent filing or an early academic study). I then cross‐validated GPT‐4o encodings with a subset of 25 down-regulated device types for which I independently determined introduction years. The correlation coefficient between the hand-coded and GPT‐4o measures was 0.76 ($p<0.001$), lending credibility to the automated approach. The prompt appears below.


\begin{quotation}
\noindent\itshape
``You are an expert in medical devices and have knowledge of when the first device of a given type was introduced. This could be through the earliest academic study that presented the device or the first patent filed by its inventor, whichever came first. I will provide you with a specific medical device, and I would like you to report back the year it was introduced, along with the relevant source.

Here are a few examples:

User: `Sterilizer, Soft-Lens, Thermal, Ac-Powered, Soft (hydrophilic) contact lens care products.'
Response: [`Year: 1974', `source: first patent filed by Cybron Corp. likely used some time before.']

User: `Stimulator, Salivary System, Electrical salivary stimulatory system.'
Response: [`Year: 1986', `source: first submitted PMA. Weiss 1986. subsequent study in 1988 by Steller M.']

User: `Device, Vascular, For Promoting Embolization, Vascular embolization device.'
Response: [`Year: 1975', `Wooly tail embolization coil invented by Cesare Gianturco.']

User: `Intervertebral Fusion Device With Bone Graft, Lumbar, Intervertebral body fusion device.'
Response: [`Year: 1953', `source: The Treatment of Ruptured Lumbar Intervertebral Discs by Vertebral Body Fusion: I. Indications, Operative Technique, After Care.']''
\end{quotation}

The resulting market‐age measure proxies the depth of knowledge about a device’s failure modes and risk profile. Higher values indicate that regulators and manufacturers can rely more on targeted (special or general) controls in lieu of intensive premarket evaluation.



 \clearpage
\newpage

 \section{Supplemental Figures and Tables}\label{sec:SuppFigTab}



    



%  \clearpage
% \newpage

%   \begin{figure}[H]
%   \begin{center}
%   \centering\caption{\textbf{Class II to I Theoretical Change in Safety Effort}}
%    \includegraphics[trim=0cm 0cm 0cm 0cm, scale = .65]{../Figures/LiabilityFunction.png}
%      \centering \caption*{\footnotesize{Note: This figure presents the comparative statics discussed in section \ref{sec:Background} of the change in safety effort after deregulation. Deregulation events considered in this figure are those that move device types from Class II to I. The value $R_{II}$ represents the regulatory stringency of Class II. The value $R_{I}$ represents the stringency of Class I regulations. The x-axis indicates the level of safety effort exerted. The y-axis denotes the dollar values of the marginal abatement of damages from safety effort and the marginal cost of safety effort. The counterfactual dotted section of the Class II marginal abatement curve represents the marginal abatement of damages from exerting effort below mandated levels while still achieving FDA approval. Shifts in the marginal benefit curve occur as fewer damages are precluded through federal preemption when regulatory stringency declines.}}
%   \label{fig:ClassII_I_Safety_Theory}
% \end{center}
% \end{figure}  
    


% \clearpage
% \newpage


% \begin{figure}[!ht]
% \centering \caption{\small{\textbf{Class III to II Theoretical Change in Safety Effort}}}

% \centering \includegraphics[trim=0cm 0cm 0cm 0cm,scale=0.6]{../Figures/Liability3_2.png}
% \caption*{\footnotesize{Note: This figure presents the comparative statics discussed in section \ref{sec:Background} of the change in safety effort with regulation changes. Deregulation events considered in this figure are those that move device types from Class III to II. The x-axis indicates the level of safety effort exerted. The y-axis denotes the dollar values of the marginal abatement of damages from safety effort and the marginal cost of safety effort. The counterfactual dotted section of the marginal abatement curves represents the marginal abatement of damages from exerting effort below mandated levels while still achieving FDA approval. The value $R_{III}$ represents the regulatory stringency of Class III. The value $R_{II}$ represents the stringency of Class II regulations. Shifts in the marginal benefit curve occur as fewer damages are precluded through federal preemption when regulatory stringency declines. Class III device manufacturers do not expect damages as approval guarantees preemption and thus exhibit no returns to safety efforts above mandated levels.}}
% \label{fig:Comparative32_Safety}
% \end{figure}

\clearpage
\newpage
\global\pdfpageattr\expandafter{/Rotate 0}

\begin{figure}[!ht]
    \centering
    \caption{\small{\textbf{Temporal Distribution of Reclassification Events}}}
    \subfloat[Down-Regulation Events]{\includegraphics[trim=0cm 0cm 0cm 1.5cm, clip, scale=0.45]{../replication_package/output/figures/Distribution_of_3_2_over_years.png}}
    \hfill
    \subfloat[Deregulation Events]{\includegraphics[trim=0cm 0cm 0cm 1.5cm, clip, scale=0.45]{../replication_package/output/figures/Distribution_of_2_1_over_years.png}}
    \caption*{\footnotesize{Note: These histograms depict the distribution of down-regulation and deregulation events over time. Panel (a) presents the yearly counts of down-regulation events, while Panel (b) displays the yearly counts of deregulation events.}}
    \label{fig:event_distribution}
\end{figure}
% 243



\clearpage
\newpage
\global\pdfpageattr\expandafter{/Rotate 90}


\begin{sidewaysfigure}[!ht]
\centering \caption{\small{\textbf{Data Catalog}}}

\centering \includegraphics[trim=0cm 0cm 0cm 0cm,scale=0.57]{../replication_package/output/figures/Data_new_latest.png}
\caption*{\footnotesize{Note: This figure presents a catalog of the various data sources used in this study. The three broad outcomes are represented by the three colored boxes: blue innovation, orange market structure, and red product safety. Each broad outcome contains various specific outcomes measured, in most cases, by two different data sources.  Buttons on the exterior represent data sources. The blue arrows connect the data sources to outcome measures. The cogs indicate when algorithms were used to process the data into an outcome measure. The green ``Text Analysis'' cog represents the word2vec algorithm used to extract safety-related keywords from patent claims data.}}
\label{fig:DataCatalog}
\end{sidewaysfigure}
% 281





\clearpage
\newpage
\global\pdfpageattr\expandafter{/Rotate 0}


\begin{figure}[!ht]
\centering \caption{\small{\textbf{Petitioned Down-Regulation Events (Not FDA Initiative)}}}
\centering 
  \begin{overpic}[trim=0cm 0cm 0cm 0cm,scale=0.50]{../replication_package/output/figures/3_2_nolog_real_pats_petitioned_oct22_24.png}
     \put(15,60){\footnotesize Base rate: 7.5 patents/yr}
  \end{overpic}
\caption*{\footnotesize{Note: This figure presents the estimates of the $\beta_k$ coefficients from event-study Equation \ref{eq:EventStudy} for the patent filing rate measure and illustrates the potential biases that stem from industry petition of reclassification. Outcome data are derived from \textit{Dimensions} patent data. Only Class III to II reclassification events petitioned by industry (not by the FDA's own initiative) are considered. Controls are device types matched on baseline averages of the outcome. The coefficient $\beta_{-1}$ is omitted and serves as the reference period. Data are analyzed at an annual frequency. 95\% confidence intervals are calculated following \citet{conley2011inference}.}}
\label{fig:Petitioned_IIItoII}
\end{figure}
% 383

%  \clearpage
% \newpage

% \begin{figure}%
% \centering\caption{\textbf{Mean Yearly Adverse Event Counts by Device Type Class}}
%     \centering\includegraphics[trim=0cm 0cm 0cm 0cm,scale=0.5]{../Figures/Event_to_RegScrutiny.png} 
%   \centering \caption*{\footnotesize{Note: This figure presents the annualized average counts of the specified adverse events for medical device types within the respective classification. The x-axis indicates the device type Class. The x-axis includes reclassified devices from Class III to II and Class II to I events separately. The y-axis details the average annualized count for a given class and adverse event type. The red bar represents the average number of yearly deaths across device types and years. The orange bar calculates a similar average for life-threatening events, and the blue bar calculates the average number of hospitalizations. These three variables are derived from the FDA MAUDE adverse event data. Standard error bands also overlay the average estimates.}}
%   \label{fig:AdverseEventMeansbyClass}

% \end{figure}
% 359



%  \clearpage
% \newpage

%  \begin{figure}
%     \caption{\textbf{FDA Regulatory Hierarchy for Medical Devices}}
% \vspace{-4mm}

%    \includegraphics[trim=0cm 4cm 0cm 0cm,scale = .7]{FDA_Product_Classification.pdf}
% \caption*{\footnotesize{Note: This figure presents the hierarchy with which the FDA regulates medical devices. Devices are regulated at the highest level by ``Parts,'' consisting of 17 advisory committees in distinct areas of medical expertise (e.g., Ophthalmic devices). Within a ``Part,'' devices are categorized into ``Subparts,'' like therapeutic or prosthetic Ophthalmic devices. Medical devices are then further distinguished by ``Regulation Numbers.'' Regulation numbers are typically the level at which regulation occurs and consist of granular device types like soft contact lenses (a therapeutic Ophthalmic device). Lastly, devices are categorized into ``product codes,'' representing the FDA's most granular device categorization (e.g., a daily-wear soft contact lens). Some Class III to II reclassifications target certain product codes within a regulation number.}}
%         \label{pic:DeviceHeirarchy}

% \end{figure}  

% \clearpage
% \newpage

% \global\pdfpageattr\expandafter{/Rotate 90}


% \begin{sidewaysfigure}
% \centering
% \centering \caption{\small{\textbf{Innovation Quality Event Study Class III to II}}}

%  \noindent
%  \subfloat{
%   \includegraphics[scale = .4]{../Figures/3_2_nolog_times_cited_oct22_24.png}
%  }
%   \subfloat{
%   \includegraphics[scale = .4]{../Figures/3_2_nolog_xi_real_2019_oct22_24.png}
%  }

% \caption*{\footnotesize{Note: This figure presents the estimates of the $\beta_k$ coefficients from the event-study equation \ref{eq:EventStudy} for the innovation quality outcomes. Only Class III to II reclassification events are considered. Controls are device types matched on baseline averages of the outcome. The coefficient $\beta_{-1}$ is omitted and serves as the reference period.  Data are analyzed at an annual frequency. The left subfigure describes the evolution of the average citations-per-patent rate. When no patents are filed in a given year, the citations-per-patent rate is set to zero. The right subfigure presents the evolution of the average patent value in treated device types relative to controls. Patent values are derived from \citet{kogan:2017}, who calculate the change in a firm's stock market valuation upon patent grant announcements to measure patent value. Patent values are deflated to millions of 2019 dollars using the CPI. Standard errors are calculated following \citet{conley2011inference}.}}
% \label{fig:InnovationES_IIItoII_Quality}
% \end{sidewaysfigure}
% % 439

% \clearpage
% \newpage

% \global\pdfpageattr\expandafter{/Rotate 90}


% \begin{sidewaysfigure}[!ht]
% \centering \caption{\small{\textbf{Innovation Quality Event Study Class II to I}}}

%  \noindent
%  \subfloat{
%   \includegraphics[scale = .4]{../Figures/2_1_nolog_times_cited_oct22_24.png}
%  }
%   \subfloat{
%   \includegraphics[scale = .4]{../Figures/2_1_nolog_xi_real_2019_oct22_24.png}

%  }\caption*{\footnotesize{Note: This figure presents the estimates of the $\beta_k$ coefficients from event-study equation \ref{eq:EventStudy} for my innovation quality measures. Only Class II to I reclassification events are considered. Controls are device types matched on baseline averages of the outcome. The coefficient $\beta_{-1}$ is omitted and serves as the reference period. Data are analyzed at an annual frequency. The left subfigure describes the evolution of the average citations-per-patent rate. When no patents are filed in a given year, the citations-per-patent rate is set to zero. The right subfigure presents the evolution of the average patent value in treated device types relative to controls. Patent values are derived from \citet{kogan:2017}, who calculate the change in a firm's stock market valuation upon patent grant announcements to measure patent value. Patent values are deflated to millions of 2019 dollars using the CPI.}}
% \label{fig:Innovation_IItoI_Quality}
% \end{sidewaysfigure}
% 439


\clearpage
\newpage








% \clearpage
% \newpage

% \global\pdfpageattr\expandafter{/Rotate 0}


% \begin{figure}[!ht]
% \centering \caption{\small{\textbf{Utilization Rates Event Study}}}

% \centering \includegraphics[trim=0cm 0cm 0cm 0cm,scale=0.5]{../Figures/3_2_nolog_q_solo_CTs.png}
% \caption*{\footnotesize{Note: This figure presents the estimates of the $\beta_k$ coefficients from event-study equation \ref{eq:EventStudy} for the utilization rates of procedures that use treated or control medical device types. I do not have claims data before 2005; Thus, I only consider post-2005 Class III to II reclassification events. Controls are device types matched on baseline average innovation rates. The coefficient $\beta_{-1}$ is omitted and serves as the reference period. Data are analyzed at an annual frequency. Utilization is measured by the yearly number of paid claims for a given procedure. Claims data come from the UCSD healthcare system. Conley--Taber 95\% confidence intervals are provided.}}
% \label{fig:Utilization_IIItoII}
% \end{figure}

% \clearpage
% \newpage


% \global\pdfpageattr\expandafter{/Rotate 0}

% \begin{figure}[!ht]
% \centering \caption{\small{\textbf{Procedure Price Event Study Class III to II}}}

% \subfloat{
%   \includegraphics[width = \textwidth]{../Figures/3_2_nolog_PaidAmt_solo_CTs.png}
%  }


% %  \subfloat{
% %   \includegraphics[scale = .39]{../Figures/3_2_nolog_BilledAmt_solo_CTs.png}
% % }
 
% \caption*{\footnotesize{Note: This figure presents the estimates of the $\beta_k$ coefficients from event-study equation \ref{eq:EventStudy} for the price component of my market structure measures. Controls are device types matched on baseline outcome averages. The coefficient $\beta_{-1}$ is omitted and serves as the reference period. I do not have UCSDH claims data before 2005; Thus, I only consider post-2005 Class III to II reclassification events. Data are analyzed at an annual frequency. The price is determined by the amount insurers pay for a given procedure. The figure describes the evolution of the prices of procedures that use treated device types relative to control groups matched using pre-event price averages. Conley--Taber 95\% confidence intervals are provided.}}
% \label{fig:MktCompostitionPrices_IIItoII}
% \end{figure}






\clearpage
\newpage



\global\pdfpageattr\expandafter{/Rotate 90}




\begin{sidewaysfigure}[!ht]
\centering\caption{\small{\textbf{Reclassification and Market Structure (Incumbent Entry and Other Measures)}}}
\subfloat[New Entry (Device Data), Down-Regulation (III to II)]{
  \begin{overpic}[trim=0cm 0cm 0cm 1.5cm,clip,scale=0.35]{../replication_package/output/figures/3_2_nolog_absnew_firm_count_oct22_24.png}
     \put(15,60){\tiny Base rate: 0.1 firms/yr}
  \end{overpic}
 }\hspace{6mm}
  \subfloat[Incumbent Entry (Device Data), Down-Regulation]{\begin{overpic}[trim=0cm 0cm 0cm 1.5cm,clip,scale=0.35]{../replication_package/output/figures/3_2_nolog_firm_expansion_oct22_24.png}
     \put(15,60){\tiny Base rate: 0.3 firms/yr}
  \end{overpic}
}
 \\
 \noindent 
 \subfloat[Incumbent Entry (Patent Data), Deregulation (II to I)]{ \begin{overpic}[trim=0cm 0cm 0cm 1.5cm,clip,scale=0.35]{../replication_package/output/figures/2_1_nolog_firm_expansion_patents_oct22_24.png}
     \put(15,60){\tiny Base rate: 2.2 firms/yr}
  \end{overpic}
 }\hspace{6mm}
  \subfloat[Incumbent Entry (Patent Data), Down-Regulation]{\begin{overpic}[trim=0cm 0cm 0cm 1.5cm,clip,scale=0.35]{../replication_package/output/figures/3_2_nolog_firm_expansion_patents_oct22_24.png}
     \put(15,60){\tiny Base rate: 1.5 firms/yr}
  \end{overpic}
 }\caption*{\footnotesize{Note: The figure displays the estimates of the $\beta_k$ coefficients from the event-study equation (see Equation \ref{eq:EventStudy}) for incumbent entry and new entry using patent or FDA device submission data. Control groups comprise device types matched on pre-event outcome averages. The $\beta_{-1}$ coefficient is the reference and thus excluded. Analyses utilize annual data. Panel (a) traces new entry among down-regulated devices, measured using FDA device submission data, against matched controls. Panel (b) does the same for incumbent entry. Panel (c) plots changes in incumbent entry among deregulated device types, measured using patent data. Panel (d) does the same for down-regulated device types. Confidence intervals are set at 95\%.}}
\label{fig:MktCompostition_IIItoII_Patents}
\end{sidewaysfigure}
% 390

\clearpage
\newpage

\global\pdfpageattr\expandafter{/Rotate 0}




% \begin{sidewaysfigure}[!ht]
% \centering\caption{\small{\textbf{Reclassification and Market Structure (Incumbent vs. Entrant Patenting)}}}
% \subfloat[Incumbent Patenting, Down-Regulation (III to II)]{
%   \begin{overpic}[trim=0cm 0cm 0cm 1.5cm,clip,scale=0.35]{../Figures/3_2_nolog_patents_incumbent_oct22_24.png}
%      \put(15,60){\tiny Base rate: 6.3 pats/yr}
%   \end{overpic}
%  }\hspace{6mm}
%   \subfloat[Entrant Patenting, Down-Regulation]{\begin{overpic}[trim=0cm 0cm 0cm 1.5cm,clip,scale=0.35]{../Figures/3_2_nolog_patents_entrants_oct22_24.png}
%      \put(15,60){\tiny Base rate: 0 pats/yr}
%   \end{overpic}
% }
%  \\
%  \noindent 
%  \subfloat[Incumbent Patenting, Deregulation (II to I)]{ \begin{overpic}[trim=0cm 0cm 0cm 1.5cm,clip,scale=0.35]{../Figures/2_1_nolog_patents_incumbent_oct22_24.png}
%      \put(15,60){\tiny Base rate: 14.8 pats/yr}
%   \end{overpic}
%  }\hspace{6mm}
%   \subfloat[Entrant Patenting, Deregulation]{\begin{overpic}[trim=0cm 0cm 0cm 1.5cm,clip,scale=0.35]{../Figures/2_1_nolog_patents_entrants_oct22_24.png}
%      \put(15,60){\tiny Base rate: 0 pats/yr}
%   \end{overpic}
%  }\caption*{\footnotesize{Note: The figure displays the estimates of the $\beta_k$ coefficients from the event-study equation (see Eq. \ref{eq:EventStudy}) for incumbent and new entrant patenting (firms that filed patents within the product codes before reclassification vs those that did not). Control groups comprise device types matched on pre-event outcome averages, and use the reclassification year from the treated group to which they are matched as the reference point to identify incumbents and entrants. The $\beta_{-1}$ coefficient is the reference and thus excluded. Analyses utilize annual data. Panel (a) traces incumbent patenting among down-regulated devices, against matched controls. Panel (b) does the same for entrant patenting. Panel (c) and (d) are similar but are for deregulation events. Confidence intervals are set at 95\%.}}
% \label{fig:MktCompostition_IIItoII_Patents_Incumbent}
% \end{sidewaysfigure}
% % 462

% \clearpage
% \newpage



% \global\pdfpageattr\expandafter{/Rotate 0}



% % and Durable Medical Equipment (DME)
% \begin{figure}[!ht]
% \centering\caption{\small{\textbf{Deregulation and Prices (Labs)}}}
% \subfloat[Logged Price Event Study, Labs]{
%   \begin{overpic}[trim=0cm 0cm 0cm 1.5cm,clip,scale=0.45]{../Figures/2_1_nolog_Price_labs_solo_CTs.png}
%      \put(60,62){\tiny Static estimate: $-0.16^*$}
%   \end{overpic}
%  }%\hspace{6mm}
% %   \subfloat[Logged Price Event Study, DME]{\begin{overpic}[trim=0cm 0cm 0cm 1.5cm,clip,scale=0.35]{../Figures/2_1_nolog_Price_DME_solo_CTs.png}
% %      \put(60,62){\tiny Static estimate: $-0.007$}
% %   \end{overpic}
% % }
%  \\
%  \noindent 
%  \subfloat[Logged Price Raw Trends, Labs]{ \begin{overpic}[trim=0cm 0cm 0cm 1.5cm,clip,scale=0.45]{../Figures/2_1_labs_price_raw_trends.png}
%      \put(15,60){}
%   \end{overpic}
%  }%\hspace{6mm}
%   % \subfloat[Logged Price Raw Trends, DME]{\begin{overpic}[trim=0cm 0cm 0cm 1.5cm,clip,scale=0.32]{../Figures/2_1_DME_price_raw_trends.png}
%   %    \put(15,60){}
%   % \end{overpic}
%  % }
%  % \caption*{\footnotesize{Panels (a) and (b) show the estimates of the $\beta_k$ coefficients from the event-study equation (see Eq. \ref{eq:EventStudy}) that measure the impact of deregulation on outpatient procedure prices. Panel (a) focuses on estimates derived from a subset of outpatient lab procedures, while Panel (b) provides estimates based on HCPCS Level II procedures (durable medical equipment, DME). Panels (c) and (d) display the average logged prices for lab procedures and DME, respectively, by treatment group over time. Control groups consist of device types matched to treatment groups based on pre-event outcome averages. The $\beta_{-1}$ coefficient is used as the reference and is therefore omitted. All confidence intervals are set at 95\%.}}
%  \caption*{\footnotesize{Panel (a) shows the estimates of the $\beta_k$ coefficients from the event-study equation (see Eq. \ref{eq:EventStudy}) that measure the impact of deregulation on outpatient procedure prices. Panel (a) focuses on estimates derived from a subset of outpatient lab procedures. Panel (b) displays the average logged prices for lab procedures, by treatment group over time. Control groups consist of device types matched to treatment groups based on pre-event outcome averages. The $\beta_{-1}$ coefficient is used as the reference and is therefore omitted. All confidence intervals are set at 95\%.}}
% \label{fig:restricted_price_analysis}
% \end{figure}
% 470

% \clearpage
% \newpage

% \global\pdfpageattr\expandafter{/Rotate 90}

% \begin{sidewaysfigure}[!ht]
% \centering \caption{\small{\textbf{Market Structure Interactions}}}
%     \includegraphics[trim=0cm 0cm 0cm 0cm,scale=0.5]{../Figures/mkt_structure_interactions.png}
%     \caption*{\footnotesize{Note: This figure illustrates market structure interactions. The first column of panels shows the relationship between price impacts and changes in outcomes, categorized by above- or below-median price impacts for treated device types. The second row of panels display the relationship between entry effects and changes in outcomes, similarly categorized by above- or below-median entry effects for treated device types. The top row focuses on interactions for down-regulation events (Class III to Class II), while the bottom row focuses on interactions for deregulation events (Class II to Class I). Estimates are presented with 95\% confidence intervals and the pre-event mean of the outcome displayed above the dots. The entry effect is defined by above- and below-median changes in new firms patenting.}}
%     \label{fig:MktStructureInt}%
% \end{sidewaysfigure}
% % 472


% \clearpage
% \newpage

% \global\pdfpageattr\expandafter{/Rotate 0}

% \begin{figure}[!ht]
% \centering \caption{\small{\textbf{Market Structure Interactions (Initial Competition)}}}
%     \includegraphics[trim=0cm 0cm 0cm 0cm,scale=0.5]{../Figures/Initial_Competition_Innovation.png}
%     \caption*{\footnotesize{Note: This figure illustrates market structure interactions. The figure shows how estimated impacts on entry and innovation depend on initial levels of competition (above and below the median count of unique firms patenting in treated device types). Panel (a) shows this relationship for deregulation events (Class II to I), while Panel (b) shows the same for down-regulation events (Class III to Class II). Estimates are presented with 95\% confidence intervals. Outcomes are logged (plus one) due to vastly different baseline levels of the outcomes between low- and high-competition environments.}}
%     \label{fig:Initial_Competition_Innovation}%
% \end{figure}
% % 474

\clearpage
\newpage

\global\pdfpageattr\expandafter{/Rotate 0}

\begin{figure}[!ht]
\centering \caption{\small{\textbf{Product Variety, Complexity, and Quality}}}
    \subfloat[Change in Patent Topics, Down-Regulation]{\includegraphics[trim=0cm 0cm 0cm 0.8cm,clip,scale=.45]{../replication_package/output/figures/variety_3_2.png}}

    \subfloat[Change in Patent Topics, Deregulation]{\includegraphics[trim=0cm 0cm 0cm 0.8cm,clip,scale=.45]{../replication_package/output/figures/variety_2_1.png}}

    \caption*{\footnotesize{Note: The figure shows estimated percent changes in patents by topic following reclassification, using Equation \ref{eq:OLS} estimated separately for each topic cluster identified via LDA topic modeling on patent texts. Panel (a) shows the effects of down-regulation; Panel (b) shows the effects of deregulation. DiD estimates compare changes in topic emphasis between treated and matched control units before and after reclassification. 95\% confidence intervals are shown.}}
    \label{fig:ProductVariety}%
\end{figure}
% 405

\clearpage
\newpage

\pdfpageattr\expandafter{/Rotate 0}



\begin{figure}[!ht]
\centering
\centering \caption{\small{\textbf{Effects on Innovation by Experience}}}

\subfloat[Cumulative Approval Time (Baseline)]{
  \includegraphics[width = 0.48\textwidth]{../replication_package/output/figures/LBDEstimation.png}
}
\subfloat[Cumulative Product Count (Robustness)]{
  \includegraphics[width = 0.48\textwidth]{../replication_package/output/figures/LBDEstimation_ProductCount.png}
}







% \centering \caption{\small{\textbf{Estimated Learning Curves and Simulated Experience-Specific Effects}}}
%    \begin{minipage}{0.5\textwidth}
%         \centering
%         \caption*{Panel A: Estimated Learning Curves}
%             % \vspace{-10pt}

%         \includegraphics[width=\textwidth]{../Figures/Simulation.png} % first figure itself
%     \end{minipage}\hfill
%     \begin{minipage}{0.5\textwidth}
%         \centering
%         \caption*{Panel B: Simulated Change in New Devices by Proficiency}
%             % \vspace{-10pt}


%         \includegraphics[width=\textwidth]{../Figures/LBDsimulationIIItoII.png} % second figure itself
%     \end{minipage}
    \caption*{\footnotesize{Note: This figure illustrates how Class~III to~II reclassification impacts the rates of new device submissions, segmented by firm experience. Panel~(a) measures proficiency as cumulative days in the FDA approval process; Panel~(b) uses cumulative product count---the total number of prior FDA submissions by each firm. Both panels report DiD estimates of percent changes in new device submission rates for treated device types relative to controls, categorized by experience terciles. Baseline submission rates are similar across terciles; percent changes normalize each tercile's estimate by its pre-event baseline, facilitating comparison across bins. The estimation excludes zero-experience observations to avoid biases from extensive-margin entry decisions. Conley--Taber 95\% confidence intervals are shown.}}
    \label{fig:FirmExperienceHeteroSimulation}%
\end{figure}
% 425




% \clearpage
% \newpage

% \global\pdfpageattr\expandafter{/Rotate 0}

% \begin{figure}%
% \centering
% \caption{\textbf{Percent Change in New Devices After Down-Regulation by PMA and 510(k) Experience Levels}}
%     \centering\includegraphics[trim=0cm 0cm 0cm 1.5cm,clip,scale=0.65]{../Figures/LBDEstimation_Separated.png} 
%   \centering \caption*{\footnotesize{Note: This figure illustrates the percent change in new medical devices submitted after down-regulation according to prior PMA (Premarket Approval) and 510(k) experience levels. The y-axis represents prior PMA Experience (High or Low), and the x-axis represents 510(k) Experience (High or Low). Each cell within the heatmap displays the estimated percentage change in new devices, with standard errors presented in parentheses below each estimate. Significance levels are indicated by stars: + for p$<$ 0.1, * for p$<$ 0.05, and ** for p$<$ 0.01. The color intensity of each cell reflects the magnitude of the percentage change, with darker shades representing higher changes.}}
%   \label{fig:PercentChangeByExperience}
% \end{figure}
% 498


% \clearpage
% \newpage

% \begin{figure}[!ht]
% \centering \caption{\small{\textbf{Market Composition Event-Study Class II to I}}}

% \subfloat{
%   \includegraphics[scale = .39]{../Figures/2_1_nolog_unique_firms_solo_CTs.png}
%  }


%  \subfloat{
%   \includegraphics[scale = .39]{../Figures/2_1_nolog_firm_expansion_patents_solo_CTs.png}
% }\caption*{\footnotesize{Note: This figure presents the estimates of the $\beta_k$ coefficients from event-study equation \ref{eq:EventStudy} for my market composition measures. Only Class II to I reclassification events are considered. Data are analyzed at an annual frequency. Controls are device types matched on baseline outcome averages. The coefficient $\beta_{-1}$ is omitted and serves as the reference period. The top subfigure illustrates the evolution of the rate of new firm entry (measured with patent data) relative to matched control groups. New firm entry represents firms that have never been granted a patent. The bottom subfigure illustrates the evolution of the rate of incumbent firm entry (firms that have received a granted patent), entering treated device types relative to matched controls. I do not include FDA-approved device measures of new and incumbent entry as I do not have reliable data on new Class I devices from FDA sources. 95\% confidence intervals are provided.}}
% \label{fig:MarketComposition_IItoI}
% \end{figure}


% \clearpage
% \newpage

% \global\pdfpageattr\expandafter{/Rotate 0}

% \begin{figure}[!ht]
% \centering \caption{\small{\textbf{Flattening the Learning Curve Simulation}}}
%     \includegraphics[trim=0cm 0cm 0cm 0cm,scale=0.5]{../Figures/LBDGammaSim.png}
%     \caption*{\footnotesize{Note: This figure presents the simulation exercise of flattening the Class III learning curve estimated in equation \ref{eq:loglinlearning}. I flatten the learning curve relative to the most experienced firm. The results of this simulation are provided in Table \ref{tab:FlattenCurveSimul}. Above, $\gamma$ begins at its initial starting point estimated in equation \ref{eq:loglinlearning}. Subsequent lines show the change in the learning curve as $\gamma$ is reduced while maintaining the approval time of the top quartile of experienced firms. $T_{Sum,25}$ represents the bottom 25th percentile of cumulative FDA experience (in days), $T_{Sum,50}$ represents the 25-50th percentile, and $T_{Sum,75}$ represents the 50-75th percentile.}}
%     \label{fig:GammaSimulation}%
% \end{figure}
% 501



\clearpage
\newpage



\pdfpageattr\expandafter{/Rotate 90}

    \begin{sidewaysfigure}
    % \centering  \caption{Effect of reclassification on Patenting Rates by Firm Asset Terciles}
\captionsetup[subfigure]{labelformat=empty}

\centering
\centering \caption{\small{\textbf{Effect of Reclassification on Patenting Rates by Asset Terciles}}}
 \subfloat[Panel (a): Class III to II]{
 % \vspace{-20pt}
  \includegraphics[width = 0.48\textwidth]{../replication_package/output/figures/3_2_smallvlarge.png}
 }
 \subfloat[Panel (b): Class II to I]{
  \includegraphics[width = 0.48\textwidth]{../replication_package/output/figures/2_1_smallvlarge.png}
}
 % \\
 % \noindent
 % \subfloat[Panel C: Simulated Class III to II]{
 %  \includegraphics[scale = .3]{../Figures/FirmSizesimulationIIItoII.png}
 % }
 %  \subfloat[Panel D: Simulated Class II to I]{
 %  \includegraphics[scale = .3]{../Figures/FirmSizesimulationIItoI.png}
 % }

    % \begin{minipage}{0.5\textwidth}
    %     \centering
    %     \caption*{Panel A: Class III to II}
    %         \vspace{-14pt}

    %     \includegraphics[width=\textwidth]{../Figures/3_2_SideGig_smallvlarge.png} % first figure itself
    % \end{minipage}\hfill
    % \begin{minipage}{0.5\textwidth}
    %     \centering
    %     \caption*{Panel B: Class II to I}
    %         \vspace{-14pt}


    %     \includegraphics[width=\textwidth]{../Figures/2_1_SideGig_smallvlarge.png} % second figure itself
    % \end{minipage}
     \caption*{\footnotesize{Note: This figure presents the DiD estimates from Equation \ref{eq:OLS} for the patenting rate across reclassification types and firm asset terciles. I exclude patent data for private firms since I only observe firm asset data for publicly traded firms. Panel (a) presents the change in patenting rates in my Class III to II treated medical device types, relative to matched control groups, across asset terciles. The first tercile represents the bottom 33rd percentile of assets, the second represents the 33rd--66th percentile, and the third represents the 66th--100th percentile.  Panel (b) presents the change in patenting rates in my Class II to I treated medical device types, relative to matched control groups, across asset terciles. 
     95\% confidence intervals overlay the estimates.}}
     \label{fig:FirmSizeHetero}%
\end{sidewaysfigure}
% 428















% \clearpage
% \newpage
% \global\pdfpageattr\expandafter{/Rotate 90}


% \begin{sidewaysfigure}[!ht]
% \centering \caption{\small{\textbf{Adverse Event Event Study Class III to II}}}

% \subfloat{
%   \includegraphics[scale = .35]{../Figures/3_2_nolog_L_oct22_24.png}
%  }
%  \subfloat{
%   \includegraphics[scale = .35]{../Figures/3_2_nolog_H_oct22_24.png}
% }
%  \\
%  \noindent
%  \subfloat{
%   \includegraphics[scale = .35]{../Figures/3_2_nolog_D_oct22_24.png}
%  }
%   \subfloat{
%   \includegraphics[scale = .35]{../replication_package/output/figures/3_2_nolog_v_serious_oct22_24.png}

%  }\caption*{\footnotesize{Note: This figure presents the estimates of the $\beta_k$ coefficients from event-study equation \ref{eq:EventStudy} for my adverse event measures. Only Class III to II reclassification events are considered. Data are analyzed at an annual frequency. Controls are device types matched on baseline outcome averages. The coefficient $\beta_{-1}$ is omitted and serves as the reference period. The top-left subfigure describes the evolution of the rate of life-threatening events stemming from the use of treated device types relative to control groups matched using baseline averages. The top-right subfigure describes the evolution of the rate of hospitalizations of treated device types relative to control groups. The bottom-left subfigure describes the evolution of the death rate. The bottom-right subfigure presents the evolution of the sum of all serious adverse events (life-threatening, death, hospitalizations, and disability) in treated device types relative to controls. Adverse events are derived from the FDA MAUDE database. Conley--Taber 95\% confidence intervals are provided.}}
% \label{fig:AdverseES_IIItoII}
% \end{sidewaysfigure}
% % 521

% \clearpage
% \newpage


% \begin{sidewaysfigure}[!ht]
% \centering \caption{\small{\textbf{Adverse Event Event Study Class II to I}}}

% \subfloat{
%   \includegraphics[scale = .35]{../Figures/2_1_nolog_L_oct22_24.png}
%  }
%  \subfloat{
%   \includegraphics[scale = .35]{../Figures/2_1_nolog_H_oct22_24.png}
% }
%  \\
%  \noindent
%  \subfloat{
%   \includegraphics[scale = .35]{../Figures/2_1_nolog_D_oct22_24.png}
%  }
%   \subfloat{
%   \includegraphics[scale = .35]{../replication_package/output/figures/2_1_nolog_v_serious_oct22_24.png}

%  }\caption*{\footnotesize{Note: This figure presents the estimates of the $\beta_k$ coefficients from event-study equation \ref{eq:EventStudy} for my adverse event measures. Only Class II to I reclassification events are considered. Only four pre-periods are included because there are no prior adverse event data. Data are analyzed at an annual frequency. Controls are device types matched on baseline outcome averages. The coefficient $\beta_{-1}$ is omitted and serves as the reference period. The top-left subfigure describes the evolution of the rate of life-threatening events stemming from the use of treated device types relative to control groups matched using baseline averages. The top-right subfigure illustrates the evolution of the rate of hospitalizations of treated device types relative to matched control groups. The bottom-left subfigure describes the relative evolution of the death rate. The bottom-right subfigure presents the relative evolution of the sum of all serious adverse events (life-threatening, death, hospitalizations, and disability) in treated device types. Adverse events are derived from the FDA MAUDE database. 95\% confidence intervals are provided.}}
% \label{fig:AdverseEvents_IItoI}
% \end{sidewaysfigure}
% 521


\clearpage
\newpage

\global\pdfpageattr\expandafter{/Rotate 0}

\begin{figure}[!ht]
\centering \caption{\small{\textbf{Impact of Class II to I Events on Safety Emphasis}}}

\subfloat[Safety Emphasis, Raw Trends]{\includegraphics[trim=0cm 0cm 0cm 1.5cm,clip,width = 0.68\textwidth]{../replication_package/output/figures/Safety_emphasis_raw_trends.png}}

    \subfloat[Safety Emphasis, Event Study]{\includegraphics[trim=0cm 0cm 0cm 1.5cm,clip,width = 0.7\textwidth]{../replication_package/output/figures/2_1_nolog_safety_prop_oct22_24.png}}

\caption*{\footnotesize{Note: Panel (a) displays raw trends in the average probability of observing mentions of safety in patents for treated device types compared to control device types. Panel (b) presents estimates of the $\beta_k$ coefficients from event-study Equation \ref{eq:EventStudy}, capturing inventors' emphasis on safety. The analysis focuses exclusively on Class II to Class I reclassification events, using annual data. Control device types are matched based on baseline averages of the innovation outcome variables. The coefficient $\beta_{-2}$ is omitted as the reference period because of noise in the pre-event data. The figure illustrates the evolution of the proportion of patents emphasizing safety in their text, with 95\% confidence intervals included.}}
\label{fig:Safety_IItoI}
\end{figure}
% 443


\clearpage
\newpage

\global\pdfpageattr\expandafter{/Rotate 0}


% [Moved to main text: panels (c)--(d) of fig:ESProductSafety]
% Previously Figure D.9 (Normalizable Subset, Raw Counts vs. Normalized)
\clearpage
\newpage

\pdfpageattr\expandafter{/Rotate 0}

\begin{figure}
\begin{center}

   \caption[\small{\textbf{Evolution of Medical Device Legal Environment}}]{\textbf{Evolution of Medical Device Legal Environment}}\label{fig:DeviceTortCasesOverTime}

 % First Row: Two Subfigures Side-by-Side with Captions on Top
    \subfloat[Medical Device and Drug Tort Cases]{%
        \includegraphics[width=0.65\textwidth]{../replication_package/output/figures/cases_preemption.png}
    }
    
    \vspace{1em} % Adds vertical space between rows
    
    % Second Row: Single Subfigure with Caption on Top
    \subfloat[Ratio of Bankruptcy Cases to Tort Claims for Medical Device Firms]{%
        \includegraphics[trim=0cm 0cm 0cm 1.5cm, clip, width=0.65\textwidth]{../replication_package/output/figures/Bankruptcy_1996_raw_trend.png}
    }
    

    

\end{center}

\caption*{\footnotesize{Note: Panel (a) depicts yearly counts of medical device tort claims and highlights those referencing ``preemption.'' It also includes counts for pharmaceutical tort cases mentioning ``preemption.'' Data for this panel are sourced from Nexis Uni and are based on searches for ``Medical device tort,'' ``medical device tort preempt,'' and ``pharmaceutical tort preempt.'' The vertical dotted line marks 1996, the year of the Supreme Court's Medtronic, Inc. v. Lohr decision, which significantly impacted the regulatory environment by withdrawing legal protections for Class II devices. Panel (b) plots the ratio of bankruptcy case counts to tort claim counts in the medical device industry by year. Data for this panel are sourced from Nexis Uni and are based on searches for ``medical device Chapter 11 bankruptcy.''}}
\end{figure}  
% 443


\clearpage
\newpage


\global\pdfpageattr\expandafter{/Rotate 90}

\begin{sidewaysfigure}[!ht]
\centering
\caption{\small{\textbf{How Impacts Vary with Device Market Age}}}

% Top row: Class III -> II
\subfloat[Down-Regulation: Serious-Event Effect]{
  \includegraphics[scale = .32]{../replication_package/output/figures/exp_vs_impact_v_serious_est_3_2.png}
}
\subfloat[Down-Regulation: Patenting Effect]{
  \includegraphics[scale = .32]{../replication_package/output/figures/exp_vs_impact_patent_est_3_2.png}
}
\\
% Bottom row: Class II -> I
\subfloat[Deregulation: Serious-Event Effect]{
  \includegraphics[scale = .32]{../replication_package/output/figures/exp_vs_impact_v_serious_est_2_1.png}
}
\subfloat[Deregulation: Patenting Effect]{
  \includegraphics[scale = .32]{../replication_package/output/figures/exp_vs_impact_patent_est_2_1.png}
}

\caption*{\footnotesize{Note: Each panel plots device-level DiD estimates against market age at reclassification. Points are covariate-adjusted predictions in the original units, evaluated at representative values for other regressors (continuous controls at sample means; indicators at sample shares; reference product area). Curves show the conditional OLS fit. The dashed vertical line marks the policy-relevant cutoff---the age at which the fitted safety effect equals zero. Panels (a)–(b) show down-regulation (Class III to II); panels (c)–(d) show deregulation (Class II to I).
}}
\label{fig:Impact_by_experience}
\end{sidewaysfigure}
% 495



\clearpage
\newpage

\pdfpageattr\expandafter{/Rotate 0}


\begin{figure}[htbp]
    \centering
    \caption{\textbf{Histogram, Market Age of Current Class III and II Devices}}\label{fig:DeviceExperienceCurrent}
    
    % First Row: Two Subfigures Side-by-Side
    \subfloat[Class III Devices]{%
        \includegraphics[trim=0cm 0cm 0cm 1.5cm, clip, width=0.7\textwidth]{../replication_package/output/figures/Distribution_of_3_over_years.png}
    }

    % \hspace{0.05\textwidth} % Horizontal space between subfigures
    \subfloat[Class II Devices]{%
        \includegraphics[trim=0cm 0cm 0cm 1.5cm, clip, width=0.7\textwidth]{../replication_package/output/figures/Distribution_of_2_over_years.png}
    }
    
    \vspace{1em} % Adds vertical space before the note
    
    \caption*{\footnotesize{Note: Panel (a) illustrates the distribution of market age for current Class III medical devices. Panel (b) presents a similar distribution for Class II medical devices. The data were collected using GPT-4o to gather invention dates for each device. The correlation coefficient between market age determined through a hand-coded correspondence for a subset of 25 down-regulated device types and that generated by GPT-4o is 0.76, with $p < 0.001$, indicating a strong positive correlation between the two measures.
    }}
\end{figure}
% 497

% \clearpage
% \newpage
% \global\pdfpageattr\expandafter{/Rotate 0}


% \begin{figure}[!ht]
% \centering \caption{\small{\textbf{Denovo Filings By Treatment Status}}}
% \centering 
%   \begin{overpic}[trim=0cm 0cm 0cm 1.5cm,clip,scale=0.50]{../Figures/Denovos_stability.png}
%      \put(40,60){}
%   \end{overpic}
% \caption*{\footnotesize{Note: This figure displays the annual number of de novo filings, categorized by whether a filing pertains to a new device type in the same part or subpart as a treatment device type or to one in the same part or subpart as a control device type.
% }}
% \label{fig:denovo_stability}
% \end{figure}
% % 1538



%   \begin{figure}[H]
%   \begin{center}
%   \centering\caption{\textbf{Theoretical Change in Safety Effort after Deregulation}}
%    \includegraphics[width = .9\textwidth]{../Figures/Safety_Effort_Theory.png}
%      \centering \caption*{\footnotesize{Note: This figure presents a possible change in the level of safety effort after deregulation. This scenario is one in which deregulation could lead to an increase in safety efforts, given a sufficient increase in damages, as described in Section \ref{sec:Background}. The x-axis indicates the level of safety effort exerted. The y-axis denotes the monetary value. The marginal cost of effort curve indicates a marginal cost of $\psi$ at initial values of safety effort before financing costs are incurred, at which point marginal costs increase with effort. The marginal abatement of damages curve under regulation is always equal to zero due to federal preemption. The counterfactual dotted section of the marginal abatement curve under regulation represents the marginal abatement of damages from exerting effort below mandated levels while still achieving FDA approval. Deregulation shifts the marginal abatement curve as legal damages are no longer prevented by federal preemption. The value $x^*$ represents the optimal level of safety effort after deregulation (i.e., where the marginal cost of safety effort is equal to the marginal abatement of expected damages). The value $\underline{x}$ represents the mandated level of safety effort. The vector $\vec{Z}$ contains other factors that affect a firm's legal damages in expectation, which might be specific to the given legal system, like damage caps.}}
%   \label{fig:Safety_Theory}
% \end{center}
% \end{figure}  
% 1712












\clearpage
\newpage


% \pdfpageattr\expandafter{/Rotate 90}


% \begin{sidewaysfigure}
%       \centering
% \centering
% \caption[short for lof]{Class II to I Changes in Adverse Event Rates at Margin of Decision Rule}

%  {
%   \includegraphics[width=0.49\linewidth]{../Figures/DPM_Deaths.png}
%   \label{subfig:fig1}
%  }
%  {
%   \includegraphics[width=0.49\linewidth]{../Figures/DPM_LifeThreat.png}
%   \label{subfig:fig2}
% }
 
% {
%   \includegraphics[width=0.49\linewidth]{../Figures/DPM_Hospital.png}
%   \label{subfig:fig3}
%  }

%  \caption*{\footnotesize{Note: This figure presents separate DiD estimates of equation \ref{eq:OLS} for each adverse event measure and each treated device type with a given proxy DPM score relative to matched controls. The DPM score is primarily an increasing function of the baseline average annual incidence of adverse events before deregulation (see Appendix \ref{sec:FDADecisionRule}). When the rightmost outlier is removed, the slopes of the fitted lines are still negative or zero. Controls for each treated device type are selected by matching based on DPM scores across both Class I and II devices that were not down-classified in the given period. The x-axes describe the same proxy DPM score across the three adverse event outcomes. The y-axes describe the change in the rate of the given adverse event type in the treated device type relative to the matched control device type. The top-left figure shows the difference-in-differences estimates for the change in death rates across device types, the top-right figure shows the same for life-threatening events, and the bottom figure shows the same for hospitalizations. Adverse event data are from the FDA's MAUDE database. 95\% confidence intervals are provided.}}
%          \label{fig:IItoIAdverseTopDecile}%

%     \end{sidewaysfigure}




% \pdfpageattr\expandafter{/Rotate 0}

% \begin{figure}
%      \centering   \caption{\textbf{Contact Lens Use Case---III to II reclassification}}
%   \centering \begin{overpic}[trim=0cm 0cm 0cm 0cm,scale=0.5]{../Figures/Contact_Firms.png}
%      \put(15,30){\includegraphics[scale=0.27]{../FDA_Presentation_Latest/cooperVision}} 
%      \put(60,20){\includegraphics[scale=0.25 ]{../FDA_Presentation_Latest/dailies}}  
%   \end{overpic}

% \caption*{\footnotesize{Note: This figure presents an example of a Class III to II reclassification event. In 1994, the FDA down-classified daily-wear soft contact lenses to Class II but kept extended-wear soft contact lenses in Class III. The x-axis measures the year, and the y-axis measures the number of unique contact lens devices submitted to the FDA for approval in a given year. The green line represents daily-wear contact lenses submitted for approval/clearance (deregulated), and the blue line represents extended-wear soft contact lenses submitted for approval (remained in Class III). The vertical black line represents the year of reclassification. The left-imposed picture shows an example of a soft contact lens invented before reclassification. The right-imposed picture shows an example of a soft contact lens invented after reclassification.}}
% \label{fig:ContactLensUseCase}

% \end{figure}  





\clearpage
\newpage

\global\pdfpageattr\expandafter{/Rotate 0}


\begin{figure}[!ht]
\centering \caption{\small{\textbf{Effects of Class III to II Events on Patenting Rates:\\Restricted Patent Sample}}}

\centering \includegraphics[trim=0cm 0cm 0cm 0cm,scale=0.5]{../replication_package/output/figures/3_2_nolog_patents_pats_nodrugs_oct22_24.png}
\caption*{\footnotesize{Note: This figure presents the estimates of the $\beta_k$ coefficients from event-study Equation \ref{eq:EventStudy} for patenting rates using the restricted patent sample described in Appendix \ref{sec:appendixPatentCollection}. Compare to Figure \ref{fig:InnovationES}, Panel (a). Controls are device types matched on baseline average innovation rates. The coefficient $\beta_{-1}$ is omitted and serves as the reference period. Data are analyzed at an annual frequency. The patenting rate is measured by the yearly number of patents filed in a given device type. Patent data come from the \textit{Dimensions} patent database. Conley--Taber 95\% confidence intervals are provided.}}
\label{fig:Innovation_3_2_Nodrugs}
\end{figure}
% 1704

\clearpage
\newpage

\pdfpageattr\expandafter{/Rotate 0}

\begin{figure}
 \centering
    \caption{\textbf{Device Recall Analysis---Class II to I}}\label{fig:Recall_Analysis}
    
    % First Row: Two Subfigures Side-by-Side with Captions on Top
    \subfloat[Raw Trend, Matched Controls]{%
        \includegraphics[trim=0cm 0cm 0cm 1.5cm, clip, width=0.49\textwidth]{../replication_package/output/figures/Recalls_1996_raw_trend.png}
    }\hfill
    \subfloat[Raw Trend---Intuitive Controls]{%
        \includegraphics[trim=0cm 0cm 0cm 1.5cm, clip, width=0.49\textwidth]{../replication_package/output/figures/Recalls_1996_raw_trend_intuitive.png}
    }
    
    \vspace{1em} % Adds vertical space between rows
    
    % Second Row: Single Subfigure with Caption on Top
    \subfloat[Recall Count Event Study, Matched Controls]{%
        \includegraphics[trim=0cm 0cm 0cm 1.5cm, clip, width=0.65\textwidth]{../replication_package/output/figures/2_1_no_log_recalls_1996.png}
      
    }
    

\caption*{\footnotesize{Note: This figure illustrates the device recall analysis for the Class II to I events. Panels (a) and (b) display the average device recall counts by treatment status from 2000 to 2023 for the matched and intuitive control groups, respectively. Panel (c) shows the event-study estimates of the change in device recalls over time. Recall counts for event times $-2 \leq t \leq 3$ are constant at zero due to the missing recall data for those specific years.}}
\end{figure}  
% 1746



\clearpage
\newpage

\pdfpageattr\expandafter{/Rotate 0}


% Requires: \usepackage{graphicx, subfig}
\begin{figure}[htbp]
\centering
\caption{\small{\textbf{Safety and Recalls After Deregulation---Class II to I (2015)}}}
\label{fig:SafetyRecalls_IItoI_2015}

% Top row: (a) serious AEs, (b) recalls event study
\begin{minipage}[t]{0.49\textwidth}
  \centering
  \subfloat[Serious adverse events (event study)]{
    \includegraphics[trim=0cm 0cm 0cm 1.5cm,clip,width=.9\linewidth]{../replication_package/output/figures/2_1_no_logv_serious_2016_oct22_24.png}
  }
\end{minipage}\hfill
\begin{minipage}[t]{0.49\textwidth}
  \centering
  \subfloat[Device recalls (event study)]{
    \includegraphics[trim=0cm 0cm 0cm 1.5cm,clip,width=.9\linewidth]{../replication_package/output/figures/Recalls_2017.png}
  }
\end{minipage}

\vspace{0.8em}

% Bottom row: (c) raw recall trends (full width)
\begin{minipage}[t]{\textwidth}
  \centering
  \subfloat[Device recalls (raw mean trends by treatment status)]{
    \includegraphics[trim=0cm 0cm 0cm 1.5cm,clip,width=0.5\textwidth]{../replication_package/output/figures/Recalls_2017_raw_trends.png}
  }
\end{minipage}

\caption*{\footnotesize
Note: Panels (a)–(b) report event-study coefficients $\beta_k$ from Equation~\ref{eq:EventStudy} for the 2015 Class II to I reclassification; $\beta_{-1}$ is omitted (reference). Outcomes are annual. Treated device types are compared to matched controls based on pre-event baseline outcome averages. Panel (a) shows serious adverse events---life-threatening, death, hospitalization, and disability---from FDA's MAUDE database. Panel (b) shows device recalls. Panel (c) plots raw mean recall counts by treatment status (2000--2023). 95\% confidence intervals are shown in Panels (a)–(b).
}
\end{figure}
% 1750

\clearpage
\newpage

\global\pdfpageattr\expandafter{/Rotate 0}


\begin{figure}[!ht]
\centering \caption{\small{\textbf{Effects of 2015 Class II to I Event on New Entry}}}\label{fig:Innovation_2_1_2017}


\centering \includegraphics[trim=0cm 0cm 0cm 0cm,scale=0.5]{../replication_package/output/figures/NewFirmsPatenting_2017.png}
\caption*{\footnotesize{Note: This figure presents the estimates of the $\beta_k$ coefficients from event-study Equation~\ref{eq:EventStudy}. Controls are device types matched on baseline average innovation rates. The coefficient $\beta_{-1}$ is omitted and serves as the reference period. Data are analyzed at an annual frequency. The entry rate is measured by the yearly number of new firms patenting in a given device type. Patent data come from the \textit{Dimensions} patent database. 95\% confidence intervals are provided.}}
\end{figure}
% 1750




% \clearpage
% \newpage
% \global\pdfpageattr\expandafter{/Rotate 0}


% \begin{figure}[!ht]
% \centering \caption{\small{\textbf{Effect of Reclassification on Product Safety \\(Hospital Beds Only)}}}
% \centering 
%   \begin{overpic}[trim=0cm 0cm 0cm 0cm,scale=0.50]{../Figures/Hospital_bed.png}
%      \put(40,60){\footnotesize Base rate: 50 events/yr}
%   \end{overpic}
% \caption*{\footnotesize{Note: This figure presents the estimates of the $\beta_k$ coefficients from event-study equation \ref{eq:EventStudy} for the serious adverse event count measure. Only Class the hospital bed treatment group is considered for Class II to I deregulation events are considered. Controls are all device types matched on baseline averages of the outcome (for all treated groups). The coefficient $\beta_{-1}$ is omitted and serves as the reference period. Data are analyzed at an annual frequency. 95\% confidence intervals are calculated following \citet{conley2011inference}.}}
% \label{fig:HospitalBedExample}
% \end{figure}
% % 2227

\clearpage
\newpage

\global\pdfpageattr\expandafter{/Rotate 0}

\begin{figure}
\begin{center}

   \caption[\small{\textbf{Serious Adverse Event Changes by Reporter Type---Class II to I}}]{\textbf{Serious Adverse Event Changes by Reporter Type---Class II to I}}\label{fig:AdvEvent_by_Reporter}
  \subfloat[Serious Adverse Events, Excluding Manufacturers]{\includegraphics[trim=0cm 0cm 0cm 1.5cm,clip,scale=0.4]{../replication_package/output/figures/2_1_nolog_v_serious_voluntary_oct22_24.png}}

    \subfloat[Serious Adverse Events, Manufacturers]{\includegraphics[trim=0cm 0cm 0cm 1.5cm,clip,scale=0.4]{../replication_package/output/figures/2_1_nolog_v_serious_manuf_oct22_24.png}}
\end{center}

\caption*{\footnotesize{Note: The figure displays the estimates of the $\beta_k$ coefficients from the event-study equation (see Equation \ref{eq:EventStudy}) for serious adverse events, encompassing life-threatening incidents, mortality, hospitalizations, and disability, differentiated by the source of the adverse event report. Results are presented separately for manufacturer reports and reports from all other sources. Control groups consist of device types matched on pre-event outcome averages. The $\beta_{-1}$ coefficient is the reference and thus excluded. Analyses utilize annual data. Panel (a) traces serious adverse events for deregulated devices against matched controls based on reports from non-manufacturer sources. Panel (b) offers a parallel analysis but exclusively considers serious adverse event reports originating from manufacturers. Confidence intervals are set at 95\%.}}


\end{figure}  
% 1775













% \clearpage
% \newpage

% \begin{figure}%
%  \centering  \caption{Small vs. Large Firm Innovativeness}

%      \centering{\includegraphics[scale=0.2]{../Figures/Patent_Citation_Dist.png} }
  
%   \caption*{\footnotesize{Note: This figure presents the relationship between firm size and innovativeness as measured by citations-per-patent. Data are binned according to an optimal split criterion. The dotted green line represents a nonlinear least-squares regression line. The x-axis represents a given firms' inflation-adjusted total assets (in millions). The y-axis represents the number of patent citations for a given patent. Each dot represents a binned average of citations for a given firm size.}}
%      \label{fig:SmallvLargeCitations}%
% \end{figure}


% \clearpage
% \newpage

% \pdfpageattr\expandafter{/Rotate 0}



% \begin{figure}[!ht]
% \centering \caption{\small{\textbf{Binscatter Regression of Log Cumulative Experience vs. Residualized Days to Approval}}}

% \centering \includegraphics[trim=0cm 0cm 0cm 0cm ,scale=0.2]{../Figures/LBDResidPlot.png}
% \caption*{\footnotesize{Note: This figure presents the estimated regression line, which estimates equation \ref{eq:loglinlearning} using OLS, together with a binscatter of residualized approval times against the log of the submitting firm's cumulative experience. Residualized approval times account for firm and device-type fixed effects. Each dot represents the average approval time of a cluster of document submissions for which submitting firms, on average, have a given level of FDA experience (after accounting for device type and firm fixed effects). Only firms with positive cumulative experience are included in the regression.}}
% \label{fig:LBDResiduals}
% \end{figure}



% \begin{figure}%
%     \centering  \caption{Effect of III to II reclassification on Adverse Events\\ Top Quartile vs. Other}
% % \vspace{-14pt}
%      \centering {\includegraphics[scale=0.5]{../Figures/Event_Estimation_3_2.png} }
%   \caption*{\footnotesize{Note: This figure presents separate estimates of equation \ref{eq:OLS} for each adverse event measure and for device types whose baseline event rate is in the bottom 75th percentile and the top quartile. The x-axis describes the percentile bin, bottom 75th ($X = 0$) or top quartile ($X > 0$), for each of the three adverse event types. The y-axis describes the change in the rate of the given adverse event type in the treated groups relative to matched control groups. The first section of the graph presents results for hospitalization events, the middle section presents results for mortality events, and the last section presents results for life-threatening events. Adverse event data are derived from the FDA MAUDE database. Standard error bars are provided.}}
%      \label{fig:IIItoIIAdverseTopDecile}%
% \end{figure}



% \clearpage
% \newpage

% \clearpage
% \newpage




% \clearpage
% \newpage


% \begin{table}[H]
% \begin{center}
% \caption{\textbf{Reclassification Events Gathered}}
% \begin{tabular}{lr}
% \hline
% \hline
% Reclassification Event    &         Number of Device Types               \\
% \midrule
% Unclassified $\implies$ class I                                        &                      14 \\

% Unclassified $\implies$ class II (no previous PMAs)                     &                    106 \\
% Unclassified $\implies$ class II (w/ previous PMAs)                                   &                      3 \\
% Unclassified $\implies$ class II (but 510(k) exempt)                      &                      5 \\
% Class I Official Exemptions        &                   2104 \\
% Reclassification class I $\implies$ II (Special Controls Only)                                 &                      3 \\
% Reclassification class II $\implies$ I                                        &                      293 \\
% Reclassification class II $\implies$ III                                        &                      2 \\
% 510(k) Exemption for class II                        &                    123 \\
% Reclassification class III $\implies$ II                                       &                     30 \\
% Reclassification class III $\implies$ II (no previous PMAs)                  &                     50 \\
% Required PMA from Class III                                     &                     64 \\
% \bottomrule

% \end{tabular}
% \label{table:ReclassificationEvents}
%  \caption*{ \begin{footnotesize}Note: The table presents the reclassification events I identified over the past 40 years of the MDA history. I study reclassifications from Class III to Class II and from Class II to Class I. Substantive reclassifications from Class III to II are a mix of pre- and post-amendment medical devices. Pre-amendment devices were those that had existed before the MDA enactment.
% \end{footnotesize} }
% \end{center}
% \end{table}

\clearpage
\newpage

\input{../replication_package/output/tables/Formatted/TableD1_SampleConstruction.tex}

\clearpage
\newpage

\begin{table}[h!]
\begin{center}
\caption{\textbf{Class III to II Down-Regulated Device Types\\(FDA-Initiated and Petitioned)}}
\label{tab:ClassIIIRegnums_combined}

% -------- Panel A --------
\textbf{A. Down-Regulated Device Types Considered (FDA-Initiated; Main Sample)}\\[3pt]
\resizebox{\textwidth}{!}{%
\begin{tabular}{llcr}
\toprule
Device Type Description & Part & Year Announced & FR Citation \\
\midrule
Tacrolimus test system. & Clinical Chemistry And Clinical Toxicology Dev... & 2002 & 67 FR 58329 \\
DNA-based TB diagnostic devices. & Immunology And Microbiology Devices & 2014 & 79 FR 31027 \\
Electrical salivary stimulatory system. & Dental Devices & 2015 & 80 FR 72586 \\
Extracorporeal shock wave lithotripter. & Gastroenterology-Urology Devices & 2000 & 65 FR 48612 \\
Rigid gas permeable contact lens. & Ophthalmic Devices & 1994 & 59 FR 10284 \\
Soft contact lens. & Ophthalmic Devices & 1994 & 59 FR 10284 \\
% Soft contact lens care products. & Ophthalmic Devices & 1997 & 62 FR 30988 \\
Bone sonometer. & Radiology Devices & 2008 & 73 FR 40969 \\
Full-field digital mammography system. & Radiology Devices & 2010 & 75 FR 68203 \\
Vascular embolization device. & Cardiovascular Devices & 2002 & 69 FR 77899 \\
Cutaneous carbon dioxide (PcCO2) monitor. & Anesthesiology Devices & 1988 & 54 FR 27160 \\
Vascular graft prosthesis. & Cardiovascular Devices & 1999 & 64 FR 12774 \\
Intervertebral body fusion device. & Orthopedic Devices & 2007 & 72 FR 32172 \\
\bottomrule
\end{tabular}%
}

\vspace{8pt}

% -------- Panel B (combined description) --------
\textbf{B. Petitioned Down-Regulated Device Types}\\[3pt]
\resizebox{\textwidth}{!}{%
\begin{tabular}{llrll}
\toprule
Device Type Description & Part & Year Announced & Petitioning Entity & FR Citation \\
\midrule
Cyclosporine test system. & Clinical Chemistry And Clinical Toxicology Dev... & 2001 & Dade Behring, Inc., and Microgenics, Inc. & 67 FR 58329 \\
Bone grafting material (dental bone repair). & Dental Devices & 2003 & Bicon, Inc. & 70 FR 21949 \\
Nonabsorbable stainless-steel suture. & General And Plastic Surgery Devices & 1986 & Alto Development Corp. & 68 FR 32984 \\
Absorbable synthetic suture (polyglycolic acid). & General And Plastic Surgery Devices & 1990 & Advanced Biosearch Associates & 68 FR 32984 \\
Nonabsorbable synthetic polyester suture. & General And Plastic Surgery Devices & 1990 & Advanced Biosearch Associates & 68 FR 32984 \\
Nonabsorbable silk suture. & General And Plastic Surgery Devices & 1990 & Advanced Biosearch Associates & 68 FR 32985 \\
Tissue adhesive (topical skin approximation). & General And Plastic Surgery Devices & 2007 & Regulatory \& Clinical Research Institute, Inc. & 73 FR 31033 \\
Tumor-associated antigen test system (alpha-fetoprotein). & Immunology And Microbiology Devices & 1996 & Centocor, Inc. & 62 FR 66005 \\
Tumor-associated antigen test system (epithelial ovarian). & Immunology And Microbiology Devices & 1996 & Centocor, Inc. & 62 FR 66005 \\
Tumor-associated antigen test system (prostate-specific antigen). & Immunology And Microbiology Devices & 1996 & Centocor, Inc. & 62 FR 66005 \\
Automated antimicrobial susceptibility test system (short-term). & Immunology And Microbiology Devices & 2000 & BioMerieux Vitek, Inc. & 68 FR 5827 \\
Hepatitis A serological assays (antibody and IgM). & Immunology And Microbiology Devices & 2004 & Beckman Coulter, Inc. & 71 FR 6679 \\
Home uterine activity monitor. & Obstetrical And Gynecological Devices & 1999 & GE Marquette Medical Systems, Inc. & 66 FR 14076 \\
Bone cement (PMMA). & Orthopedic Devices & 1999 & OSMA & 67 FR 46855 \\
Hip joint metal/polymer constrained prosthesis. & Orthopedic Devices & 2001 & OSMA & 67 FR 21173 \\
Stair-climbing wheelchair. & Physical Medicine Devices & 2013 & DEKA Research \& Dev. Corp. & 79 FR 20782 \\
Nuclear magnetic resonance diagnostic device. & Radiology Devices & 1988 & Thomson-CGR Medical Corp. & 53 FR 5078 \\
\bottomrule
\end{tabular}%
}

\caption*{\begin{footnotesize}
Note: Panel A lists device types down-regulated from Class III to II by FDA initiative; Panel B lists device types down-regulated via reclassification petitions. In Panel B, the device description consolidates the original ``Broad'' and ``Narrow'' descriptions: if the narrow term adds specificity, it appears as a brief parenthetical qualifier; if it is redundant, the broad term alone is used. Entries include regulatory part, announcement year, petitioning entity (Panel B), and Federal Register citation. All rows correspond to device types with at least one prior PMA submission.
\end{footnotesize}}
\end{center}
\end{table}
% 283

\clearpage
\newpage



\begin{table}[h!]
\begin{center}
\caption{\textbf{Treated Device Type to Procedure Linkages}}
\resizebox{\textwidth}{!}{%
\begin{tabular}{l p{6.0cm} p{2.6cm} p{2.6cm} p{7.8cm}}
\hline \hline
\addlinespace
\textbf{Product Code} & \textbf{Device Name} & \textbf{CPT (L1)} & \textbf{HCPCS (L2)} & \textbf{Procedure} \\
\hline
\addlinespace
\multicolumn{5}{l}{\textbf{A. Class III to II}} \\
\addlinespace
LNS & Lithotriptor, Extracorporeal Shock-Wave, Urological & 50590 &  & Extracorporeal shock wave lithotripsy (ESWL) for renal calculi (kidney stones). \\
MAX & Intervertebral Fusion Device With Bone Graft, Lumbar & 22558 &  & Anterior lumbar interbody fusion (ALIF). \\
MLM & Enzyme Immunoassay, Tacrolimus & 80197 &  & Therapeutic drug assay of tacrolimus (FK506) in blood. \\
\\
\addlinespace
\multicolumn{5}{l}{\textbf{B. Class II to I}} \\
\addlinespace
CEK & Biuret (Colorimetric), Total Protein & 84155 &  & Total protein, serum -- quantitative biuret colorimetric assay. \\
CER & 2,4-Dinitrophenylhydrazine, Lactate Dehydrogenase & 83615 &  & Lactate dehydrogenase (LDH) quantitative assay. \\
CFH & Tetrazolium Int Dye-Diaphorase, Lactate Dehydrogenase & 83615 &  & Lactate dehydrogenase (LDH) quantitative assay. \\
CFJ & NAD Reduction/NADH Oxidation, Lactate Dehydrogenase & 83615 &  & Lactate dehydrogenase (LDH) quantitative assay. \\
DAD & Haptoglobin, Antigen, Antiserum, Control & 83010 &  & Haptoglobin assay (quantitative immunologic serum haptoglobin test). \\
FFL & Dislodger, Stone, Basket, Ureteral, Metal & 52352 &  & Ureteroscopy with removal of ureteral calculus (stone extraction using basket device). \\
FNL & Bed, AC-Powered Adjustable Hospital &  & E0260 & Hospital bed, semi-electric (head and foot adjustment), with any type side rails, with mattress. \\
FNM & Mattress, Air Flotation, Alternating Pressure &  & E0277 & Powered pressure-reducing air mattress (alternating/low air loss) furnished as DME. \\
IKK & System, Isokinetic Testing And Evaluation & 97750 &  & Physical performance test or measurement with written report, each 15 minutes. \\
ILO & Pack, Hot Or Cold, Water Circulating & 97010 &  & Application of hot or cold packs (superficial thermal modality). \\
JCA & Device, Bleeding Time & 85002 &  & Bleeding time test. \\
JGP & Lowry (Colorimetric), Total Protein & 84155 &  & Measurement of total protein concentration in blood (serum, plasma, or whole blood). \\
JGQ & Turbidimetric, Total Protein & 84155 &  & Total protein; quantitative assay (serum/plasma) -- turbidimetric method. \\
JPI & Device, Hematocrit Measuring & 85014 &  & Hematocrit (microhematocrit) test -- measurement of packed cell volume. \\
KSX & Substance, Blood Grouping Of Non-Human Origin For In Vitro Diagnostic Use & 86900, 86901 &  & Blood grouping (ABO) and Rh(D) typing immunohematology tests. \\
\\
\bottomrule
\end{tabular}}
\caption*{\footnotesize Note: The table provides the device-to-procedure linkages for the two reclassification pathways where (i) procedures are linked to the device, (ii) the reclassification event occurred after 1996 (the earliest year in my data), and (iii) there are outpatient claims attributable to that procedure. Procedure code description text is abbreviated; code descriptors were verified against published references.}
\label{tab:DeviceToProcedureLinkages}
\end{center}
\end{table}
% 301

\clearpage
\newpage

\begin{table}[t]
\centering
\caption{\textbf{Robustness Matrix: Innovation (All estimators)}}
\begin{threeparttable}

% === add double rule across full width ===
\noindent\makebox[\textwidth]{\rule{\textwidth}{0.5pt}}

\vspace{-13pt}

\noindent\makebox[\textwidth]{\rule{\textwidth}{0.5pt}}
% =========================================

\begin{minipage}[t]{0.49\textwidth}
\centering\footnotesize
\textbf{A. Class III to II}\par\vspace{4pt}

\emph{Borusyak et al.}\par
\begin{adjustbox}{max width=\linewidth}
\begin{tabular}{@{}lcccc@{}}
\toprule
 & Pre-mean & Matched & Intuitive & Later \\
\cmidrule(lr){2-5}
\input{../replication_package/output/tables/Formatted/DIDimpute_PaperRegTable3_2_Main_Panel_NoLog_oct22_24.tex}
\bottomrule
\end{tabular}
\end{adjustbox}\vspace{6pt}

\emph{Log(+1)}\par
\begin{adjustbox}{max width=\linewidth}
\begin{tabular}{@{}lcccc@{}}
\toprule
 & Pre-mean & Matched & Intuitive & Later \\
\cmidrule(lr){2-5}
\input{../replication_package/output/tables/Formatted/PaperRegTable3_2_Main_Panel_Log_oct22_24.tex}
\bottomrule
\end{tabular}
\end{adjustbox}\vspace{6pt}

\emph{Quasi-Poisson}\par
\begin{adjustbox}{max width=\linewidth}
\begin{tabular}{@{}lcccc@{}}
\toprule
 & Pre-mean & Matched & Intuitive & Later \\
\cmidrule(lr){2-5}
\input{../replication_package/output/tables/Formatted/PaperRegTable3_2_Main_Panel_NoLog_poisson_oct22_24.tex}
\bottomrule
\end{tabular}
\end{adjustbox}
\end{minipage}\hfill
\begin{minipage}[t]{0.49\textwidth}
\centering\footnotesize
\textbf{B. Class II to I}\par\vspace{4pt}

\emph{Borusyak et al.}\par
\begin{adjustbox}{max width=\linewidth}
\begin{tabular}{@{}lcccc@{}}
\toprule
 & Pre-mean & Matched & Intuitive & Later \\
\cmidrule(lr){2-5}
\input{../replication_package/output/tables/Formatted/DIDimpute_PaperRegTable2_1_Main_Panel_NoLog_oct22_24.tex}
\bottomrule
\end{tabular}
\end{adjustbox}\vspace{6pt}

\emph{Log(+1)}\par
\begin{adjustbox}{max width=\linewidth}
\begin{tabular}{@{}lcccc@{}}
\toprule
 & Pre-mean & Matched & Intuitive & Later \\
\cmidrule(lr){2-5}
\input{../replication_package/output/tables/Formatted/PaperRegTable2_1_Main_Panel_Log_oct22_24.tex}
\bottomrule
\end{tabular}
\end{adjustbox}\vspace{6pt}

\emph{Quasi-Poisson}\par
\begin{adjustbox}{max width=\linewidth}
\begin{tabular}{@{}lcccc@{}}
\toprule
 & Pre-mean & Matched & Intuitive & Later \\
\cmidrule(lr){2-5}
\input{../replication_package/output/tables/Formatted/PaperRegTable2_1_Main_Panel_NoLog_poisson_oct22_24.tex}
\bottomrule
\end{tabular}
\end{adjustbox}
\end{minipage}
\begin{tabular*}{\textwidth}{@{}l@{}}
\end{tabular*}

\begin{tablenotes}[flushleft,para]\footnotesize
\item[] Note: Panel A shows Class III to II reclassifications; Panel B shows Class II to I. Column (1) reports the 5-year pre-event mean for treated device types (``Pre-mean''). Columns (2)–(4) report difference-in-differences estimates using three control sets: (i) \emph{Matched} controls, (ii) \emph{Intuitive} (device types treating similar conditions), and (iii) \emph{Later-treated} device types (treated after the analysis window). 
Estimation is reported three ways: \emph{Borusyak et al. (2021)} imputation estimator (levels), \emph{Log(+1)} where count outcomes are transformed by $\log(y+1)$, and \emph{Quasi-Poisson} for overdispersed counts. Confidence intervals are based on Conley–Taber test statistics; $+,*,**,***$ denote statistical significance at the 0.10, 0.05, 0.01, and 0.001 levels, respectively. Where present, a ``Sample Size'' row refers to the estimation sample used in Columns (2)–(4).
\end{tablenotes}
\end{threeparttable}
\label{tab:Robustness_Innovation}
\end{table}
% 309

\clearpage
\newpage

\begin{table}[t]
\centering
\caption{\textbf{Robustness Matrix: Market Structure (All estimators)}}
\begin{threeparttable}

% === add double rule across full width ===
\noindent\makebox[\textwidth]{\rule{\textwidth}{0.5pt}}

\vspace{-13pt}

\noindent\makebox[\textwidth]{\rule{\textwidth}{0.5pt}}
% =========================================

\begin{minipage}[t]{0.49\textwidth}
\centering\scriptsize
\textbf{A. Class III to II}\par\vspace{4pt}
\emph{Prices (Borusyak et al. only)}\par
\begin{adjustbox}{max width=\linewidth}
\begin{tabular}{@{}lccccc@{}}
\toprule
 &  & \multicolumn{4}{c}{DiD Estimates} \\
\cmidrule(lr){3-6}
 & Pre-mean & All CPTs & Matched & Intuitive & Later \\
\midrule
\input{../replication_package/output/tables/Formatted/DIDimpute_RegTable_Prices_3_2_sep25.tex}
\bottomrule
\end{tabular}
\end{adjustbox}\vspace{6pt}

\emph{Market composition --- Borusyak et al.}\par
\begin{adjustbox}{max width=\linewidth}
\begin{tabular}{@{}lccccc@{}}
\toprule
 &  & \multicolumn{4}{c}{DiD Estimates} \\
\cmidrule(lr){3-6}
 & Pre-mean & All CPTs & Matched & Intuitive & Later \\
\midrule
\input{../replication_package/output/tables/Formatted/DIDimpute_PaperRegTable3_2_MarketComposition_NoLog_oct22_24.tex}
\bottomrule
\end{tabular}
\end{adjustbox}\vspace{6pt}

\emph{Market composition --- Log(+1)}\par
\begin{adjustbox}{max width=\linewidth}
\begin{tabular}{@{}lccccc@{}}
\toprule
 &  & \multicolumn{4}{c}{DiD Estimates} \\
\cmidrule(lr){3-6}
 & Pre-mean & All CPTs & Matched & Intuitive & Later \\
\midrule
\input{../replication_package/output/tables/Formatted/PaperRegTable3_2_MarketComposition_Log_oct22_24.tex}
\bottomrule
\end{tabular}
\end{adjustbox}\vspace{6pt}

\emph{Market composition --- Quasi-Poisson}\par
\begin{adjustbox}{max width=\linewidth}
\begin{tabular}{@{}lccccc@{}}
\toprule
 &  & \multicolumn{4}{c}{DiD Estimates} \\
\cmidrule(lr){3-6}
 & Pre-mean & All CPTs & Matched & Intuitive & Later \\
\midrule
\input{../replication_package/output/tables/Formatted/PaperRegTable3_2_MarketComposition_NoLog_poisson_oct22_24.tex}
\bottomrule
\end{tabular}
\end{adjustbox}
\end{minipage}\hfill
\begin{minipage}[t]{0.49\textwidth}
\centering\scriptsize
\textbf{B. Class II to I}\par\vspace{4pt}

\emph{Prices (Borusyak et al. only)}\par
\begin{adjustbox}{max width=\linewidth}
\begin{tabular}{@{}lccccc@{}}
\toprule
 &  & \multicolumn{4}{c}{DiD Estimates} \\
\cmidrule(lr){3-6}
 & Pre-mean & All CPTs & Matched & Intuitive & Later \\
\midrule
\input{../replication_package/output/tables/Formatted/DIDimpute_RegTable_Prices_2_1_sep25.tex}
\bottomrule
\end{tabular}
\end{adjustbox}\vspace{6pt}

\emph{Market composition --- Borusyak et al.}\par
\begin{adjustbox}{max width=\linewidth}
\begin{tabular}{@{}lccccc@{}}
\toprule
 &  & \multicolumn{4}{c}{DiD Estimates} \\
\cmidrule(lr){3-6}
 & Pre-mean & All CPTs & Matched & Intuitive & Later \\
\midrule
\input{../replication_package/output/tables/Formatted/DIDimpute_PaperRegTable2_1_MarketComposition_NoLog_oct22_24.tex}
\bottomrule
\end{tabular}
\end{adjustbox}\vspace{6pt}

\emph{Market composition --- Log(+1)}\par
\begin{adjustbox}{max width=\linewidth}
\begin{tabular}{@{}lccccc@{}}
\toprule
 &  & \multicolumn{4}{c}{DiD Estimates} \\
\cmidrule(lr){3-6}
 & Pre-mean & All CPTs & Matched & Intuitive & Later \\
\midrule
\input{../replication_package/output/tables/Formatted/PaperRegTable2_1_MarketComposition_Log_oct22_24.tex}
\bottomrule
\end{tabular}
\end{adjustbox}\vspace{6pt}

\emph{Market composition --- Quasi-Poisson}\par
\begin{adjustbox}{max width=\linewidth}
\begin{tabular}{@{}lccccc@{}}
\toprule
 &  & \multicolumn{4}{c}{DiD Estimates} \\
\cmidrule(lr){3-6}
 & Pre-mean & All CPTs & Matched & Intuitive & Later \\
\midrule
\input{../replication_package/output/tables/Formatted/PaperRegTable2_1_MarketComposition_NoLog_poisson_oct22_24.tex}
\bottomrule
\end{tabular}
\end{adjustbox}
\end{minipage}
\begin{tabular*}{\textwidth}{@{}l@{}}
\end{tabular*}

\begin{tablenotes}[flushleft,para]\footnotesize
\item[] Note: Panel A shows Class III to II; Panel B shows Class II to I. Each block reports the 5-year pre-event mean for treated device types (Column 1) and DiD estimates by control set: \emph{All CPTs} (matching baseline prices among all procedure codes), \emph{Matched} (matching outcomes to medical-device controls), \emph{Intuitive} (similar device types), and \emph{Later-treated}. 
The \emph{Prices} blocks are available only for the Borusyak et al. estimator, as the price outcome is continuous and positive; thus, the log price is always the outcome used. Some control-set columns (e.g., Intuitive/Later) are blank for price estimates. For prices, pre-event means are reported in levels, while estimates are on the log scale. The \emph{Market composition} estimation is reported three ways: \emph{Borusyak et al. (2021)} imputation estimator (levels), \emph{Log(+1)} where count outcomes are transformed by $\log(y+1)$, and \emph{Quasi-Poisson} for overdispersed counts.  Confidence intervals use Conley–Taber test statistics; $+,*,**,***$ mark 0.10, 0.05, 0.01, and 0.001 significance. Any ``Sample Size'' rows refer to the estimation sample for the corresponding control-set columns.
\end{tablenotes}
\end{threeparttable}
\label{tab:Robustness_MarketStructure}
\end{table}
% 309


\clearpage
\newpage


\begin{table}[t]
\centering
\caption{\textbf{Robustness Matrix: Adverse Events (All estimators)}}
\begin{threeparttable}

% === add double rule across full width ===
\noindent\makebox[\textwidth]{\rule{\textwidth}{0.5pt}}

\vspace{-13pt}

\noindent\makebox[\textwidth]{\rule{\textwidth}{0.5pt}}
% =========================================

\begin{minipage}[t]{0.49\textwidth}
\centering\footnotesize
\textbf{A. Class III to II}\par\vspace{4pt}

\emph{Borusyak et al.}\par
\begin{adjustbox}{max width=\linewidth}
\begin{tabular}{@{}lcccc@{}}
\toprule
 & Pre-mean & Matched & Intuitive & Later \\
\cmidrule(lr){2-5}
\input{../replication_package/output/tables/Formatted/DIDimpute_PaperRegTable3_2_AdverseEvent_NoLog_oct22_24.tex}
\bottomrule
\end{tabular}
\end{adjustbox}\vspace{6pt}

\emph{Log(+1)}\par
\begin{adjustbox}{max width=\linewidth}
\begin{tabular}{@{}lcccc@{}}
\toprule
 & Pre-mean & Matched & Intuitive & Later \\
\cmidrule(lr){2-5}
\input{../replication_package/output/tables/Formatted/PaperRegTable3_2_AdverseEvent_Log_oct22_24.tex}
\bottomrule
\end{tabular}
\end{adjustbox}\vspace{6pt}

\emph{Quasi-Poisson}\par
\begin{adjustbox}{max width=\linewidth}
\begin{tabular}{@{}lcccc@{}}
\toprule
 & Pre-mean & Matched & Intuitive & Later \\
\cmidrule(lr){2-5}
\input{../replication_package/output/tables/Formatted/PaperRegTable3_2_AdverseEvent_NoLog_poisson_oct22_24.tex}
\bottomrule
\end{tabular}
\end{adjustbox}
\end{minipage}\hfill
\begin{minipage}[t]{0.49\textwidth}
\centering\footnotesize
\textbf{B. Class II to I}\par\vspace{4pt}

\emph{Borusyak et al.}\par
\begin{adjustbox}{max width=\linewidth}
\begin{tabular}{@{}lcccc@{}}
\toprule
 & Pre-mean & Matched & Intuitive & Later \\
\cmidrule(lr){2-5}
\input{../replication_package/output/tables/Formatted/DIDimpute_PaperRegTable2_1_AdverseEvent_NoLog_oct22_24.tex}
\bottomrule
\end{tabular}
\end{adjustbox}\vspace{6pt}

\emph{Log(+1)}\par
\begin{adjustbox}{max width=\linewidth}
\begin{tabular}{@{}lcccc@{}}
\toprule
 & Pre-mean & Matched & Intuitive & Later \\
\cmidrule(lr){2-5}
\input{../replication_package/output/tables/Formatted/PaperRegTable2_1_AdverseEvent_Log_oct22_24.tex}
\bottomrule
\end{tabular}
\end{adjustbox}\vspace{6pt}

\emph{Quasi-Poisson}\par
\begin{adjustbox}{max width=\linewidth}
\begin{tabular}{@{}lcccc@{}}
\toprule
 & Pre-mean & Matched & Intuitive & Later \\
\cmidrule(lr){2-5}
\input{../replication_package/output/tables/Formatted/PaperRegTable2_1_AdverseEvent_NoLog_poisson_oct22_24.tex}
\bottomrule
\end{tabular}
\end{adjustbox}
\end{minipage}
\begin{tabular*}{\textwidth}{@{}l@{}}
\end{tabular*}

\begin{tablenotes}[flushleft,para]\footnotesize
\item[] Note: Panel A shows Class III to II; Panel B shows Class II to I. Column (1) reports the 5-year pre-event mean for treated device types. Columns (2)–(4) report DiD estimates using \emph{Matched}, \emph{Intuitive}, and \emph{Later-treated} control sets. For the Later-treated specification in the Class III to II case, control device types are treated after 2015; consequently, post-2015 observations are excluded by construction. Estimation is shown using the \emph{Borusyak et al. (2021)} imputation estimator (levels), \emph{Log(+1)} ($\log(y+1)$), and \emph{Quasi-Poisson}. Confidence intervals follow Conley–Taber; significance codes are $+,*,**,***$ for 0.10, 0.05, 0.01, and 0.001 levels. ``Sample Size'' rows, where present, refer to the estimation sample for Columns (2)–(4).
\end{tablenotes}
\end{threeparttable}
\label{tab:Robustness_Adverse}
\end{table}
% 309



\clearpage
\newpage

\begin{table}[h!]
\begin{center}
\caption{\textbf{Reclassification Spillovers (Innovation)}}
\begin{tabular}{lcc}
\hline \hline &   & \multicolumn{1}{c}{DiD Estimates}\\
\cmidrule{3-3}
 & Pre-mean & Matched\\
 Reclassification& \multicolumn{1}{c}{(1)}  & \multicolumn{1}{c}{(2)} \\
\hline 
&&\\
\vspace{5mm}
\textbf{A. Class III to II} & &  \\
\input{../replication_package/output/tables/Formatted/PaperRegTableSpillovers_3_2_Spillovers_NoLog.tex}
&&\\
\hline

&&\\
\textbf{B. Class II to I} & &  \\
&&\\
\input{../replication_package/output/tables/Formatted/PaperRegTableSpillovers_2_1_Spillovers_NoLog.tex}

&&\\

\bottomrule 
\end{tabular}
\caption*{ \begin{footnotesize}Note: The table presents estimates of Equation \ref{eq:OLS}, which is a difference-in-differences (DiD) style OLS regression model for device types that are closely related to treated medical device types. Column (1) presents the 5-year baseline average of closely related device types for the outcomes listed on the left-hand side. Column (2) presents OLS estimates of the effect of reclassifications on device types closely related to treated device types, using matched controls. Confidence intervals in Column (2) are calculated using Conley--Taber test statistics. Standard errors allow for clustering at the product-code level. +, *, **, and *** correspond to statistical significance at the 0.10, 0.05, 0.01, and 0.001 levels, respectively.
\end{footnotesize} }
\label{tab:Innovation_Spillovers}
\end{center}
\end{table}
% 342



\clearpage
\newpage



\begin{table}[h!]
\begin{center} 
\caption{\textbf{Class III to II Device Types by Broad Device Category:\\
 Treated Group versus Intuitive Control Group}}

\resizebox{\textwidth}{!}{\begin{tabular}{llrr}
\hline \hline
                                Device Description &                                  Broad Description &  Treatment &  Implantable \\
\midrule
               Compensated Thorpe tube flowmeter.  &        Anesthesiology Devices---Monitoring Devices &          0 &            0 \\
         Cutaneous carbon dioxide (PcCO2) monitor. &        Anesthesiology Devices---Monitoring Devices &          1 &            0 \\
        Pacemaker repair or replacement material.  &  Cardiovascular Devices---Cardiovascular Prosth... &          0 &            1 \\
                    Vascular embolization device.  &  Cardiovascular Devices---Cardiovascular Prosth... &          1 &            1 \\
                       Vascular graft prosthesis---$<$6 mm diameter.  &  Cardiovascular Devices---Cardiovascular Prosth... &          0 &            1 \\
                       Vascular graft prosthesis---Synthetic/Biologic Composite  &  Cardiovascular Devices---Cardiovascular Prosth... &          1 &            1 \\
                Blood specimen collection device.  &  Clinical Chemistry And Clinical Toxicology Dev... &          0 &            0 \\
                          Tacrolimus test system.  &  Clinical Chemistry And Clinical Toxicology Dev... &          1 &            0 \\
          Electrical salivary stimulatory system.  &               Dental Devices---Therapeutic Devices &          1 &            0 \\
 Intranasal Expiratory Resistance Valve (Sleep Apnea) &               Dental Devices---Therapeutic Devices &          0 &            0 \\
           Extracorporeal shock wave lithotripter. &  Gastroenterology-Urology Devices---Therapeutic... &          1 &            0 \\
             Hemodialysis system and accessories.  &  Gastroenterology-Urology Devices---Therapeutic... &          0 &            0 \\
 Mycobacterium tuberculosis immunofluorescent reagents &  Immunology And Microbiology Devices---Serologi... &          0 &            0 \\
 Nucleic acid-based in vitro diagnostic devices... &  Immunology And Microbiology Devices---Serologi... &          1 &            0 \\
                Rigid gas permeable contact lens (daily wear)  &           Ophthalmic Devices---Therapeutic Devices &          1 &            0 \\
                Rigid gas permeable contact lens (extended wear)  &           Ophthalmic Devices---Therapeutic Devices &          0 &            0 \\
                 Soft (hydrophilic) contact lens (daily wear)  &           Ophthalmic Devices---Therapeutic Devices &          1 &            0 \\
                 Soft (hydrophilic) contact lens (extended wear) &           Ophthalmic Devices---Therapeutic Devices &          0 &            0 \\
               Intervertebral body fusion device.  &            Orthopedic Devices---Prosthetic Devices &          1 &            1 \\
    Wrist joint carpal lunate polymer prosthesis.  &            Orthopedic Devices---Prosthetic Devices &          0 &            1 \\
                                  Bone sonometer.  &             Radiology Devices---Diagnostic Devices &          1 &            0 \\
             Emission computed tomography system.  &             Radiology Devices---Diagnostic Devices &          0 &            0 \\
           Full-field digital mammography system.  &             Radiology Devices---Diagnostic Devices &          1 &            0 \\
                       Mammographic X-ray system.  &             Radiology Devices---Diagnostic Devices &          0 &            0 \\
\bottomrule
\end{tabular}
}
\caption*{ \begin{footnotesize}Note: The table presents the devices and their broad categories used in the treatment and intuitive control groups. No ``life-sustaining'' devices are included in the treatment or control groups. The column ``Implantable'' indicates whether a device type is implantable.
\end{footnotesize} }
\label{tab:ClassIIIRegnumswControls}
\end{center}
\end{table}
% 342

\clearpage
\newpage





\begin{table}[h!]
\begin{center} 
\caption{\textbf{Class II to I Device Types by Broad Category,\\ with Intuitive Controls}}

\begin{adjustbox}{max width=\textwidth}
\begin{tabular}{@{}lrrrr@{}}
\hline \hline
\textbf{Category Description} & \textbf{Treated: Count} & \textbf{Treated: Implants} & \textbf{Control: Count} & \textbf{Control: Implants} \\
\midrule
Anesthesiology---Diagnostic Devices & 2 & 0 & 3 & 0 \\
Anesthesiology---Miscellaneous & 2 & 0 & 3 & 0 \\
Anesthesiology---Monitoring Devices & 6 & 0 & 11 & 0 \\
Anesthesiology---Therapeutic Devices & 17 & 0 & 23 & 0 \\
Cardiovascular---Monitoring Devices & 5 & 0 & 5 & 0 \\
Cardiovascular---Prosthetic Devices & 2 & 0 & 2 & 1 \\
Cardiovascular---Surgical Devices & 0 & 0 & 2 & 0 \\
Clinical Chemistry/Toxicology---Test Systems & 6 & 0 & 6 & 0 \\
Clinical Chemistry/Toxicology---Laboratory Instruments & 3 & 0 & 3 & 0 \\
Dental---Diagnostic Devices & 1 & 0 & 2 & 0 \\
Dental---Miscellaneous Devices & 1 & 0 & 1 & 0 \\
Dental---Surgical Devices & 1 & 0 & 2 & 0 \\
Ear, Nose, and Throat---Diagnostic Devices & 0 & 0 & 2 & 0 \\
Ear, Nose, and Throat---Surgical Devices & 2 & 0 & 6 & 0 \\
Gastroenterology-Urology---Diagnostic Devices & 4 & 0 & 20 & 0 \\
Gastroenterology-Urology---Monitoring Devices & 1 & 0 & 1 & 0 \\
Gastroenterology-Urology---Surgical Devices & 2 & 0 & 10 & 0 \\
Gastroenterology-Urology---Therapeutic Devices & 2 & 0 & 19 & 1 \\
General \& Plastic Surgery---Surgical Devices & 0 & 0 & 1 & 0 \\
General Hospital \& Personal Use---Miscellaneous & 7 & 0 & 14 & 0 \\
General Hospital \& Personal Use---Monitoring & 4 & 0 & 5 & 0 \\
General Hospital \& Personal Use---Therapeutic & 4 & 0 & 7 & 0 \\
Hematology \& Pathology---Manual Devices & 0 & 0 & 4 & 0 \\
Hematology \& Pathology---Used by Blood Manufacturer & 0 & 0 & 4 & 0 \\
Immunology \& Microbiology---Diagnostic Devices & 0 & 0 & 1 & 0 \\
Immunology \& Microbiology---Immunological Test Systems & 4 & 0 & 14 & 0 \\
Immunology \& Microbiology---Microbiology Devices & 0 & 0 & 1 & 0 \\
Immunology \& Microbiology---Serological Reagents & 12 & 0 & 47 & 0 \\
Neurological---Diagnostic Devices & 0 & 0 & 1 & 0 \\
Neurological---Therapeutic Devices & 1 & 0 & 1 & 0 \\
Obstetrical \& Gynecological---Diagnostic Devices & 0 & 0 & 1 & 0 \\
Obstetrical \& Gynecological---Surgical Devices & 0 & 0 & 6 & 0 \\
Obstetrical \& Gynecological---Therapeutic Devices & 0 & 0 & 2 & 0 \\
Ophthalmic---Diagnostic Devices & 2 & 0 & 4 & 0 \\
Ophthalmic---Prosthetic Devices & 5 & 2 & 4 & 4 \\
Ophthalmic---Surgical Devices & 0 & 0 & 3 & 0 \\
Orthopedic---Diagnostic Devices & 0 & 0 & 1 & 0 \\
Orthopedic---Surgical Devices & 0 & 0 & 1 & 0 \\
Physical Medicine---Diagnostic Devices & 3 & 0 & 5 & 0 \\
Physical Medicine---Prosthetic Devices & 4 & 0 & 6 & 0 \\
Physical Medicine---Therapeutic Devices & 12 & 0 & 19 & 0 \\
Radiology---Diagnostic Devices & 5 & 0 & 9 & 0 \\
Radiology---Therapeutic Devices & 1 & 0 & 12 & 0 \\
\bottomrule
\end{tabular}
\end{adjustbox}
\caption*{ \begin{footnotesize}Note: The table reports, for each broad device category, the number of device types in the \emph{Treated} set and in the \emph{Intuitive Control} set (\emph{Count}), alongside how many of those device types are implantable (\emph{Implants}). No life-sustaining devices are included in any set. Category labels are harmonized for readability (minor punctuation/capitalization differences consolidated).
\end{footnotesize} }
\label{tab:ClassIIRegnumsTreated}
\end{center}
\end{table}
% 342


\clearpage
\newpage


\begin{table}[t]
\centering
\caption{\textbf{Effects of Reclassification on Innovation, Market Structure, and Adverse Events (Zero-Count Units Excluded)}}
\label{tab:Combined_DropNoCounts}
\begin{threeparttable}

\noindent\makebox[\textwidth]{\rule{\textwidth}{0.5pt}}

\vspace{-13pt}

\noindent\makebox[\textwidth]{\rule{\textwidth}{0.5pt}}
% =========================
% I. Innovation
% =========================
\textbf{I. Innovation}\par\vspace{4pt}

\begin{minipage}[t]{0.49\textwidth}
\centering\footnotesize
\textbf{A. Class III to II}\par\vspace{2pt}
\begin{adjustbox}{max width=\linewidth}
\begin{tabular}{@{}lcccc@{}}
\toprule
 &  & \multicolumn{3}{c}{DiD Estimates}\\
\cmidrule(lr){3-5}
 & Pre-mean & Matched & Intuitive & Later \\
\cmidrule(lr){2-5}
\input{../replication_package/output/tables/Formatted/PaperRegTable3_2_Main_Panel_NoLog_No0_oct22_24.tex}
\bottomrule
\end{tabular}
\end{adjustbox}
\end{minipage}\hfill
\begin{minipage}[t]{0.49\textwidth}
\centering\footnotesize
\textbf{B. Class II to I}\par\vspace{2pt}
\begin{adjustbox}{max width=\linewidth}
\begin{tabular}{@{}lcccc@{}}
\toprule
 &  & \multicolumn{3}{c}{DiD Estimates}\\
\cmidrule(lr){3-5}
 & Pre-mean & Matched & Intuitive & Later \\
\cmidrule(lr){2-5}
\input{../replication_package/output/tables/Formatted/PaperRegTable2_1_Main_Panel_NoLog_No0_oct22_24.tex}
\bottomrule
\end{tabular}
\end{adjustbox}
\end{minipage}

\vspace{10pt}

% =========================
% II. Market Structure
% =========================
\textbf{II. Market Structure}\par\vspace{4pt}

\begin{minipage}[t]{0.49\textwidth}
\centering\footnotesize
\textbf{A. Class III to II}\par\vspace{2pt}
\begin{adjustbox}{max width=\linewidth}
\begin{tabular}{@{}lccccc@{}}
\toprule
 &  & \multicolumn{4}{c}{DiD Estimates}\\
\cmidrule(lr){3-6}
 & Pre-mean & Price & Matched & Intuitive & Later \\
\cmidrule(lr){2-6}
\input{../replication_package/output/tables/Formatted/PaperRegTable3_2_MarketComposition_NoLog_No0_oct22_24.tex}
\bottomrule
\end{tabular}
\end{adjustbox}
\end{minipage}\hfill
\begin{minipage}[t]{0.49\textwidth}
\centering\footnotesize
\textbf{B. Class II to I}\par\vspace{2pt}
\begin{adjustbox}{max width=\linewidth}
\begin{tabular}{@{}lccccc@{}}
\toprule
 &  & \multicolumn{4}{c}{DiD Estimates}\\
\cmidrule(lr){3-6}
 & Pre-mean & Price & Matched & Intuitive & Later \\
\cmidrule(lr){2-6}
\input{../replication_package/output/tables/Formatted/PaperRegTable2_1_MarketComposition_NoLog_No0_oct22_24.tex}
\bottomrule
\end{tabular}
\end{adjustbox}
\end{minipage}

\vspace{10pt}

% =========================
% III. Adverse Events
% =========================
\textbf{III. Adverse Events}\par\vspace{4pt}

\begin{minipage}[t]{0.49\textwidth}
\centering\footnotesize
\textbf{A. Class III to II}\par\vspace{2pt}
\begin{adjustbox}{max width=\linewidth}
\begin{tabular}{@{}lcccc@{}}
\toprule
 &  & \multicolumn{3}{c}{DiD Estimates}\\
\cmidrule(lr){3-5}
 & Pre-mean & Matched & Intuitive & Later \\
\cmidrule(lr){2-5}
\input{../replication_package/output/tables/Formatted/PaperRegTable3_2_AdverseEvent_NoLog_No0_oct22_24.tex}
\bottomrule
\end{tabular}
\end{adjustbox}
\end{minipage}\hfill
\begin{minipage}[t]{0.49\textwidth}
\centering\footnotesize
\textbf{B. Class II to I}\par\vspace{2pt}
\begin{adjustbox}{max width=\linewidth}
\begin{tabular}{@{}lcccc@{}}
\toprule
 &  & \multicolumn{3}{c}{DiD Estimates}\\
\cmidrule(lr){3-5}
 & Pre-mean & Matched & Intuitive & Later \\
\cmidrule(lr){2-5}
\input{../replication_package/output/tables/Formatted/PaperRegTable2_1_AdverseEvent_NoLog_No0_oct22_24.tex}
\bottomrule
\end{tabular}
\end{adjustbox}
\end{minipage}

\begin{tabular*}{\textwidth}{@{}l@{}}
\end{tabular*}


\begin{tablenotes}[flushleft,para]\footnotesize
\item[] Note: Each panel reports estimates from Equation~\ref{eq:OLS}, a difference-in-differences (DiD) OLS specification. Device types with only zero counts for a given outcome over the sample are excluded (\emph{Zero-Count Units Excluded}). Column~(1) in every subpanel reports the 5-year pre-event mean for treated device types. For the Innovation and Adverse Events blocks, Columns~(2)–(4) report DiD estimates using three control sets: \emph{Matched}, \emph{Intuitive} (treat similar diseases), and \emph{Later-treated}. For the Market Structure block, Columns~(2)–(5) report DiD estimates using \emph{Price-matched} controls, \emph{Matched} controls (matched on baseline innovation and adverse-event levels), \emph{Intuitive} controls, and \emph{Later-treated} controls, respectively. For the Later-treated specification in the Class~III to II case, control device types are treated after 2015; accordingly, post-2015 observations are dropped by construction. Sample sizes reflect the minimum N across outcomes within each panel; N may vary slightly by outcome because the zero-count exclusion is applied separately for each outcome variable. Confidence intervals are computed using Conley–Taber test statistics. Significance codes: $+$, $*$, $**$, and $***$ denote 0.10, 0.05, 0.01, and 0.001 levels, respectively.
\end{tablenotes}

\end{threeparttable}
\end{table}
% 354

\clearpage
\newpage

\input{../replication_package/output/tables/Formatted/TableD11_Utilization.tex}

\clearpage
\newpage

%  % &  Mean at && \multicolumn{2}{c}{DID Estimates} \\
%  %  \cmidrule{4-5}
%  %  & $ -5\leq t < 0 $ &&\multicolumn{1}{c}{Matched} & \multicolumn{1}{c}{Full} \\
%  %  \cmidrule{2-2} \cmidrule{4-5}
%  %  & \multicolumn{1}{c}{(1)}  &&\multicolumn{1}{c}{(2)} & \multicolumn{1}{c}{(3)}

% \begin{table}[h!]
% \begin{center}
% \caption{\textbf{Effect of Reclassification on Innovation \\(Drop No Counts)}}
% \resizebox{\textwidth}{!}{\begin{tabular}{lcccc}
% \hline \hline &   & \multicolumn{3}{c}{DiD Estimates}\\
% \cmidrule{3-5}
%  & Pre-mean & Matched & Intuitive& Later\\
%  Reclassification& \multicolumn{1}{c}{(1)}  & \multicolumn{1}{c}{(2)}&\multicolumn{1}{c}{(3)} &\multicolumn{1}{c}{(4)}\\
% \hline 
% &&&&\\
% \vspace{5mm}
% A. Class III to II: &&&&\\

% \input{../replication_package/output/tables/Formatted/PaperRegTable3_2_Main_Panel_NoLog_No0_oct22_24.tex}
% &&&&\\
% \hline

% &&&&\\
% B. Class II to I:&&&&\\
% &&&&\\
% \input{../replication_package/output/tables/Formatted/PaperRegTable2_1_Main_Panel_NoLog_No0_oct22_24.tex}

% &&&&\\

% \hline 
% \end{tabular}}
% \caption*{\begin{footnotesize}Note: The table presents estimates of equation \ref{eq:OLS}, which is a difference-in-differences (DiD) style OLS regression model. In this analysis, I drop all device types that do not exhibit any positive quantity of the given outcome. Column (1) presents the 5-year baseline average of treated device types for the outcomes listed on the left-hand side. Columns (2)--(4) present DiD estimates for the listed outcomes using different control groups: namely, a matched control group, intuitively similar device types (treat similar diseases), and ``later-treated'' device types (treated after the sample window), respectively. Confidence intervals are calculated using Conley--Taber test statistics. +, *, **, and *** correspond to statistical significance at 0.10, 0.05, 0.01, and 0.001 levels, respectively.
% \end{footnotesize} }
% \label{tab:Innovation_DropNoCounts}
% \end{center}
% \end{table}
% % 422



% \clearpage
% \newpage



% \begin{table}[h!]
% \begin{center}
% \caption{\textbf{Market Structure Estimates (Drop No Counts)}}
% \resizebox{\textwidth}{!}{
% \small
% \begin{tabular}{lccccc}
% \hline \hline 
% &   & \multicolumn{4}{c}{DiD Estimates}\\
% \cmidrule{3-6}
%  & Pre-mean & Price & Matched & Intuitive & Later \\
%  Reclassification& \multicolumn{1}{c}{(1)}  & \multicolumn{1}{c}{(2)}&\multicolumn{1}{c}{(3)} & \multicolumn{1}{c}{(4)} & \multicolumn{1}{c}{(5)} \\
%  \addlinespace

% \midrule 
% % &&&&\\
% \addlinespace
% A. Class III to II: & & & & & \\
% \addlinespace

% % &&&&\\
% % \input{../Tables/Formatted/RegTable_Prices_3_2.tex}
% % &&&&\\
% % \addlinespace

% % \cmidrule{2-6}\\

% \input{../replication_package/output/tables/Formatted/PaperRegTable3_2_MarketComposition_NoLog_No0_oct22_24.tex}

% % &&&&\\
% \addlinespace

% \hline

% \addlinespace


% B. Class II to I: & & & & & \\
% \addlinespace

% % &&&&\\
% % \input{../Tables/Formatted/RegTable_Prices_2_1.tex}
% % &&&&\\
% % \addlinespace

% % \cmidrule{2-6}\\

% % &&&&\\
% \input{../replication_package/output/tables/Formatted/PaperRegTable2_1_MarketComposition_NoLog_No0_oct22_24.tex}

% % &&&\\
% \addlinespace

% \hline 
% \end{tabular}}

% \caption*{ \begin{footnotesize}Note: The table presents estimates of equation \ref{eq:OLS}, which is a difference-in-differences (DiD) style OLS regression model. In this analysis, I drop all device types that do not exhibit any positive quantity of the given outcome. Column (1) presents the 5-year baseline average of treated device types for the outcomes listed on the left-hand side. Columns (2)--(5) present DiD estimates for a given outcome using different control groups: namely, a group matched on baseline prices, a group matched on baseline innovation and adverse event levels, an intuitively comparable group, and a later-treated group, respectively. Confidence intervals are calculated using Conley--Taber test statistics.  +, *, **, and *** correspond to statistical significance at 0.10, 0.05, 0.01, and 0.001 levels, respectively.
% \end{footnotesize} }
% \label{tab:MarketComposition_NoZeros}
% \end{center}
% \end{table}
% % 422



% \clearpage
% \newpage


% \begin{table}[h!]
% \begin{center}
% \caption{\textbf{Effect of Reclassifications on Adverse Events\\ (Drop No Counts)}}
% \resizebox{\textwidth}{!}{\begin{tabular}{lcccc}
% \hline \hline &   & \multicolumn{3}{c}{DiD Estimates}\\
% \cmidrule{3-5}
%  & Pre-mean & Matched & Intuitive & Later \\
%  Reclassification& \multicolumn{1}{c}{(1)}  & \multicolumn{1}{c}{(2)}&\multicolumn{1}{c}{(3)} &\multicolumn{1}{c}{(4)}\\
% \hline 
% &&&&\\
% \vspace{5mm}
% A. Class III to II: &&&&\\
% % \hspace{-1mm}\input{../replication_package/output/tables/Formatted/PaperRegTableSafety_IIItoII_NoLog.tex}
% \input{../replication_package/output/tables/Formatted/PaperRegTable3_2_AdverseEvent_NoLog_No0_oct22_24.tex}
% &&&&\\
% \hline

% &&&&\\
% B. Class II to I: &&&&\\
% &&&&\\
% % \hspace{-1mm}\input{../replication_package/output/tables/Formatted/PaperRegTableSafety_IItoI_NoLog.tex}
% \input{../replication_package/output/tables/Formatted/PaperRegTable2_1_AdverseEvent_NoLog_No0_oct22_24.tex}

% &&&&\\

% \hline 
% \end{tabular}}
% \caption*{ \begin{footnotesize}Note: The table presents estimates of equation \ref{eq:OLS}, which is a difference-in-differences (DiD) style OLS regression model. In this analysis, I drop all device types that do not exhibit any positive quantity of the given outcome. Column (1) presents the 5-year baseline average of treated device types for the outcomes listed on the left-hand side. Columns (2)--(4) present DiD estimates for the listed outcomes using different control groups: namely, a matched control group, intuitively similar device types (treat similar diseases), and ``later-treated'' device types (treated after the sample window), respectively. For Column (4), Class III to II,  control device types are treated after 2015; thus, all observations after 2015 are dropped. Confidence intervals are calculated using Conley--Taber test statistics. +, *, **, and *** correspond to statistical significance at 0.10, 0.05, 0.01, and 0.001 levels, respectively.
% \end{footnotesize} }
% \label{tab:AdverseEvent_NoZeros}
% \end{center}
% \end{table}
% % 422






% \clearpage
% \newpage


% \begin{table}[h!]
% \begin{center}
% \caption{\textbf{Effect of Reclassification on Logged Market Size}}
% \resizebox{0.6\textwidth}{!}{%
% \begin{tabular}{lcc}
% \hline \hline
% \textbf{Outcome} & \textbf{Estimate} & \textbf{p-value} \\
% \hline

% \addlinespace
% B. Class III to II & & \\
% \addlinespace
% \hspace{6mm}Market Size & 0.94** & 0.001 \\

% \addlinespace

% A. Class II to I & & \\
% \addlinespace
% \hspace{6mm}Market Size & 0.01 & 0.92 \\
% \hline
% \end{tabular}}
% \caption*{
% \begin{footnotesize}
% Note: This table presents the estimated effects of reclassifications on logged market size. Market size is defined as the procedure price multiplied by quantity. Panel A corresponds to the transition from Class III to Class II (3 to 2), and Panel B corresponds to the transition from Class II to Class I (2 to 1). Estimates are rounded to two decimal places, and p-values are rounded to three decimal places. Statistical significance is indicated by * ($p < 0.05$) and ** ($p < 0.01$).
% \end{footnotesize}
% }
% \label{tab:MarketSizeEffects}
% \end{center}
% \end{table}
% 483


% \clearpage
% \newpage


% \begin{table}[!ht]
% \centering \caption{\small{\textbf{Flattening the Learning Curve Simulation---Unique Devices Approved}}}
% \begin{tabular}{lrrrrc}
% \hline \hline

%  & \multicolumn{4}{c}{Percent Changes}\\
%   \cmidrule{2-5}

%  $\gamma$ &  $T_{Sum,25}$ &  $T_{Sum,50}$ &  $T_{Sum,75}$ & $T_{Sum,100}$ & Total \%$\Delta$\\
% \midrule
% 0.072 & 0.0 & 0.0 & 0.0 & 0.0 & 0.0\\
%   & (0.0) & (0.0) & (0.0) & (0.0) & (0.0)\\
% 0.064 & 12.08 & 7.26 & 3.57 & 1.76 & 5.04\\
%   & (8.13) & (5.4) & (2.85) & (2.18) & (4.16)\\
% 0.056 & 25.97 & 13.11 & 6.41 & 3.98 & 9.99\\
%   & (12.47) & (6.87) & (3.87) & (3.59) & (6.87)\\
% 0.048 & 45.94 & 20.44 & 12.41 & 5.15 & 16.31\\
%   & (17.58) & (11.25) & (4.97) & (3.28) & (10.02)\\
% 0.04 & 54.46 & 24.44 & 13.28 & 6.71 & 19.77\\
%   & (18.76) & (8.15) & (5.88) & (3.65) & (12.05)\\
% 0.032 & 68.43 & 33.07 & 15.16 & 7.72 & 25.03\\
%   & (25.16) & (11.68) & (5.62) & (4.7) & (15.52)\\
% 0.024 & 84.04 & 39.42 & 21.22 & 9.66 & 30.95\\
%   & (30.03) & (13.0) & (8.46) & (4.91) & (18.29)\\
% 0.016 & 90.02 & 44.71 & 24.39 & 11.19 & 34.65\\
%   & (33.71) & (13.69) & (8.83) & (4.53) & (19.69)\\
% 0.008 & 112.41 & 58.62 & 29.23 & 15.46 & 43.87\\
%   & (38.99) & (18.13) & (10.68) & (5.91) & (23.23)\\
% 0.0 & 132.36 & 59.07 & 34.21 & 16.48 & 48.35\\
%   & (41.23) & (19.01) & (11.87) & (6.71) & (25.7)\\

% \bottomrule
% \end{tabular}

            
%     \caption*{\footnotesize{Note: This table presents the results of the simulation exercise described in Appendix \ref{sec:appendixGammaSimulation}, which simulates the effect of flattening the learning curve on the rate of unique device submissions at an annual frequency by asset quartiles. Figure \ref{fig:GammaSimulation} illustrates this flattening exercise. Standard errors generated from a Monte Carlo procedure are presented in parenthesis below the estimates. This procedure produces estimates across repeated random draws from the empirical distribution of firm characteristics to calculate confidence intervals. I express changes as percent changes relative to the $\gamma = 0.075$ baseline. I flatten the learning curve relative to the firm with the highest experience in the data. In the Table, $\gamma$ begins at its initial starting point estimated in equation \ref{eq:loglinlearning}. Subsequent rows in the table show the percent change in the rate of unique device submissions as $\gamma$, the learning rate, is reduced. These changes are presented for each experience quartile for Class III device manufacturers. $T_{Sum,25}$ represents the bottom 25th percentile of cumulative FDA experience (in days), $T_{Sum,50}$ represents the 25--50th percentile, $T_{Sum,75}$ represents the 50--75th percentile, and $T_{Sum,100}$ represents the 75th--100th percentile. The far-right column presents the total percent change in unique device submissions from a flattening of the learning curve relative to the baseline frequency of unique device submissions.}}
%     \label{tab:FlattenCurveSimul}%
% \end{table}
% 510

% \clearpage
% \newpage


\begin{table}[h!]
\centering
\caption{\textbf{Effect of Reclassification on Probability of Any Safety Mention (Robustness Across Alternative Approaches)}}
\resizebox{\textwidth}{!}{%
\begin{tabular}{llcc}
\hline \hline
Outcome & Approach & Pre-mean & Estimate \\
\midrule

\input{../replication_package/output/tables/Formatted/TableD12_SafetyRobustness.tex}
\bottomrule
\end{tabular}}
\caption*{
\begin{footnotesize}
Note: This table presents the estimated effects of reclassification events on the probability of safety mentions, analyzed through two distinct approaches. Panel A reports results for Class II to I reclassifications (sample size: 9,856), while Panel B reports results for Class III to II reclassifications (sample size: 1,056). The "pre-mean" column indicates the average outcome prior to reclassification. The estimates are rounded to three decimal places, with standard errors (SE) shown in parentheses directly below the estimates. For the Safety-Only Approach, the outcome measures the probability of observing any mention of the keyword "safety" within a device type-year pair. For the Embedding Distance Approach, the outcome measures whether any patent was identified as semantically similar to a set of five safety-related reference sentences using the BERT embedding model, within the same device type-year pair. Statistical significance is denoted as follows: + ($p < 0.10$), * ($p < 0.05$), ** ($p < 0.01$), and *** ($p < 0.001$).
\end{footnotesize}
}
\label{tab:ProbabilitySafetyMentionEffects_Robustness}
\end{table}
% 437


\clearpage
\newpage

\input{../replication_package/output/tables/Formatted/TableD11a_DeviceExperience_3_2.tex}
% 495

\newpage

\input{../replication_package/output/tables/Formatted/TableD11b_DeviceExperience_2_1.tex}
% 495


\clearpage
\newpage

\input{../replication_package/output/tables/Formatted/TableD15_CostBenefit.tex}
% 497




\clearpage
\newpage

\input{../replication_package/output/tables/Table_D14_summary_stats_by_class.tex}

% 908
  


% \begin{sidewaystable}[h!]
% \begin{center} 
% \caption{\textbf{Down-Classified Class III Device Types Considered}}

% \resizebox{\textwidth}{!}{\begin{tabular}{rrlllrl}
% \toprule
%  Group &  D &                              Description &                Subpart &             Part &  Implant &                             Reason \\
% \midrule
%      1 &          1 &                           Tacrolimus test system.  &                      Test Systems  &           Clinical chemistry &             &     Both immunosuppressant enzymes \\
% 1 &          0 &                            Sirolimus test system.  &                       Test Systems &           Clinical chemistry &             &                                 \\
%      2 &          0 &  TB immunofluorescent &                           Reagents &  Immunology &             &           Both target tuberculosis \\
%      2 &          1 &  TB nucleic acid respiratory &                           Reagents &  Immunology &             &                                 \\
%      3 &          1 &          Cutaneous C02 monitor. &                 Monitoring  &               Anesthesiology &             &          Both related gas monitors \\
%      3 &          0 &                   Cutaneous 02 monitor. &                 Monitoring  &               Anesthesiology &             &                                 \\
%      4 &          0 &                                    Vascular clip.  &   Prosthetic  &               Cardiovascular &            \checkmark &               Both treat aneurysms \\
%      4 &          1 &                       Vascular embolization device &   Prosthetic  &               Cardiovascular &            \checkmark &                                 \\
%      5 &          1 &  Vascular graft---w/ coating &   Prosthetic  &               Cardiovascular &            \checkmark &  Same device, different composites \\
%      5 &          0 &  Vascular graft prosthesis---w/o coating&   Prosthetic  &               Cardiovascular &            \checkmark &                                 \\
%      6 &          0 &                   Orthodontic headgear.  &                Therapeutic  &                       Dental &             &    Both therapeutic dental devices \\
%      6 &          1 &           Electrical salivary stimulator.  &                Therapeutic  &                       Dental &             &                                 \\
%      7 &          1 &            Shock wave lithotripter. &                Therapeutic  &     Gastro. &             &           Both treat kidney stones \\
%      7 &          0 &                          Laser surgical instrument &                   Surgical  &  General &             &             (lasers or shockwaves) \\
%      8 &          0 &     Hard contact lens---extended wear &                Therapeutic  &                   Ophthalmic &             &         Daily versus extended-wear \\
%      8 &          1 &       Hard contact lens---daily wear. &                Therapeutic  &                   Ophthalmic &             &                                 \\
%      9 &          0 &      Soft  contact lens---extended wear &                Therapeutic  &                   Ophthalmic &             &         Daily versus extended-wear \\
%      9 &          1 &        Soft contact lens---daily wear. &                Therapeutic  &                   Ophthalmic &             &                                 \\
%     10 &          0 &          Thoracolumbosacral pedicle screw system.  &                 Prosthetic  &                   Orthopedic &            \checkmark &           Both are spinal implants \\
%     10 &          1 &                Intervertebral body fusion device.  &                 Prosthetic  &                   Orthopedic &            \checkmark &                                 \\
%     11 &          0 &                                Bone densitometer.  &                 Diagnostic  &                    Radiology &             &           Both measure bone health \\
%     11 &          1 &                                   Bone sonometer.  &                 Diagnostic  &                    Radiology &             &    (ultrasound vs. x-ray) \\
%     12 &          0 &                        Mammographic x-ray system.  &                 Diagnostic  &                    Radiology &            &                   Both mammography \\
%     12 &          1 &            Full-field digital mammography system.  &                 Diagnostic  &                    Radiology &             &       (full-field vs. traditional) \\
% \bottomrule
% \end{tabular}}

% \caption*{ \begin{footnotesize}Note: The table presents the device types which experienced a reclassification from Class III to II by the FDA's own initiative and for which there was at least one submitted PMA document prior to reclassification. These device types form my treated group in the Class III to II reclassification analysis. I provide device type descriptions, in which part they pertain, and the year they were down-classified. No life-sustaining devices are considered in treated or control groups (e.g., ventilators).
% \end{footnotesize} }
% \label{tab:ClassIIIRegnums}
% \end{center}
% \end{sidewaystable}





\clearpage
\newpage

\pdfpageattr\expandafter{/Rotate 0}

\begin{table}[h!]
\centering
\caption{\textbf{Impact of Reclassification on Patenting Rates:\\Restricted Samples}}
\label{tab:Innovation_nodrugs_A61only}
\begin{tabular}{lcc}
\hline \hline
 & \multicolumn{2}{c}{Patenting Rate} \\
\cmidrule(lr){2-3}
\hspace{6mm}Restriction & Pre-mean & DiD Estimate \\
\hline
\input{../replication_package/output/tables/Formatted/TableD15_RestrictedSamples.tex}
\bottomrule
\end{tabular}
\caption*{
\begin{footnotesize}
Note: The table presents estimates of Equation \ref{eq:OLS}, a difference-in-differences (DiD) OLS regression model analyzing patenting rates for restricted patent samples described in Appendix \ref{sec:appendixPatentCollection}. Panel A examines reclassifications from Class III to II, while Panel B examines Class II to I. \textit{GPT-4o Identified Devices} refers to patents determined as devices by GPT-4o, and \textit{Health Patents (No Drugs)} refers to patents classified under the CPC system as "health" (A61) excluding drug patents (not A61P). The ``Pre-mean'' column presents the 5-year baseline average of treated device types for the restrictions listed on the left-hand side. The ``DiD Estimate'' column provides DiD estimates when the treated groups are compared to matched controls for the restrictions listed on the left-hand side. Confidence intervals are calculated using Conley--Taber test statistics. Standard errors are presented in parentheses below the estimates. Significance levels are denoted as follows: + p$<$0.10, * p$<$0.05, ** p$<$0.01, *** p$<$0.001.
\end{footnotesize}
}
\end{table}
% 1704



 

\clearpage
\newpage
\begin{table}[h!]
\centering
\caption{\textbf{Effect of Reclassification on Breakthrough Patents}}
\label{tab:BreakthroughPatentsEffects}
\begin{tabular}{lcccc}
\toprule
 & Pre-mean & Q-Poisson & OLS & Log \\
\midrule
\input{../replication_package/output/tables/Formatted/TableD16_BreakthroughPatents.tex}
\bottomrule
\end{tabular}
\caption*{\footnotesize
Note: This table presents the estimated effects of reclassification on the yearly count of breakthrough patents across multiple models (Quasi-Poisson, OLS linear model, and OLS log outcome). Panel~A corresponds to down-regulation (III to II) with a sample size of 1,584, and Panel~B corresponds to deregulation (II to I) with a sample size of 11,748. Standard errors in parentheses. Significance: $+$ $p<0.10$, $*$ $p<0.05$, $**$ $p<0.01$.}
\end{table}
% 1710





% \clearpage
% \newpage

% \pdfpageattr\expandafter{/Rotate 0}


% \begin{table}[H]
% \begin{center}
% \caption{\textbf{Keywords Used in Text Analysis of Patent Claims}}
% \begin{tabular}{ll}
% \multicolumn{2}{l}{Safety Advancement Keywords} \\
% \midrule
% \midrule
% safety                 & hazard                 \\
% safe                   & danger                 \\
% safer                  & risking              \\
% endangering            & harming                \\
% precautions            & injuring               \\
% unsafe                 & injury                 \\
% hazardous              & jeopardizing           \\
% failsafe               & risk                   \\
% safely                 & complication           \\
% dangerous              & life-threatening          \\
% \bottomrule

% \end{tabular}
% \label{table:safety_keywords}
%  \caption*{ \begin{footnotesize}Note: The table presents the keywords related to product safety that were extracted using the Word2Vec algorithm. I label a patent as advancing safety if any of the above words are included in its claims section. Interestingly, some keywords indicate safety advancements in what the product is not: some inventors claim advancements in product safety by moving away from constructions that are ``hazardous,'' ``unsafe,'' or ``dangerous.'' It is important to note that inventors would not reasonably claim a product advancement that would lead to more injuries. Thus, one can assume that these negative mentions can still be attributable to safety improvements.
% \end{footnotesize} }
% \end{center}
% \end{table}
% % 2489




















% \clearpage
% \newpage



% \input{../Tables/Corr_Size_Experience.tex} %from Corr_Size_Experience.do











% \begin{table}[h!]
% \begin{center}
% \caption{\textbf{Effect of reclassifications on Innovation:\\ Restricted Patent Sample}}
% \resizebox{.9\textwidth}{!}{\begin{tabular}{lcccc}
% \hline \hline &   & \multicolumn{3}{c}{DiD Estimates}\\
% \cmidrule{3-5}
%  & Pre-mean & Matched & Intuitive& Later  \\
%  reclassification& \multicolumn{1}{c}{(1)}  & \multicolumn{1}{c}{(2)}&\multicolumn{1}{c}{(3)} &\multicolumn{1}{c}{(4)} \\
% \hline 
% &&&&\\
% \vspace{5mm}
% A. Class III to II: &&&&\\

% \input{../Tables/Formatted/PaperRegTable3_2_Main_Panel_NoLog_NoDrugs_A61_oct22_24.tex}
% &&&&\\
% \hline

% &&&&\\
% B. Class II to I:&&&&\\
% &&&&\\
% \input{../Tables/Formatted/PaperRegTable2_1_Main_Panel_NoLog_NoDrugs_A61_oct22_24.tex}

% &&&&\\

% \hline 
% \end{tabular}}
% \caption*{ \begin{footnotesize}Note: The table presents estimates of equation \ref{eq:OLS}, which is a difference-in-differences (DiD) style OLS regression model. This table differs from Table \ref{tab:Innovation} in that it presents estimates from an estimation that uses a restricted patent sample described in Appendix \ref{sec:appendixPatentCollection}, and only presents the patenting rate outcome. Simply put, patents in this analysis include only those labeled as health-related and non-drug by patent examiners. Patents are derived from the \textit{Dimensions} patent database. Column (1) presents the 5-year baseline average of treated device types for the outcomes listed on the left-hand side. Columns (2)--(4) present DiD estimates for the listed outcomes using different control groups: namely, a matched control group, intuitively similar device types (treat similar diseases), and ``later-treated'' device types (treated after the sample window), respectively. For Column (4), Class III to II,  control device types are treated after 2015; thus, all observations after 2015 are dropped. Confidence intervals are calculated using Conley--Taber test statistics. +, *, **, and *** correspond to statistical significance at 0.10, 0.05, 0.01, and 0.001 levels, respectively.
% \end{footnotesize} }
% \label{tab:Innovation_nodrugs_A61only}
% \end{center}
% \end{table}


% \clearpage
% \newpage



% \clearpage
% \newpage



% \begin{table}[h!]
% \begin{center} 
% \caption{\textbf{Estimation of Learning Curve Parameters (in Days)}}
% \input{../Tables/LBDGammaEst_oct22_24.tex}

% \caption*{ \begin{footnotesize}Note: The table presents the estimates of equation \ref{eq:loglinlearning}, which estimates the learning coefficient $\gamma$ and the baseline time requirement $\beta(R_c)$ for both Class III original PMA approvals (Column 1) and Class II 510(k) clearance (Column 2) of unique devices via OLS. The estimates for Class III devices are calculated by only considering the approval times of filed original PMAs by firms with at least one day of prior experience navigating FDA regulations. The estimates for Class II devices are calculated by only considering the clearance times of 510(k) documents for unique devices that were submitted by firms with at least one day of prior experience navigating FDA regulations. Prior experience is calculated using approval times when filing any prior documentation type (510(k) or PMAs). Standard errors are clustered at the firm level. +, *, **, and *** correspond to statistical significance at 0.10, 0.05, 0.01, and 0.001 levels, respectively.
% \end{footnotesize} }
% \label{tab:LearningParameter}
% \end{center}
% \end{table}
% % 2563













  



\end{document}
